%%=============================================================
%% preamble_ai3.tex
%% Preamble dùng chung cho chapter4_4.tex và chapter5_4.tex
%%
%% Đưa vào main.tex bằng:  %%=============================================================
%% preamble_ai3.tex
%% Preamble dùng chung cho chapter4_4.tex và chapter5_4.tex
%%
%% Đưa vào main.tex bằng:  %%=============================================================
%% preamble_ai3.tex
%% Preamble dùng chung cho chapter4_4.tex và chapter5_4.tex
%%
%% Đưa vào main.tex bằng:  %%=============================================================
%% preamble_ai3.tex
%% Preamble dùng chung cho chapter4_4.tex và chapter5_4.tex
%%
%% Đưa vào main.tex bằng:  \input{preamble_ai3}
%% Hoặc copy thẳng các \usepackage cần thiết vào preamble của main.
%%=============================================================

%% Tiếng Việt
\usepackage[T5]{fontenc}
\usepackage[utf8]{inputenc}
\usepackage[vietnamese]{babel}

%% TikZ và các thư viện vẻ hình
\usepackage{tikz}
\usetikzlibrary{
  shapes.geometric,
  arrows.meta,
  positioning,
  fit,
  backgrounds,
  calc,
  decorations.pathreplacing,
  patterns,
}

%% PGFPlots cho biểu đồ
\usepackage{pgfplots}
\pgfplotsset{compat=1.18}
\usepackage{pgfplotstable}

%% Code listings
\usepackage{listings}
\usepackage{xcolor}

%% Màu dùng chung (cần định nghĩa trước khi \input chapter)
\definecolor{cppblue}{RGB}{30,100,200}
\definecolor{pygreen}{RGB}{40,150,80}
\definecolor{warnred}{RGB}{200,60,60}
\definecolor{orng}{RGB}{220,130,30}

%% Cài đặt listings mặc định
\lstset{
  basicstyle=\ttfamily\small,
  keywordstyle=\color{cppblue},
  commentstyle=\color{gray!70},
  stringstyle=\color{orng},
  breaklines=true,
  frame=single,
  rulecolor=\color{gray!40},
  framexleftmargin=4pt,
  xleftmargin=8pt,
  numbers=none,
}

%% Hình ảnh
\usepackage{graphicx}

%% Bảng đẹp hơn
\usepackage{booktabs}
\usepackage{multirow}
\usepackage{subcaption}

%% Đảm bảo float không trọn ra ngoài section
\usepackage[section]{placeins}

%%=============================================================
%% Cách sử dụng trong main.tex:
%%
%%   \documentclass[12pt,a4paper]{report}
%%   \input{preamble_ai3}
%%   \begin{document}
%%     ...
%%     \chapter{Thiết kế và Hiện thực}
%%     ...
%%     \input{chapter4_4}   %% mục 4.4
%%     ...
%%     \chapter{Kết quả và Đánh giá}
%%     ...
%%     \input{chapter5_4}   %% mục 5.4
%%     ...
%%   \end{document}
%%
%% Compile: pdflatex -> pdflatex (chạy 2 lần để có label/ref đúng)
%%=============================================================

%% Hoặc copy thẳng các \usepackage cần thiết vào preamble của main.
%%=============================================================

%% Tiếng Việt
\usepackage[T5]{fontenc}
\usepackage[utf8]{inputenc}
\usepackage[vietnamese]{babel}

%% TikZ và các thư viện vẻ hình
\usepackage{tikz}
\usetikzlibrary{
  shapes.geometric,
  arrows.meta,
  positioning,
  fit,
  backgrounds,
  calc,
  decorations.pathreplacing,
  patterns,
}

%% PGFPlots cho biểu đồ
\usepackage{pgfplots}
\pgfplotsset{compat=1.18}
\usepackage{pgfplotstable}

%% Code listings
\usepackage{listings}
\usepackage{xcolor}

%% Màu dùng chung (cần định nghĩa trước khi \input chapter)
\definecolor{cppblue}{RGB}{30,100,200}
\definecolor{pygreen}{RGB}{40,150,80}
\definecolor{warnred}{RGB}{200,60,60}
\definecolor{orng}{RGB}{220,130,30}

%% Cài đặt listings mặc định
\lstset{
  basicstyle=\ttfamily\small,
  keywordstyle=\color{cppblue},
  commentstyle=\color{gray!70},
  stringstyle=\color{orng},
  breaklines=true,
  frame=single,
  rulecolor=\color{gray!40},
  framexleftmargin=4pt,
  xleftmargin=8pt,
  numbers=none,
}

%% Hình ảnh
\usepackage{graphicx}

%% Bảng đẹp hơn
\usepackage{booktabs}
\usepackage{multirow}
\usepackage{subcaption}

%% Đảm bảo float không trọn ra ngoài section
\usepackage[section]{placeins}

%%=============================================================
%% Cách sử dụng trong main.tex:
%%
%%   \documentclass[12pt,a4paper]{report}
%%   %%=============================================================
%% preamble_ai3.tex
%% Preamble dùng chung cho chapter4_4.tex và chapter5_4.tex
%%
%% Đưa vào main.tex bằng:  \input{preamble_ai3}
%% Hoặc copy thẳng các \usepackage cần thiết vào preamble của main.
%%=============================================================

%% Tiếng Việt
\usepackage[T5]{fontenc}
\usepackage[utf8]{inputenc}
\usepackage[vietnamese]{babel}

%% TikZ và các thư viện vẻ hình
\usepackage{tikz}
\usetikzlibrary{
  shapes.geometric,
  arrows.meta,
  positioning,
  fit,
  backgrounds,
  calc,
  decorations.pathreplacing,
  patterns,
}

%% PGFPlots cho biểu đồ
\usepackage{pgfplots}
\pgfplotsset{compat=1.18}
\usepackage{pgfplotstable}

%% Code listings
\usepackage{listings}
\usepackage{xcolor}

%% Màu dùng chung (cần định nghĩa trước khi \input chapter)
\definecolor{cppblue}{RGB}{30,100,200}
\definecolor{pygreen}{RGB}{40,150,80}
\definecolor{warnred}{RGB}{200,60,60}
\definecolor{orng}{RGB}{220,130,30}

%% Cài đặt listings mặc định
\lstset{
  basicstyle=\ttfamily\small,
  keywordstyle=\color{cppblue},
  commentstyle=\color{gray!70},
  stringstyle=\color{orng},
  breaklines=true,
  frame=single,
  rulecolor=\color{gray!40},
  framexleftmargin=4pt,
  xleftmargin=8pt,
  numbers=none,
}

%% Hình ảnh
\usepackage{graphicx}

%% Bảng đẹp hơn
\usepackage{booktabs}
\usepackage{multirow}
\usepackage{subcaption}

%% Đảm bảo float không trọn ra ngoài section
\usepackage[section]{placeins}

%%=============================================================
%% Cách sử dụng trong main.tex:
%%
%%   \documentclass[12pt,a4paper]{report}
%%   \input{preamble_ai3}
%%   \begin{document}
%%     ...
%%     \chapter{Thiết kế và Hiện thực}
%%     ...
%%     \input{chapter4_4}   %% mục 4.4
%%     ...
%%     \chapter{Kết quả và Đánh giá}
%%     ...
%%     \input{chapter5_4}   %% mục 5.4
%%     ...
%%   \end{document}
%%
%% Compile: pdflatex -> pdflatex (chạy 2 lần để có label/ref đúng)
%%=============================================================

%%   \begin{document}
%%     ...
%%     \chapter{Thiết kế và Hiện thực}
%%     ...
%%     %%=============================================================
%% chapter4_4.tex
%% Phần 4.4 — AI #3: Phát hiện vật thể cho pick & place
%%
%% Đây là \input file — thêm vào preamble của main.tex:
%%   \usepackage{tikz}
%%   \usetikzlibrary{shapes.geometric,arrows.meta,positioning,
%%     fit,backgrounds,calc,decorations.pathreplacing}
%%   \usepackage{pgfplots}
%%   \pgfplotsset{compat=1.18}
%%   \usepackage{listings,xcolor,booktabs,multirow,subcaption}
%%=============================================================

\definecolor{cppblue}{RGB}{30,100,200}
\definecolor{pygreen}{RGB}{40,150,80}
\definecolor{warnred}{RGB}{200,60,60}
\definecolor{orng}{RGB}{220,130,30}

\lstset{
  basicstyle=\ttfamily\small,
  keywordstyle=\color{cppblue},
  commentstyle=\color{gray!70},
  stringstyle=\color{orng},
  breaklines=true,
  frame=single,
  rulecolor=\color{gray!40},
  framexleftmargin=4pt,
  xleftmargin=8pt,
  numbers=none,
}

% ================================================================
\section{AI \#3 --- Tự động hoá khâu khởi tạo trong hệ thống pick \& place}
\label{sec:ai3}
% ================================================================

% ----------------------------------------------------------------
\subsection{Bối cảnh: vấn đề của hệ thống visual servoing ban đầu}
\label{subsec:ai3-context}
% ----------------------------------------------------------------

Hệ thống pick \& place chạy trên cánh tay robot Denso VS-6577E 6 bậc tự do,
sử dụng thư viện ViSP (Visual Servoing Platform) để điều khiển bằng thị giác.
Camera USB được gắn trực tiếp lên đầu công tác theo cấu hình \textit{eye-in-hand} ---
camera di chuyển cùng robot và nhìn thẳng về phía vật thể cần gắp trong suốt
quá trình tiếp cận.

Vấn đề cốt lõi của hệ thống ban đầu nằm ở bước khởi tạo: sau khi robot về
home pose cố định, \textbf{người vận hành phải click chuột vào vị trí cylinder}
trên cửa sổ hiển thị camera thì blob tracker (\texttt{vpDot2}) mới bắt đầu
hoạt động được. Visual servo chỉ chạy sau thao tác thủ công đó. Ngoài ra,
nếu tracker mất target trong lúc servo --- điều hoàn toàn có thể xảy ra khi
cylinder bị che khuất hoặc di chuyển ra ngoài vùng sáng --- hệ thống reset
hoàn toàn và lại đòi hỏi click thủ công ở lần tiếp theo. Mỗi chu kỳ
pick \& place đều cần sự can thiệp của người vận hành, loại bỏ khả năng
vận hành tự động không giám sát.

\begin{figure}[htbp]
  \centering
  \begin{tikzpicture}[
    node distance=0.8cm and 0.6cm,
    bx/.style={draw, rounded corners=4pt, text width=3.0cm, align=center,
               minimum height=0.9cm, font=\small},
    robot/.style={bx, draw=cppblue!70, fill=cppblue!8},
    manual/.style={bx, draw=warnred!70, fill=warnred!8},
    aibox/.style={bx, draw=pygreen!70, fill=pygreen!8},
    arr/.style={->, >=Stealth, thick},
    rarr/.style={arr, draw=warnred!70},
    garr/.style={arr, draw=pygreen!70},
  ]

  %% --- CỘT TRÁI: hệ thống gốc ---
  \node[robot]                    (L1) {PREINIT\\{\scriptsize robot $\to$ home pose}};
  \node[manual, below=of L1]      (L2) {Click thủ công\\{\scriptsize chọn cylinder}};
  \node[robot,  below=of L2]      (L3) {JOINT\\{\scriptsize vpDot2 tracking}};
  \node[manual, below=of L3]      (L4) {Mất tracking\\{\scriptsize $\to$ reset toàn bộ}};

  \draw[arr]  (L1) -- (L2);
  \draw[arr]  (L2) -- (L3);
  \draw[rarr] (L3) -- (L4);
  \draw[rarr, dashed] (L4.west)
        to[out=190,in=170,looseness=1.3] (L1.west)
        node[midway, left, font=\scriptsize, text=warnred!70] {reset};

  \node[font=\bfseries\small, text=warnred, above=0.5cm of L1]
       {\textit{Trước cải tiến}};

  %% --- CỘT PHẢI: hệ thống mới ---
  \node[robot,  right=3.8cm of L1] (R1) {PREINIT\\{\scriptsize robot $\to$ home pose}};
  \node[aibox, below=of R1]        (R2) {AI detect\\{\scriptsize EfficientDet-Lite0}};
  \node[robot,  below=of R2]       (R3) {JOINT\\{\scriptsize AI tracking mỗi frame}};
  \node[aibox, below=of R3]        (R4) {AI re-detect\\{\scriptsize tự phục hồi}};

  \draw[arr]  (R1) -- (R2);
  \draw[garr] (R2) -- (R3);
  \draw[garr] (R3) -- (R4);
  \draw[garr, dashed] (R4.east)
        to[out=0,in=0,looseness=1.5] (R3.east)
        node[midway, right, font=\scriptsize, text=pygreen!70] {tiếp tục};

  %% fallback arrow từ AI detect
  \node[manual, right=1.2cm of R3] (FB)
       {{\scriptsize fallback:}\\click thủ công};
  \draw[rarr, dashed] (R2.east) -| (FB.north)
        node[near start, above, font=\scriptsize, text=warnred] {AI fail};

  \node[font=\bfseries\small, text=pygreen!80!black, above=0.5cm of R1]
       {\textit{Sau cải tiến}};

  %% Đường phân cách
  \draw[thick, dashed, gray!50]
      ($(L1.east)!0.5!(R1.west)+(0,0.8)$) -- ++(0,-6.5);

  \end{tikzpicture}
  \caption{So sánh luồng vận hành trước và sau khi tích hợp module AI.
           Bước yêu cầu thao tác thủ công (đỏ) được thay thế bằng AI tự động (xanh lá).
           Fallback về click thủ công vẫn được giữ lại cho trường hợp AI thất bại.}
  \label{fig:ai3-before-after}
\end{figure}

Hình~\ref{fig:ai3-before-after} minh hoạ sự thay đổi này. Module AI đảm nhiệm
toàn bộ việc phát hiện cylinder ở bước khởi tạo và tự phục hồi khi mất target
trong quá trình servo. Hệ thống vẫn giữ nguyên fallback về click thủ công để
đảm bảo tương thích ngược.


% ----------------------------------------------------------------
\subsection{Xây dựng và huấn luyện mô hình EfficientDet-Lite0}
\label{subsec:ai3-model}
% ----------------------------------------------------------------

\subsubsection*{Lựa chọn kiến trúc}

Yêu cầu đặt ra khá cụ thể: mô hình phải đủ nhỏ để triển khai được trên
Google Coral EdgeTPU, phải hỗ trợ INT8 quantization, và tương thích với
\texttt{tflite\_model\_maker} --- framework đã có sẵn trong môi trường phát triển.
\textbf{EfficientDet-Lite0} đáp ứng tất cả: đây là biến thể nhỏ nhất trong
họ EfficientDet được Google tối ưu cho thiết bị edge, input cố định
$320 \times 320$ pixel, backbone là EfficientNet-B0 đã được simplify.

\begin{table}[htbp]
  \centering
  \caption{Thông số mô hình EfficientDet-Lite0}
  \label{tab:ai3-model}
  \small
  \begin{tabular}{ll}
    \toprule
    \textbf{Thuộc tính}       & \textbf{Giá trị} \\
    \midrule
    Kiến trúc                 & EfficientDet-Lite0 \\
    Framework training        & TensorFlow~2.8.4 / \texttt{tflite\_model\_maker}~0.4.3 \\
    Input size                & $320 \times 320$~px, uint8 \\
    Số lớp đối tượng          & 2 (cylinder~--~lớp~0, cube~--~lớp~1) \\
    Số epochs / batch size    & 50 / 8 \\
    Transfer learning         & Frozen backbone, fine-tune detection head \\
    Target hardware           & Google Coral USB Accelerator (EdgeTPU) \\
    \bottomrule
  \end{tabular}
\end{table}

Mô hình được huấn luyện để phát hiện \textbf{hai lớp}: cylinder và cube nhựa PLA
in 3D. Việc đưa thêm cube vào tập huấn luyện là có chủ ý --- giúp mô hình học
cách phân biệt hai loại vật thể, từ đó giảm false positive khi cả hai cùng
xuất hiện trong cảnh. Ở phía tích hợp với ViSP, chỉ kết quả cylinder
(\texttt{class\_id~==~0}) được dùng; dự đoán cube bị lọc bỏ.

\subsubsection*{Tập dữ liệu}

Toàn bộ ảnh được thu trực tiếp từ camera eye-in-hand gắn trên robot Denso,
đảm bảo điều kiện thực tế sát với lúc triển khai. Ảnh bao gồm nhiều cấu hình:
góc chụp khác nhau, khoảng cách thay đổi, và số lượng vật thể từ 1 đến 5 trong
một cảnh.

\begin{figure}[htbp]
  \centering
  \begin{tikzpicture}
    \begin{axis}[
      xbar, bar width=0.42cm,
      xlabel={Số lượng ảnh (xấp xỉ)},
      ylabel={Danh mục},
      symbolic y coords={
        Background,
        2 cube + 3 cyl,
        2 cylinders,
        3 cylinders,
        1 cyl + 1 cube,
        1 cube,
        1 cylinder,
        Cylinder (gốc)
      },
      ytick=data,
      xmin=0, xmax=175,
      width=9.5cm, height=6.8cm,
      nodes near coords,
      nodes near coords align={horizontal},
      every node near coord/.style={font=\scriptsize},
      tick label style={font=\small},
      xlabel style={font=\small},
      ylabel style={font=\small},
      enlarge y limits=0.08,
    ]
    \addplot[fill=cppblue!55, draw=cppblue!80] coordinates {
      (38,  Background)
      (28,  2 cube + 3 cyl)
      (52,  2 cylinders)
      (43,  3 cylinders)
      (68,  1 cyl + 1 cube)
      (38,  1 cube)
      (108, 1 cylinder)
      (144, Cylinder (gốc))
    };
    \end{axis}
  \end{tikzpicture}
  \caption{Phân bố tập dữ liệu theo danh mục ảnh (~534 ảnh tổng,
           633 annotation cylinder, 170 annotation cube).
           Các cảnh đa vật thể giúp mô hình học phát hiện từng đối tượng
           riêng lẻ trong môi trường phức tạp.}
  \label{fig:ai3-dataset}
\end{figure}

Annotation được thực hiện bằng \texttt{labelme} theo format bounding box hình chữ
nhật. Workflow:
\begin{enumerate}
  \item Vẽ bounding box cho từng cylinder và cube bằng \texttt{labelme}
  \item Chuyển đổi labelme JSON $\to$ COCO JSON thống nhất
        (script \texttt{labelme\_to\_coco.py})
  \item Tập dữ liệu được split theo tỉ lệ $70/15/15$ (train/val/test),
        seed cố định để đảm bảo tính tái lập
\end{enumerate}

Phân chia ba chiều thay vì hai chiều đảm bảo tập test hoàn toàn độc lập ---
không được nhìn thấy trong suốt quá trình training hoặc điều chỉnh siêu tham số.

\subsubsection*{Pipeline huấn luyện}

\begin{figure}[htbp]
  \centering
  \begin{tikzpicture}[
    node distance=0.45cm and 0.15cm,
    st/.style={draw, rounded corners=3pt, fill=gray!10, draw=gray!55,
               text width=2.0cm, align=center, minimum height=0.85cm, font=\small},
    blue_st/.style={st, fill=cppblue!10, draw=cppblue!55},
    out_st/.style={st, fill=orng!20, draw=orng!70},
    arr/.style={->, >=Stealth, thick, gray!60},
    lbl/.style={font=\scriptsize, text=gray!80},
  ]

  \node[blue_st] (s1) {Annotation\\{\scriptsize labelme}};
  \node[blue_st, right=of s1] (s2) {COCO\\JSON};
  \node[blue_st, right=of s2] (s3) {Pascal\\VOC XML};
  \node[blue_st, right=of s3] (s4) {Training\\{\scriptsize 50 epochs}};
  \node[out_st, right=of s4] (s5) {SavedModel\\{\scriptsize Float32}};
  \node[out_st, below=0.55cm of s5] (s6) {INT8\\TFLite};
  \node[out_st, below=0.55cm of s6] (s7) {EdgeTPU\\TFLite};

  \draw[arr] (s1) -- (s2) node[lbl,above,midway] {to\_coco.py};
  \draw[arr] (s2) -- (s3) node[lbl,above,midway] {tflite\_mm};
  \draw[arr] (s3) -- (s4);
  \draw[arr] (s4) -- (s5) node[lbl,above,midway] {export};
  \draw[arr] (s5) -- (s6) node[lbl,right,midway] {PTQ};
  \draw[arr] (s6) -- (s7) node[lbl,right,midway] {edgetpu\_compiler};

  \end{tikzpicture}
  \caption{Pipeline huấn luyện và xuất mô hình. Từ ảnh thô đến ba file
           mô hình phục vụ các môi trường triển khai khác nhau.
           Bước chuyển đổi sang Pascal VOC XML là yêu cầu bắt buộc của
           \texttt{tflite\_model\_maker}.}
  \label{fig:ai3-pipeline}
\end{figure}

Transfer learning với \texttt{train\_whole\_model=False} --- chỉ fine-tune detection
head, giữ nguyên backbone đã pretrain --- phù hợp với tập dữ liệu tương đối nhỏ
($\sim$374 ảnh train) và giảm nguy cơ overfitting đáng kể so với
huấn luyện từ đầu.


% ----------------------------------------------------------------
\subsection{Quantization INT8 và biên dịch cho EdgeTPU}
\label{subsec:ai3-quant}
% ----------------------------------------------------------------

Mô hình sau training là Float32 --- không chạy được trực tiếp trên Coral EdgeTPU.
EdgeTPU chỉ thực thi mô hình INT8 quantized, trong đó trọng số và activation
được biểu diễn bằng số nguyên 8-bit. Lợi ích: kích thước giảm $\sim4\times$,
và phần cứng EdgeTPU đạt tốc độ inference cực nhanh ($\sim$4--15~ms/frame)
nhờ tận dụng phép tính số nguyên.

Post-Training Quantization (PTQ) sử dụng tập train làm \textit{representative
dataset} để calibrate scale và zero\_point cho từng layer --- đây là bước bắt
buộc để duy trì độ chính xác của INT8 model gần với Float32 baseline.

File \texttt{\_edgetpu.tflite} sau khi biên dịch chứa custom op
\texttt{edgetpu-custom-op} mà TFLite runtime thông thường
\textit{không thể thực thi trên CPU}, gây lỗi:

\begin{lstlisting}[language={}]
RuntimeError: Encountered unresolved custom op: edgetpu-custom-op
\end{lstlisting}

Do đó hệ thống \textbf{duy trì hai file riêng biệt} và lựa chọn tự động:

\begin{center}
\small
\begin{tabular}{lll}
  \toprule
  File & Runtime & Deployment \\
  \midrule
  \texttt{cylinder\_detector\_int8\_edgetpu.tflite} & pycoral   & Coral EdgeTPU \\
  \texttt{cylinder\_detector\_int8.tflite}          & TFLite/ai-edge-litert & CPU \\
  \bottomrule
\end{tabular}
\end{center}


% ----------------------------------------------------------------
\subsection{Kiến trúc tích hợp: C++ ViSP ↔ Python subprocess}
\label{subsec:ai3-arch}
% ----------------------------------------------------------------

\subsubsection*{Quyết định thiết kế}

Thay vì link trực tiếp thư viện TFLite C++ vào build system của ViSP
(đòi hỏi thay đổi lớn vào CMakeLists và phụ thuộc version build cụ thể),
hệ thống áp dụng kiến trúc \textbf{subprocess model}: C++ gọi Python script
như một tiến trình con qua \texttt{popen()}, giao tiếp một chiều qua stdout.
Cách này giúp module AI hoàn toàn tách biệt --- cập nhật mô hình không cần
recompile code robot.

\begin{figure}[htbp]
  \centering
  \begin{tikzpicture}[
    node distance=0.38cm and 0.5cm,
    comp/.style={draw, rounded corners=4pt, text width=3.1cm, align=center,
                 minimum height=0.82cm, font=\small},
    cpp/.style={comp, fill=cppblue!10, draw=cppblue!55},
    py/.style={comp,  fill=pygreen!10, draw=pygreen!60},
    tmp/.style={comp,  fill=orng!18, draw=orng!65, text width=2.5cm,
                minimum height=0.7cm, font=\scriptsize},
    arr/.style={->, >=Stealth, thick},
    ]

  %% ---- C++ side ----
  \node[cpp] (cap)   {Camera capture\\{\scriptsize vpImage{$<$}uchar{$>$}}};
  \node[cpp, below=of cap]  (save)  {Save frame\\{\scriptsize cv::imwrite()}};
  \node[cpp, below=of save] (pop)   {popen() subprocess\\{\scriptsize build cmd}};
  \node[cpp, below=of pop]  (sel)   {select() timeout 3s\\{\scriptsize POSIX wait}};
  \node[cpp, below=of sel]  (parse) {Parse stdout\\{\scriptsize sscanf SUCCESS}};
  \node[cpp, below=of parse](feat)  {vpFeatureBuilder\\{\scriptsize create feature}};

  \draw[arr](cap)--(save); \draw[arr](save)--(pop);
  \draw[arr](pop)--(sel);  \draw[arr](sel)--(parse); \draw[arr](parse)--(feat);

  \node[draw=cppblue!50, dashed, rounded corners=6pt,
        fit=(cap)(save)(pop)(sel)(parse)(feat),
        inner sep=7pt] (cbox) {};
  \node[above=2pt of cbox, font=\bfseries\small, text=cppblue]
       {C++ (ViSP main loop)};

  %% ---- Temp file ----
  \node[tmp, right=2.6cm of save] (jpg)
       {/tmp/visp\_ai\_frame.jpg};

  %% ---- Python side ----
  \node[py, right=2.6cm of pop]   (load) {Load model\\{\scriptsize fallback chain}};
  \node[py, below=of load]        (pre)  {Preprocess\\{\scriptsize resize 320$\times$320}};
  \node[py, below=of pre]         (inf)  {Inference\\{\scriptsize invoke()}};
  \node[py, below=of inf]         (post) {Post-process\\{\scriptsize dequant+filter}};
  \node[py, below=of post]        (out)  {Print stdout\\{\scriptsize SUCCESS u v conf}};

  \draw[arr](load)--(pre); \draw[arr](pre)--(inf);
  \draw[arr](inf)--(post); \draw[arr](post)--(out);

  \node[draw=pygreen!55, dashed, rounded corners=6pt,
        fit=(load)(pre)(inf)(post)(out),
        inner sep=7pt] (pbox) {};
  \node[above=2pt of pbox, font=\bfseries\small, text=pygreen!80!black]
       {Python subprocess};

  %% ---- Data flow arrows ----
  \draw[arr, orng!80]  (save.east) -- (jpg.west)  node[midway,above,font=\scriptsize]{write};
  \draw[arr, orng!80]  (jpg.east)  -- (load.west) node[midway,above,font=\scriptsize]{read};
  \draw[arr, cppblue!60](pop.east)  -- (load.west) node[midway,above,font=\scriptsize,yshift=2pt]{spawn};
  \draw[arr, pygreen!70, bend right=25](out.west)  to (parse.east)
        node[midway,below,font=\scriptsize]{stdout pipe};

  \end{tikzpicture}
  \caption{Kiến trúc tích hợp C++ ViSP $\leftrightarrow$ Python subprocess.
           C++ lưu frame hiện tại ra file tạm, spawn Python làm child process
           qua \texttt{popen()}, đọc kết quả từ stdout sau khi qua timeout guard.}
  \label{fig:ai3-arch}
\end{figure}

\subsubsection*{Protocol trao đổi dữ liệu}

Output của Python script là một dòng duy nhất trên stdout:

\begin{lstlisting}[language={}]
SUCCESS 272.86 190.74 0.8362 238.42 166.89 68.88 47.68
%        u      v      conf   x_min  y_min  width  height
\end{lstlisting}

hoặc \texttt{FAILURE~<lý~do>} nếu không phát hiện được. Tất cả debug log
đi ra \texttt{stderr} --- C++ chỉ đọc stdout --- tránh nhiễu kết quả parse.

Đáng chú ý: trong ViSP, \texttt{vpImagePoint} được định nghĩa là \texttt{(row,~col)}
tương ứng \texttt{(v,~u)}, không phải \texttt{(u,~v)} như thông thường:

\begin{lstlisting}[language=C++]
detected_center = vpImagePoint(v, u);  // row=v, col=u
\end{lstlisting}

Tham số \texttt{hint} cho phép Python chọn cylinder \textit{gần nhất với vị
trí đã biết} thay vì cylinder có confidence cao nhất --- cơ chế này tránh
hiện tượng servo ``nhảy'' sang cylinder khác khi scores thay đổi giữa các frame.


% ----------------------------------------------------------------
\subsection{State machine: PREINIT \texorpdfstring{$\to$}{→} INIT (AI) \texorpdfstring{$\to$}{→} JOINT \texorpdfstring{$\to$}{→} ...}
\label{subsec:ai3-sm}
% ----------------------------------------------------------------

Hệ thống vận hành theo vòng lặp chính liên tục với 6 trạng thái.
Mỗi iteration đọc một frame camera và thực thi logic trạng thái hiện tại.

\begin{figure}[htbp]
  \centering
  \begin{tikzpicture}[
    node distance=1.0cm and 2.0cm,
    st/.style={draw, rounded corners=4pt, fill=cppblue!10, draw=cppblue!55,
               text width=2.3cm, align=center, minimum height=1.0cm, font=\small},
    ai/.style={st, fill=pygreen!10, draw=pygreen!60},
    arr/.style={->, >=Stealth, thick},
    lbl/.style={font=\scriptsize, text=gray!70},
  ]

  %% Top row
  \node[st]  (pre)  {PREINIT\\{\scriptsize home pose}};
  \node[ai,  right=of pre]   (ini)  {INIT\\{\scriptsize AI detect}};
  \node[ai,  right=of ini]   (jnt)  {JOINT\\{\scriptsize AI tracking}};
  \node[st,  right=of jnt]   (app)  {APPROACH\\{\scriptsize chờ RC5}};

  %% Bottom row
  \node[st,  below=1.8cm of app]  (grp)  {GRIPPER\\{\scriptsize đóng gripper}};
  \node[st,  left=of grp]         (cls)  {CLASSIFIED\\{\scriptsize vận chuyển}};

  %% Arrows
  \draw[arr](pre)--(ini)  node[lbl,above,midway]{home\\reached};
  \draw[arr](ini)--(jnt)  node[lbl,above,midway]{AI init\\ok};
  \draw[arr](jnt)--(app)  node[lbl,above,midway]{$|e|<5{\times}10^{-3}$};
  \draw[arr](app)--(grp)  node[lbl,right,midway]{RC5\\signal};
  \draw[arr](grp)--(cls)  node[lbl,below,midway]{grasped};
  \draw[arr](cls.west) to[out=180,in=270,looseness=0.8] (pre.south)
        node[lbl,left,pos=0.5]{cycle\\restart};

  %% AI self-recovery
  \draw[arr, pygreen!70, dashed]
        (jnt.north) to[out=100,in=80,looseness=2.5] (jnt.north)
        node[above=0.7cm of jnt, font=\scriptsize, text=pygreen!80!black]
        {AI re-detect};

  %% Fallback from INIT
  \draw[arr, warnred!70, dashed]
        (ini.south) -- ++(0,-0.7)
        node[below, font=\scriptsize, text=warnred]
        {fallback: click thủ công};

  \end{tikzpicture}
  \caption{State machine của hệ thống pick \& place sau cải tiến.
           Trạng thái tô xanh lá là nơi AI được gọi.
           Vòng lặp tự phục hồi ở JOINT (mũi tên đứt nét) cho phép hệ thống
           tiếp tục mà không cần reset khi mất detection.}
  \label{fig:ai3-sm}
\end{figure}

\paragraph{PREINIT.}
Robot gửi lệnh joint position về home pose
$[0^\circ, 0^\circ, 90^\circ, 0^\circ, 90^\circ, 0^\circ]$.
Servo task được cấu hình với interaction matrix Desired, pseudo-inverse,
gain $\lambda = 0.4$.

\paragraph{INIT.}
Chờ robot đạt home pose (sai số từng joint $< 0.01$) và gripper mở xong.
Khi cả hai điều kiện thỏa, \texttt{detectCylinderWithAI()} được gọi
không có hint --- Python chọn cylinder có confidence cao nhất.
Kết quả trực tiếp tạo visual feature \texttt{vpFeaturePoint}
và add vào servo task. Nếu AI thất bại, \texttt{ai\_hint}
giữ giá trị mặc định (fallback thủ công).

\paragraph{JOINT.}
Khác với cách tiếp cận truyền thống dùng blob tracker mỗi frame,
hệ thống hiện tại \textbf{gọi AI trực tiếp ở mỗi iteration}
với \texttt{hint = \&ai\_hint}. Control law:
\begin{equation}
  \mathbf{v} = -\lambda\,\mathbf{L}_{s^*}^{+}\,\mathbf{e},
  \qquad \mathbf{q}_{k+1} = \mathbf{q}_k + \mathbf{v}\,\Delta t
  \label{eq:ibvs}
\end{equation}
với $\Delta t = 0.5\,$s. Điều kiện hội tụ: $|e_x| < 5{\times}10^{-3}$
và $|e_y| < 5{\times}10^{-3}$ (normalized image coordinates).
Khi đạt, gửi \texttt{"OKE\textbackslash r"} qua UART đến RC5 controller.


% ----------------------------------------------------------------
\subsection{Các vấn đề kỹ thuật và giải pháp}
\label{subsec:ai3-issues}
% ----------------------------------------------------------------

Một số vấn đề nảy sinh trong quá trình tích hợp mà khó thấy trước khi thực
sự chạy hệ thống.

\subsubsection*{P1: Xác định output tensor sai lệch}

\texttt{tflite\_model\_maker} 0.4.x đặt tên tensor theo thứ tự riêng của nó:

\begin{center}
\small
\begin{tabular}{ccc}
  \toprule
  Tensor & Shape & Nội dung \\
  \midrule
  \texttt{:0} & \texttt{[1]}     & num\_detections \\
  \texttt{:1} & \texttt{[1,\,N]} & scores \quad \textbf{← cần lấy cái này} \\
  \texttt{:2} & \texttt{[1,\,N]} & class IDs \quad \textbf{← không phải cái này} \\
  \texttt{:3} & \texttt{[1,\,N,\,4]} & boxes \\
  \bottomrule
\end{tabular}
\end{center}

Thứ tự này không khớp với API của pycoral, vốn giả định positional order
\texttt{boxes/classes/scores/count = 0/1/2/3}. Dùng API đó sẽ đọc
num\_detections (raw uint8 = 255) như scores và crash với \texttt{IndexError}.
Giải pháp: xác định tensor \textbf{theo shape động} --- boxes là
\texttt{[1,N,4]}, scores là tensor \texttt{[1,N]} có name suffix nhỏ nhất
(\texttt{:1} chứ không phải \texttt{:2}).

\subsubsection*{P2: Python 3.10 không có \texttt{tflite-runtime}}

Google không publish wheel \texttt{tflite-runtime} cho Python 3.10+.
Giải pháp là xây dựng \textbf{fallback chain 4 tầng}:

\begin{figure}[htbp]
  \centering
  \begin{tikzpicture}[
    node distance=0.4cm,
    fb/.style={draw, rounded corners=3pt, text width=7.5cm, align=left,
               minimum height=0.85cm, font=\small, inner sep=7pt},
    t1/.style={fb, fill=purple!12, draw=purple!55},
    t2/.style={fb, fill=teal!10,   draw=teal!50},
    t3/.style={fb, fill=cppblue!10, draw=cppblue!50},
    t4/.style={fb, fill=gray!12,   draw=gray!50},
    fa/.style={fb, fill=warnred!10, draw=warnred!50},
    arr/.style={->, >=Stealth, thick, dashed, gray!55},
    lbl/.style={font=\scriptsize, text=warnred!80},
  ]
  \node[t1] (a) {\textbf{1.}~\texttt{pycoral.utils.edgetpu}\quad{\scriptsize Coral EdgeTPU, Python 3.9}};
  \node[t2, below=of a] (b) {\textbf{2.}~\texttt{ai\_edge\_litert.interpreter}\quad{\scriptsize Google's new pkg, Python 3.8--3.12}};
  \node[t3, below=of b] (c) {\textbf{3.}~\texttt{tflite\_runtime.interpreter}\quad{\scriptsize Legacy, Python 3.9}};
  \node[t4, below=of c] (d) {\textbf{4.}~\texttt{tensorflow.lite.Interpreter}\quad{\scriptsize Heavy fallback}};
  \node[fa, below=of d] (e) {\textbf{FAILURE}:~\texttt{no\_tflite\_backend\_available}};

  \draw[arr](a)--(b) node[lbl,right,midway]{ImportError / Coral not connected};
  \draw[arr](b)--(c) node[lbl,right,midway]{ImportError};
  \draw[arr](c)--(d) node[lbl,right,midway]{ImportError};
  \draw[arr](d)--(e) node[lbl,right,midway]{ImportError};
  \end{tikzpicture}
  \caption{Python backend fallback chain trong \texttt{detect\_cylinder.py}.
           Script tự động thử từng backend theo thứ tự ưu tiên,
           tương thích với cả môi trường có Coral lẫn CPU-only.}
  \label{fig:ai3-fallback}
\end{figure}

\texttt{ai-edge-litert} là gói chính thức của Google thay thế \texttt{tflite-runtime},
hỗ trợ Python 3.8--3.12 với API giống hệt nhau --- là giải pháp cho vấn đề Python 3.10.

\subsubsection*{P3: JSON comment trong config}

\texttt{config.json} dùng comment \texttt{//} để ghi chú, nhưng cả Python
\texttt{json.load()} và \texttt{nlohmann::json} đều từ chối mặc định.
Phía C++: \texttt{nlohmann::json::parse(file, nullptr, true, true)} với
\texttt{ignore\_comments=true} (có từ nlohmann/json~3.9.0).
Phía Python: strip comment lines thủ công trước khi \texttt{json.loads()}.


% ----------------------------------------------------------------
\subsection{Thông số cấu hình visual servo}
\label{subsec:ai3-servo}
% ----------------------------------------------------------------

Hệ thống sử dụng \textbf{Image-Based Visual Servoing (IBVS)} với feature là
tọa độ 2D của cylinder trong không gian ảnh normalized $\mathbf{s}=(x,y)$,
desired position $\mathbf{s}^*=(0,0,1)$ --- cylinder phải nằm đúng tâm ảnh.

\begin{table}[htbp]
  \centering
  \caption{Thông số điều khiển visual servo}
  \label{tab:ai3-servo}
  \small
  \begin{tabular}{ll}
    \toprule
    \textbf{Tham số}            & \textbf{Giá trị} \\
    \midrule
    Loại điều khiển             & Eye-in-hand, \texttt{EYEINHAND\_L\_cVe\_eJe} \\
    Interaction matrix          & Desired ($\mathbf{L}_{s^*}$), pseudo-inverse \\
    Gain $\lambda$              & 0.4 \\
    Bước thời gian $\Delta t$   & 0.5\,s \\
    Ngưỡng hội tụ               & $|e_x| < 5{\times}10^{-3}$ và $|e_y| < 5{\times}10^{-3}$ \\
    Desired position            & $(x^*, y^*) = (0,\,0)$ --- tâm ảnh, $Z = 1$ \\
    \bottomrule
  \end{tabular}
\end{table}

Dùng $\mathbf{L}_{s^*}$ (tại vị trí mong muốn) thay vì $\mathbf{L}_{s_{cur}}$
giúp tránh singularity khi feature tiến gần desired point và giữ ổn định
theo ổn định hơn trong giai đoạn hội tụ. Velocity được saturate trước khi
gửi xuống robot để đảm bảo không vượt giới hạn vật lý từng khớp.
   %% mục 4.4
%%     ...
%%     \chapter{Kết quả và Đánh giá}
%%     ...
%%     %%=============================================================
%% chapter5_4.tex
%% Phần 5.4 — Kết quả đánh giá: AI #3 Phát hiện vật thể
%%
%% \input file — cần thêm vào preamble:
%%   \usepackage{tikz}
%%   \usetikzlibrary{arrows.meta,positioning,patterns}
%%   \usepackage{pgfplots}
%%   \pgfplotsset{compat=1.18}
%%   \usepackage{booktabs,multirow,subcaption,xcolor}
%%=============================================================

% ================================================================
\section{Kết quả và Đánh giá: AI \#3 --- Phát hiện vật thể}
\label{sec:ai3-results}
% ================================================================

% ----------------------------------------------------------------
\subsection{Kết quả detection model: COCO metrics}
\label{subsec:ai3-coco}
% ----------------------------------------------------------------

Mô hình được đánh giá trên cả validation set và test set độc lập
theo chuẩn \textbf{COCO evaluation protocol} --- tính trung bình mAP qua
các ngưỡng IoU từ 0.50 đến 0.95 với bước 0.05.

\begin{figure}[htbp]
  \centering
  \begin{tikzpicture}
    \begin{axis}[
      ybar, bar width=0.45cm,
      width=10cm, height=7cm,
      symbolic x coords={mAP (0.50:0.95), AP50, AP75, AR},
      xtick=data,
      x tick label style={font=\small, align=center},
      ymin=0, ymax=1.15,
      ylabel={Score},
      ylabel style={font=\small},
      legend style={at={(0.5,1.03)}, anchor=south, legend columns=2, font=\small},
      legend cell align=left,
      nodes near coords,
      every node near coord/.style={font=\scriptsize, /pgf/number format/fixed,
                                    /pgf/number format/precision=3},
      enlarge x limits=0.25,
      grid=major, grid style={dashed, gray!30},
      axis line style={gray!60},
      tick style={gray!60},
    ]
    \addplot[fill=cppblue!70, draw=cppblue!90] coordinates {
      (mAP (0.50:0.95), 0.791)
      (AP50, 0.990)
      (AP75, 0.990)
      (AR,   0.869)
    };
    \addplot[fill=orng!80, draw=orng!90] coordinates {
      (mAP (0.50:0.95), 0.799)
      (AP50, 1.000)
      (AP75, 0.969)
      (AR,   0.856)
    };
    \legend{Validation, Test}
    \end{axis}
  \end{tikzpicture}
  \caption{COCO metrics trên validation set và test set. AP50 trên test set
           đạt 1.000 cho thấy 100\% cylinder được phát hiện đúng ở ngưỡng IoU
           lỏng lẻ (50\%). mAP$\approx$0.80 phản ánh vẫn còn đây đó những trường hợp
           định vị bounding box chưa được chính xác ở ngưỡng IoU cao.}
  \label{fig:ai3-coco}
\end{figure}

Bảng~\ref{tab:ai3-metrics} tổng hợp số liệu chi tiết:

\begin{table}[htbp]
  \centering
  \caption{COCO metrics trên validation và test set}
  \label{tab:ai3-metrics}
  \small
  \begin{tabular}{lcc}
    \toprule
    \textbf{Metric}                    & \textbf{Validation} & \textbf{Test} \\
    \midrule
    mAP (IoU 0.50:0.95)               & 0.791               & \textbf{0.799} \\
    AP50 (IoU $\geq$ 0.50)            & 0.990               & \textbf{1.000} \\
    AP75 (IoU $\geq$ 0.75)            & 0.990               & 0.969 \\
    AR (maxDets=100)                   & 0.869               & 0.856 \\
    \midrule
    Precision (IoU=0.50, conf=0.50)    & ---                 & \textbf{1.000} \\
    Recall    (IoU=0.50, conf=0.50)    & ---                 & 0.990 \\
    F1-Score                           & ---                 & \textbf{0.995} \\
    TP / FP / FN                       & ---                 & 95 / 0 / 1 \\
    \bottomrule
  \end{tabular}
\end{table}

\begin{figure}[htbp]
  \centering
  \begin{tikzpicture}
    \begin{axis}[
      ybar, bar width=0.55cm,
      width=8.5cm, height=6cm,
      symbolic x coords={Precision, Recall, F1-Score},
      xtick=data,
      x tick label style={font=\small},
      ymin=0.96, ymax=1.04,
      ylabel={Score},
      ylabel style={font=\small},
      nodes near coords,
      every node near coord/.style={font=\small, /pgf/number format/fixed,
                                    /pgf/number format/precision=3},
      enlarge x limits=0.3,
      grid=major, grid style={dashed, gray!30},
      axis line style={gray!60},
    ]
    \addplot[fill=pygreen!65, draw=pygreen!80] coordinates {
      (Precision, 1.000)
      (Recall,    0.990)
      (F1-Score,  0.995)
    };
    \end{axis}
  \end{tikzpicture}
  \caption{Precision, Recall và F1-Score trên test set (PASCAL VOC protocol,
           IoU$\geq$0.50, confidence$\geq$0.50, chỉ lớp cylinder).
           Precision = 1.000 và FP = 0 chứng tỏ không có false positive
           nào ở ngưỡng 0.50. FN = 1 (một cylinder bị bỏ sót) giải thích
           Recall = 0.990.}
  \label{fig:ai3-prf}
\end{figure}

Giá trị AP50 = 1.000 trên test set à ngưỡng IoU lỏng lẻ (50\%) phản ánh
mô hình phát hiện được toàn bộ cylinder. mAP = 0.799 (trung bình
qua IoU 0.50--0.95) thấp hơn một chút do mô hình EfficientDet-Lite0
đơn giản vẫn còn giới hạn trong việc định vị bounding box chính xác
ở các ngưỡng IoU cao ($\geq$0.75) --- điều phổ biến với các kiến trúc
nhe hơn.


% ----------------------------------------------------------------
\subsection{Phân tích confidence score: true positive vs false positive}
\label{subsec:ai3-conf}
% ----------------------------------------------------------------

Một trong những điều quan trọng khi triển khai thực tế là khoảng cách
phân tách giữa true positive và false positive --- điều này quyết định
ngưỡng confidence nào thực tế dùng được mà không sai kết quả.

Kiểm tra trên 218 ảnh cylinder (true positives) và 10 ảnh cube
(false positives):

\begin{figure}[htbp]
  \centering
  \begin{tikzpicture}
    \begin{axis}[
      width=11cm, height=5.5cm,
      xlabel={Confidence score},
      ylabel={},
      xmin=0, xmax=1.05,
      ymin=0, ymax=3,
      xtick={0,0.1,...,1.0},
      ytick={0.7, 2.0},
      yticklabels={\small False Positive\\(cube, $n$=10), \small True Positive\\(cylinder, $n$=218)},
      y tick label style={align=right, font=\small},
      x tick label style={font=\small},
      xlabel style={font=\small},
      grid=both, grid style={dashed, gray!25},
      axis line style={gray!60},
    ]

    %% --- True Positive band (y=2) ---
    %% Most scores: 0.85-0.98 (majority)
    \addplot[only marks, mark=|, mark size=5pt, mark options={draw=pygreen!70, line width=1.2pt},
             opacity=0.5]
    coordinates {
      (0.29,2)(0.61,2)(0.74,2)
      (0.82,2)(0.83,2)(0.84,2)(0.85,2)(0.86,2)(0.87,2)(0.88,2)
      (0.89,2)(0.90,2)(0.91,2)(0.92,2)(0.93,2)(0.94,2)(0.94,2)
      (0.95,2)(0.95,2)(0.96,2)(0.96,2)(0.97,2)(0.97,2)(0.98,2)
    };

    %% --- False Positive band (y=0.7) ---
    \addplot[only marks, mark=|, mark size=5pt, mark options={draw=warnred!70, line width=1.2pt},
             opacity=0.6]
    coordinates {
      (0.08,0.7)(0.13,0.7)(0.17,0.7)(0.21,0.7)
      (0.25,0.7)(0.28,0.7)(0.30,0.7)(0.33,0.7)(0.35,0.7)(0.36,0.7)
    };

    %% --- Threshold line at 0.50 ---
    \addplot[thick, dashed, draw=gray!80, domain=0:3]
    coordinates {(0.50, 0) (0.50, 3)};
    \node[font=\scriptsize, text=gray!80, rotate=90]
         at (axis cs: 0.52, 1.5) {threshold = 0.50};

    %% --- Gap annotation ---
    \draw[->, >=Stealth, thick, gray!70]
         (axis cs:0.38, 2.5) -- (axis cs:0.48, 2.5);
    \draw[->, >=Stealth, thick, gray!70]
         (axis cs:0.34, 2.5) -- (axis cs:0.27, 2.5);
    \node[font=\scriptsize, text=gray!80] at (axis cs:0.37, 2.7) {gap = 0.34};

    %% --- Median annotation for TP ---
    \addplot[only marks, mark=diamond*, mark size=4pt,
             mark options={draw=pygreen!80!black, fill=pygreen!60}]
    coordinates {(0.94, 2)};
    \node[font=\scriptsize, text=pygreen!80!black, above]
         at (axis cs:0.94, 2) {median 0.94};

    %% --- Max FP annotation ---
    \addplot[only marks, mark=diamond*, mark size=4pt,
             mark options={draw=warnred!80, fill=warnred!50}]
    coordinates {(0.36, 0.7)};
    \node[font=\scriptsize, text=warnred, below]
         at (axis cs:0.36, 0.7) {max 0.36};

    \end{axis}
  \end{tikzpicture}
  \caption{Phân tách confidence score giữa true positive (cylinder, xanh lá)
           và false positive (cube, đỏ). Mỗi vạch đứng đại diện một detection.
           Ngưỡng 0.50 nằm chính giữa khoảng gap 0.34,
           loại bỏ toàn bộ false positive mà vẫn chấp nhận 99.5\% true positive.}
  \label{fig:ai3-conf}
\end{figure}

Khoảng cách 0.34 giữa điểm false positive cao nhất (0.36) và 99.1\%
true positive ($\geq$0.70) xác nhận mô hình đã học được đặc trưng
phân biệt mạnh giữa cylinder và cube. Ngưỡng hiện tại 0.50 nằm an toàn
giữa hai vùng phân bố này.

Ba ảnh outlier có score thấp (0.29, 0.61, 0.74) tương ứng với cylinder
bị nhìn từ góc bất thường hoặc bị che khuất một phần. Ỡ runtime,
các trường hợp này fallback về click thủ công --- hành vi an toàn.

\begin{table}[htbp]
  \centering
  \caption{Phân tích confidence score trên 218 ảnh cylinder + 10 ảnh cube}
  \label{tab:ai3-conf}
  \small
  \begin{tabular}{lcc}
    \toprule
    & \textbf{True Positive (cylinder)} & \textbf{False Positive (cube)} \\
    \midrule
    Score range      & 0.29 -- 0.98          & 0.08 -- 0.36 \\
    Median           & \textbf{0.94}         & --- \\
    $\geq$ 0.50      & 99.5\% (217/218)      & \textbf{0/10 (0\%)} \\
    $\geq$ 0.70      & 99.1\% (216/218)      & 0/10 (0\%) \\
    $\geq$ 0.80      & 95.9\% (209/218)      & 0/10 (0\%) \\
    \midrule
    Gap ph\^an t\'ach & \multicolumn{2}{c}{$0.70 - 0.36 = 0.34$} \\
    \bottomrule
  \end{tabular}
\end{table}


% ----------------------------------------------------------------
\subsection{Kết quả tích hợp hệ thống}
\label{subsec:ai3-integration}
% ----------------------------------------------------------------

\begin{table}[htbp]
  \centering
  \caption{So sánh khả năng trước và sau khi tích hợp AI}
  \label{tab:ai3-comparison}
  \small
  \begin{tabular}{lll}
    \toprule
    \textbf{Tiêu chí}                    & \textbf{Hệ thống gốc} & \textbf{Sau cải tiến} \\
    \midrule
    Khởi tạo mỗi chu kỳ               & Click thủ công (bắt buộc) & Tự động (AI) \\
    Tỷ lệ init thành công             & 100\% (do manual)  & $\sim$99.5\% (AI) \\
    Fallback khi AI thất bại           & Không có              & Click thủ công \\
    Phục hồi khi mất tracking          & Reset toàn bộ        & AI re-detect tự động \\
    Chạy không giám sát                & Không thể              & Có (trường hợp bình thường) \\
    \bottomrule
  \end{tabular}
\end{table}

Thiết kế hiện tại gọi AI trực tiếp \textbf{mỗi iteration} của servo
loop thay vì dùng blob tracker. Điều này có một hàm ý đáng chú ý:
\textit{không có tracking failure} theo nghĩa truyền thống --- mỗi frame
là một lần detection độc lập. Tính ổn định của servo phụ thuộc vào
sự nhất quán của AI detection qua các frame liên tiếp. Cơ chế \texttt{--hint}
(truyền vị trí frame trước) giúp được phần nào khi có nhiều cylinder
trong cảnh.


% ----------------------------------------------------------------
\subsection{Phân tích latency: subprocess CPU vs EdgeTPU}
\label{subsec:ai3-latency}
% ----------------------------------------------------------------

Mỗi lần gọi \texttt{detectCylinderWithAI()} bao gồm nhiều thành phần
latency khác nhau. Điều ít người để ý là: phần overhead lớn nhất
không phải là inference mà là \textbf{Python startup time} --- chi phí
mở một process mới mỗi lần gọi \texttt{popen()}.

\begin{figure}[htbp]
  \centering
  \begin{tikzpicture}
    \begin{axis}[
      xbar stacked,
      bar width=0.55cm,
      width=11cm, height=6.5cm,
      symbolic y coords={
        {EdgeTPU\\(persistent, tương lai)},
        {CPU\\(persistent, tương lai)},
        {EdgeTPU\\(subprocess hiện tại)},
        {CPU\\(subprocess hiện tại)}
      },
      ytick=data,
      y tick label style={align=right, font=\small},
      xlabel={Thời gian (ms)},
      xlabel style={font=\small},
      xmin=0, xmax=780,
      legend style={at={(0.5,-0.18)}, anchor=north, legend columns=4,
                    font=\scriptsize},
      legend cell align=left,
      grid=major, grid style={dashed, gray!25},
      axis line style={gray!60},
      enlarge y limits=0.18,
      nodes near coords, nodes near coords align={horizontal},
      point meta=explicit symbolic,
      every node near coord/.style={font=\scriptsize},
    ]
    %%  stacks: I/O | Inference | Model load | Python startup
    %% CPU subprocess: 10 + 150 + 100 + 350 = 610
    \addplot[fill=gray!40, draw=gray!60] coordinates {
      (10, {CPU\\(subprocess hiện tại)})
      (10, {EdgeTPU\\(subprocess hiện tại)})
      (10, {CPU\\(persistent, tương lai)})
      (5,  {EdgeTPU\\(persistent, tương lai)})
    };  %% I/O
    \addplot[fill=warnred!55, draw=warnred!70] coordinates {
      (150, {CPU\\(subprocess hiện tại)})
      (10,  {EdgeTPU\\(subprocess hiện tại)})
      (150, {CPU\\(persistent, tương lai)})
      (10,  {EdgeTPU\\(persistent, tương lai)})
    };  %% Inference
    \addplot[fill=orng!65, draw=orng!80] coordinates {
      (100, {CPU\\(subprocess hiện tại)})
      (50,  {EdgeTPU\\(subprocess hiện tại)})
      (0,   {CPU\\(persistent, tương lai)})
      (0,   {EdgeTPU\\(persistent, tương lai)})
    };  %% Model load
    \addplot[fill=cppblue!55, draw=cppblue!70] coordinates {
      (350, {CPU\\(subprocess hiện tại)})
      (350, {EdgeTPU\\(subprocess hiện tại)})
      (0,   {CPU\\(persistent, tương lai)})
      (0,   {EdgeTPU\\(persistent, tương lai)})
    };  %% Python startup
    \legend{I/O (JPEG), Inference, Model load, Python startup}
    \end{axis}
  \end{tikzpicture}
  \caption{Phân tích latency mỗi lần gọi AI detection.
           Python startup time ($\sim$350~ms) chiếm phần lẫn tổng overhead
           trong kiến trúc subprocess hiện tại --- lớn hơn nhiều so với
           thời gian inference thực tế. Kiến trúc persistent process (tương lai)
           loại bỏ hoàn toàn chi phí này.}
  \label{fig:ai3-latency}
\end{figure}

\begin{table}[htbp]
  \centering
  \caption{So sánh tổng latency giữa các cấu hình triển khai}
  \label{tab:ai3-latency}
  \small
  \begin{tabular}{lcc}
    \toprule
    \textbf{Cấu hình}                    & \textbf{Inference} & \textbf{Tổng / lần gọi} \\
    \midrule
    CPU subprocess (hiện tại)            & $\sim$100--200~ms  & $\sim$600--700~ms \\
    EdgeTPU subprocess (hiện tại)       & $\sim$4--15~ms     & $\sim$420--430~ms \\
    CPU persistent (tương lai)           & $\sim$100--200~ms  & $\sim$160--210~ms \\
    EdgeTPU persistent (tương lai)      & $\sim$4--15~ms     & $\sim$15--20~ms  \\
    \bottomrule
  \end{tabular}
\end{table}

Kết quả trên cho thấy rõ điểm nghẻn hiện tại:
EdgeTPU inference chỉ mất 4--15~ms nhưng toàn bộ lá gọi
vẫn tốn 420~ms do phải khởi động Python mỗi lần. Với
kiến trúc \textbf{persistent Python process} --- một process
duy nhất khởi động một lần, load mô hình vào bộ nhớ, sau
đó tiếp nhận request từ C++ qua stdin/stdout pipe --- latency
mỗi lần gọi giảm xuống còn $\sim$15--20~ms với EdgeTPU,
mở ra khả năng chạy AI detection \textit{mỗi frame} camera
thay vì chỉ khi cần.

% Chú thích về nguồn số liệu
\smallskip
\noindent\textit{\small
Lưu ý: COCO metrics trong Bảng~\ref{tab:ai3-metrics} và Hình~\ref{fig:ai3-coco}
lấy từ kết quả đánh giá hai lớp (cylinder + cube) trên repository huấn luyện.
Phân tích confidence score trong Bảng~\ref{tab:ai3-conf} lấy từ kiểm tra
độc lập trên 218 ảnh cylinder và 10 ảnh cube (CPU inference via ai-edge-litert).
}
   %% mục 5.4
%%     ...
%%   \end{document}
%%
%% Compile: pdflatex -> pdflatex (chạy 2 lần để có label/ref đúng)
%%=============================================================

%% Hoặc copy thẳng các \usepackage cần thiết vào preamble của main.
%%=============================================================

%% Tiếng Việt
\usepackage[T5]{fontenc}
\usepackage[utf8]{inputenc}
\usepackage[vietnamese]{babel}

%% TikZ và các thư viện vẻ hình
\usepackage{tikz}
\usetikzlibrary{
  shapes.geometric,
  arrows.meta,
  positioning,
  fit,
  backgrounds,
  calc,
  decorations.pathreplacing,
  patterns,
}

%% PGFPlots cho biểu đồ
\usepackage{pgfplots}
\pgfplotsset{compat=1.18}
\usepackage{pgfplotstable}

%% Code listings
\usepackage{listings}
\usepackage{xcolor}

%% Màu dùng chung (cần định nghĩa trước khi \input chapter)
\definecolor{cppblue}{RGB}{30,100,200}
\definecolor{pygreen}{RGB}{40,150,80}
\definecolor{warnred}{RGB}{200,60,60}
\definecolor{orng}{RGB}{220,130,30}

%% Cài đặt listings mặc định
\lstset{
  basicstyle=\ttfamily\small,
  keywordstyle=\color{cppblue},
  commentstyle=\color{gray!70},
  stringstyle=\color{orng},
  breaklines=true,
  frame=single,
  rulecolor=\color{gray!40},
  framexleftmargin=4pt,
  xleftmargin=8pt,
  numbers=none,
}

%% Hình ảnh
\usepackage{graphicx}

%% Bảng đẹp hơn
\usepackage{booktabs}
\usepackage{multirow}
\usepackage{subcaption}

%% Đảm bảo float không trọn ra ngoài section
\usepackage[section]{placeins}

%%=============================================================
%% Cách sử dụng trong main.tex:
%%
%%   \documentclass[12pt,a4paper]{report}
%%   %%=============================================================
%% preamble_ai3.tex
%% Preamble dùng chung cho chapter4_4.tex và chapter5_4.tex
%%
%% Đưa vào main.tex bằng:  %%=============================================================
%% preamble_ai3.tex
%% Preamble dùng chung cho chapter4_4.tex và chapter5_4.tex
%%
%% Đưa vào main.tex bằng:  \input{preamble_ai3}
%% Hoặc copy thẳng các \usepackage cần thiết vào preamble của main.
%%=============================================================

%% Tiếng Việt
\usepackage[T5]{fontenc}
\usepackage[utf8]{inputenc}
\usepackage[vietnamese]{babel}

%% TikZ và các thư viện vẻ hình
\usepackage{tikz}
\usetikzlibrary{
  shapes.geometric,
  arrows.meta,
  positioning,
  fit,
  backgrounds,
  calc,
  decorations.pathreplacing,
  patterns,
}

%% PGFPlots cho biểu đồ
\usepackage{pgfplots}
\pgfplotsset{compat=1.18}
\usepackage{pgfplotstable}

%% Code listings
\usepackage{listings}
\usepackage{xcolor}

%% Màu dùng chung (cần định nghĩa trước khi \input chapter)
\definecolor{cppblue}{RGB}{30,100,200}
\definecolor{pygreen}{RGB}{40,150,80}
\definecolor{warnred}{RGB}{200,60,60}
\definecolor{orng}{RGB}{220,130,30}

%% Cài đặt listings mặc định
\lstset{
  basicstyle=\ttfamily\small,
  keywordstyle=\color{cppblue},
  commentstyle=\color{gray!70},
  stringstyle=\color{orng},
  breaklines=true,
  frame=single,
  rulecolor=\color{gray!40},
  framexleftmargin=4pt,
  xleftmargin=8pt,
  numbers=none,
}

%% Hình ảnh
\usepackage{graphicx}

%% Bảng đẹp hơn
\usepackage{booktabs}
\usepackage{multirow}
\usepackage{subcaption}

%% Đảm bảo float không trọn ra ngoài section
\usepackage[section]{placeins}

%%=============================================================
%% Cách sử dụng trong main.tex:
%%
%%   \documentclass[12pt,a4paper]{report}
%%   \input{preamble_ai3}
%%   \begin{document}
%%     ...
%%     \chapter{Thiết kế và Hiện thực}
%%     ...
%%     \input{chapter4_4}   %% mục 4.4
%%     ...
%%     \chapter{Kết quả và Đánh giá}
%%     ...
%%     \input{chapter5_4}   %% mục 5.4
%%     ...
%%   \end{document}
%%
%% Compile: pdflatex -> pdflatex (chạy 2 lần để có label/ref đúng)
%%=============================================================

%% Hoặc copy thẳng các \usepackage cần thiết vào preamble của main.
%%=============================================================

%% Tiếng Việt
\usepackage[T5]{fontenc}
\usepackage[utf8]{inputenc}
\usepackage[vietnamese]{babel}

%% TikZ và các thư viện vẻ hình
\usepackage{tikz}
\usetikzlibrary{
  shapes.geometric,
  arrows.meta,
  positioning,
  fit,
  backgrounds,
  calc,
  decorations.pathreplacing,
  patterns,
}

%% PGFPlots cho biểu đồ
\usepackage{pgfplots}
\pgfplotsset{compat=1.18}
\usepackage{pgfplotstable}

%% Code listings
\usepackage{listings}
\usepackage{xcolor}

%% Màu dùng chung (cần định nghĩa trước khi \input chapter)
\definecolor{cppblue}{RGB}{30,100,200}
\definecolor{pygreen}{RGB}{40,150,80}
\definecolor{warnred}{RGB}{200,60,60}
\definecolor{orng}{RGB}{220,130,30}

%% Cài đặt listings mặc định
\lstset{
  basicstyle=\ttfamily\small,
  keywordstyle=\color{cppblue},
  commentstyle=\color{gray!70},
  stringstyle=\color{orng},
  breaklines=true,
  frame=single,
  rulecolor=\color{gray!40},
  framexleftmargin=4pt,
  xleftmargin=8pt,
  numbers=none,
}

%% Hình ảnh
\usepackage{graphicx}

%% Bảng đẹp hơn
\usepackage{booktabs}
\usepackage{multirow}
\usepackage{subcaption}

%% Đảm bảo float không trọn ra ngoài section
\usepackage[section]{placeins}

%%=============================================================
%% Cách sử dụng trong main.tex:
%%
%%   \documentclass[12pt,a4paper]{report}
%%   %%=============================================================
%% preamble_ai3.tex
%% Preamble dùng chung cho chapter4_4.tex và chapter5_4.tex
%%
%% Đưa vào main.tex bằng:  \input{preamble_ai3}
%% Hoặc copy thẳng các \usepackage cần thiết vào preamble của main.
%%=============================================================

%% Tiếng Việt
\usepackage[T5]{fontenc}
\usepackage[utf8]{inputenc}
\usepackage[vietnamese]{babel}

%% TikZ và các thư viện vẻ hình
\usepackage{tikz}
\usetikzlibrary{
  shapes.geometric,
  arrows.meta,
  positioning,
  fit,
  backgrounds,
  calc,
  decorations.pathreplacing,
  patterns,
}

%% PGFPlots cho biểu đồ
\usepackage{pgfplots}
\pgfplotsset{compat=1.18}
\usepackage{pgfplotstable}

%% Code listings
\usepackage{listings}
\usepackage{xcolor}

%% Màu dùng chung (cần định nghĩa trước khi \input chapter)
\definecolor{cppblue}{RGB}{30,100,200}
\definecolor{pygreen}{RGB}{40,150,80}
\definecolor{warnred}{RGB}{200,60,60}
\definecolor{orng}{RGB}{220,130,30}

%% Cài đặt listings mặc định
\lstset{
  basicstyle=\ttfamily\small,
  keywordstyle=\color{cppblue},
  commentstyle=\color{gray!70},
  stringstyle=\color{orng},
  breaklines=true,
  frame=single,
  rulecolor=\color{gray!40},
  framexleftmargin=4pt,
  xleftmargin=8pt,
  numbers=none,
}

%% Hình ảnh
\usepackage{graphicx}

%% Bảng đẹp hơn
\usepackage{booktabs}
\usepackage{multirow}
\usepackage{subcaption}

%% Đảm bảo float không trọn ra ngoài section
\usepackage[section]{placeins}

%%=============================================================
%% Cách sử dụng trong main.tex:
%%
%%   \documentclass[12pt,a4paper]{report}
%%   \input{preamble_ai3}
%%   \begin{document}
%%     ...
%%     \chapter{Thiết kế và Hiện thực}
%%     ...
%%     \input{chapter4_4}   %% mục 4.4
%%     ...
%%     \chapter{Kết quả và Đánh giá}
%%     ...
%%     \input{chapter5_4}   %% mục 5.4
%%     ...
%%   \end{document}
%%
%% Compile: pdflatex -> pdflatex (chạy 2 lần để có label/ref đúng)
%%=============================================================

%%   \begin{document}
%%     ...
%%     \chapter{Thiết kế và Hiện thực}
%%     ...
%%     %%=============================================================
%% chapter4_4.tex
%% Phần 4.4 — AI #3: Phát hiện vật thể cho pick & place
%%
%% Đây là \input file — thêm vào preamble của main.tex:
%%   \usepackage{tikz}
%%   \usetikzlibrary{shapes.geometric,arrows.meta,positioning,
%%     fit,backgrounds,calc,decorations.pathreplacing}
%%   \usepackage{pgfplots}
%%   \pgfplotsset{compat=1.18}
%%   \usepackage{listings,xcolor,booktabs,multirow,subcaption}
%%=============================================================

\definecolor{cppblue}{RGB}{30,100,200}
\definecolor{pygreen}{RGB}{40,150,80}
\definecolor{warnred}{RGB}{200,60,60}
\definecolor{orng}{RGB}{220,130,30}

\lstset{
  basicstyle=\ttfamily\small,
  keywordstyle=\color{cppblue},
  commentstyle=\color{gray!70},
  stringstyle=\color{orng},
  breaklines=true,
  frame=single,
  rulecolor=\color{gray!40},
  framexleftmargin=4pt,
  xleftmargin=8pt,
  numbers=none,
}

% ================================================================
\section{AI \#3 --- Tự động hoá khâu khởi tạo trong hệ thống pick \& place}
\label{sec:ai3}
% ================================================================

% ----------------------------------------------------------------
\subsection{Bối cảnh: vấn đề của hệ thống visual servoing ban đầu}
\label{subsec:ai3-context}
% ----------------------------------------------------------------

Hệ thống pick \& place chạy trên cánh tay robot Denso VS-6577E 6 bậc tự do,
sử dụng thư viện ViSP (Visual Servoing Platform) để điều khiển bằng thị giác.
Camera USB được gắn trực tiếp lên đầu công tác theo cấu hình \textit{eye-in-hand} ---
camera di chuyển cùng robot và nhìn thẳng về phía vật thể cần gắp trong suốt
quá trình tiếp cận.

Vấn đề cốt lõi của hệ thống ban đầu nằm ở bước khởi tạo: sau khi robot về
home pose cố định, \textbf{người vận hành phải click chuột vào vị trí cylinder}
trên cửa sổ hiển thị camera thì blob tracker (\texttt{vpDot2}) mới bắt đầu
hoạt động được. Visual servo chỉ chạy sau thao tác thủ công đó. Ngoài ra,
nếu tracker mất target trong lúc servo --- điều hoàn toàn có thể xảy ra khi
cylinder bị che khuất hoặc di chuyển ra ngoài vùng sáng --- hệ thống reset
hoàn toàn và lại đòi hỏi click thủ công ở lần tiếp theo. Mỗi chu kỳ
pick \& place đều cần sự can thiệp của người vận hành, loại bỏ khả năng
vận hành tự động không giám sát.

\begin{figure}[htbp]
  \centering
  \begin{tikzpicture}[
    node distance=0.8cm and 0.6cm,
    bx/.style={draw, rounded corners=4pt, text width=3.0cm, align=center,
               minimum height=0.9cm, font=\small},
    robot/.style={bx, draw=cppblue!70, fill=cppblue!8},
    manual/.style={bx, draw=warnred!70, fill=warnred!8},
    aibox/.style={bx, draw=pygreen!70, fill=pygreen!8},
    arr/.style={->, >=Stealth, thick},
    rarr/.style={arr, draw=warnred!70},
    garr/.style={arr, draw=pygreen!70},
  ]

  %% --- CỘT TRÁI: hệ thống gốc ---
  \node[robot]                    (L1) {PREINIT\\{\scriptsize robot $\to$ home pose}};
  \node[manual, below=of L1]      (L2) {Click thủ công\\{\scriptsize chọn cylinder}};
  \node[robot,  below=of L2]      (L3) {JOINT\\{\scriptsize vpDot2 tracking}};
  \node[manual, below=of L3]      (L4) {Mất tracking\\{\scriptsize $\to$ reset toàn bộ}};

  \draw[arr]  (L1) -- (L2);
  \draw[arr]  (L2) -- (L3);
  \draw[rarr] (L3) -- (L4);
  \draw[rarr, dashed] (L4.west)
        to[out=190,in=170,looseness=1.3] (L1.west)
        node[midway, left, font=\scriptsize, text=warnred!70] {reset};

  \node[font=\bfseries\small, text=warnred, above=0.5cm of L1]
       {\textit{Trước cải tiến}};

  %% --- CỘT PHẢI: hệ thống mới ---
  \node[robot,  right=3.8cm of L1] (R1) {PREINIT\\{\scriptsize robot $\to$ home pose}};
  \node[aibox, below=of R1]        (R2) {AI detect\\{\scriptsize EfficientDet-Lite0}};
  \node[robot,  below=of R2]       (R3) {JOINT\\{\scriptsize AI tracking mỗi frame}};
  \node[aibox, below=of R3]        (R4) {AI re-detect\\{\scriptsize tự phục hồi}};

  \draw[arr]  (R1) -- (R2);
  \draw[garr] (R2) -- (R3);
  \draw[garr] (R3) -- (R4);
  \draw[garr, dashed] (R4.east)
        to[out=0,in=0,looseness=1.5] (R3.east)
        node[midway, right, font=\scriptsize, text=pygreen!70] {tiếp tục};

  %% fallback arrow từ AI detect
  \node[manual, right=1.2cm of R3] (FB)
       {{\scriptsize fallback:}\\click thủ công};
  \draw[rarr, dashed] (R2.east) -| (FB.north)
        node[near start, above, font=\scriptsize, text=warnred] {AI fail};

  \node[font=\bfseries\small, text=pygreen!80!black, above=0.5cm of R1]
       {\textit{Sau cải tiến}};

  %% Đường phân cách
  \draw[thick, dashed, gray!50]
      ($(L1.east)!0.5!(R1.west)+(0,0.8)$) -- ++(0,-6.5);

  \end{tikzpicture}
  \caption{So sánh luồng vận hành trước và sau khi tích hợp module AI.
           Bước yêu cầu thao tác thủ công (đỏ) được thay thế bằng AI tự động (xanh lá).
           Fallback về click thủ công vẫn được giữ lại cho trường hợp AI thất bại.}
  \label{fig:ai3-before-after}
\end{figure}

Hình~\ref{fig:ai3-before-after} minh hoạ sự thay đổi này. Module AI đảm nhiệm
toàn bộ việc phát hiện cylinder ở bước khởi tạo và tự phục hồi khi mất target
trong quá trình servo. Hệ thống vẫn giữ nguyên fallback về click thủ công để
đảm bảo tương thích ngược.


% ----------------------------------------------------------------
\subsection{Xây dựng và huấn luyện mô hình EfficientDet-Lite0}
\label{subsec:ai3-model}
% ----------------------------------------------------------------

\subsubsection*{Lựa chọn kiến trúc}

Yêu cầu đặt ra khá cụ thể: mô hình phải đủ nhỏ để triển khai được trên
Google Coral EdgeTPU, phải hỗ trợ INT8 quantization, và tương thích với
\texttt{tflite\_model\_maker} --- framework đã có sẵn trong môi trường phát triển.
\textbf{EfficientDet-Lite0} đáp ứng tất cả: đây là biến thể nhỏ nhất trong
họ EfficientDet được Google tối ưu cho thiết bị edge, input cố định
$320 \times 320$ pixel, backbone là EfficientNet-B0 đã được simplify.

\begin{table}[htbp]
  \centering
  \caption{Thông số mô hình EfficientDet-Lite0}
  \label{tab:ai3-model}
  \small
  \begin{tabular}{ll}
    \toprule
    \textbf{Thuộc tính}       & \textbf{Giá trị} \\
    \midrule
    Kiến trúc                 & EfficientDet-Lite0 \\
    Framework training        & TensorFlow~2.8.4 / \texttt{tflite\_model\_maker}~0.4.3 \\
    Input size                & $320 \times 320$~px, uint8 \\
    Số lớp đối tượng          & 2 (cylinder~--~lớp~0, cube~--~lớp~1) \\
    Số epochs / batch size    & 50 / 8 \\
    Transfer learning         & Frozen backbone, fine-tune detection head \\
    Target hardware           & Google Coral USB Accelerator (EdgeTPU) \\
    \bottomrule
  \end{tabular}
\end{table}

Mô hình được huấn luyện để phát hiện \textbf{hai lớp}: cylinder và cube nhựa PLA
in 3D. Việc đưa thêm cube vào tập huấn luyện là có chủ ý --- giúp mô hình học
cách phân biệt hai loại vật thể, từ đó giảm false positive khi cả hai cùng
xuất hiện trong cảnh. Ở phía tích hợp với ViSP, chỉ kết quả cylinder
(\texttt{class\_id~==~0}) được dùng; dự đoán cube bị lọc bỏ.

\subsubsection*{Tập dữ liệu}

Toàn bộ ảnh được thu trực tiếp từ camera eye-in-hand gắn trên robot Denso,
đảm bảo điều kiện thực tế sát với lúc triển khai. Ảnh bao gồm nhiều cấu hình:
góc chụp khác nhau, khoảng cách thay đổi, và số lượng vật thể từ 1 đến 5 trong
một cảnh.

\begin{figure}[htbp]
  \centering
  \begin{tikzpicture}
    \begin{axis}[
      xbar, bar width=0.42cm,
      xlabel={Số lượng ảnh (xấp xỉ)},
      ylabel={Danh mục},
      symbolic y coords={
        Background,
        2 cube + 3 cyl,
        2 cylinders,
        3 cylinders,
        1 cyl + 1 cube,
        1 cube,
        1 cylinder,
        Cylinder (gốc)
      },
      ytick=data,
      xmin=0, xmax=175,
      width=9.5cm, height=6.8cm,
      nodes near coords,
      nodes near coords align={horizontal},
      every node near coord/.style={font=\scriptsize},
      tick label style={font=\small},
      xlabel style={font=\small},
      ylabel style={font=\small},
      enlarge y limits=0.08,
    ]
    \addplot[fill=cppblue!55, draw=cppblue!80] coordinates {
      (38,  Background)
      (28,  2 cube + 3 cyl)
      (52,  2 cylinders)
      (43,  3 cylinders)
      (68,  1 cyl + 1 cube)
      (38,  1 cube)
      (108, 1 cylinder)
      (144, Cylinder (gốc))
    };
    \end{axis}
  \end{tikzpicture}
  \caption{Phân bố tập dữ liệu theo danh mục ảnh (~534 ảnh tổng,
           633 annotation cylinder, 170 annotation cube).
           Các cảnh đa vật thể giúp mô hình học phát hiện từng đối tượng
           riêng lẻ trong môi trường phức tạp.}
  \label{fig:ai3-dataset}
\end{figure}

Annotation được thực hiện bằng \texttt{labelme} theo format bounding box hình chữ
nhật. Workflow:
\begin{enumerate}
  \item Vẽ bounding box cho từng cylinder và cube bằng \texttt{labelme}
  \item Chuyển đổi labelme JSON $\to$ COCO JSON thống nhất
        (script \texttt{labelme\_to\_coco.py})
  \item Tập dữ liệu được split theo tỉ lệ $70/15/15$ (train/val/test),
        seed cố định để đảm bảo tính tái lập
\end{enumerate}

Phân chia ba chiều thay vì hai chiều đảm bảo tập test hoàn toàn độc lập ---
không được nhìn thấy trong suốt quá trình training hoặc điều chỉnh siêu tham số.

\subsubsection*{Pipeline huấn luyện}

\begin{figure}[htbp]
  \centering
  \begin{tikzpicture}[
    node distance=0.45cm and 0.15cm,
    st/.style={draw, rounded corners=3pt, fill=gray!10, draw=gray!55,
               text width=2.0cm, align=center, minimum height=0.85cm, font=\small},
    blue_st/.style={st, fill=cppblue!10, draw=cppblue!55},
    out_st/.style={st, fill=orng!20, draw=orng!70},
    arr/.style={->, >=Stealth, thick, gray!60},
    lbl/.style={font=\scriptsize, text=gray!80},
  ]

  \node[blue_st] (s1) {Annotation\\{\scriptsize labelme}};
  \node[blue_st, right=of s1] (s2) {COCO\\JSON};
  \node[blue_st, right=of s2] (s3) {Pascal\\VOC XML};
  \node[blue_st, right=of s3] (s4) {Training\\{\scriptsize 50 epochs}};
  \node[out_st, right=of s4] (s5) {SavedModel\\{\scriptsize Float32}};
  \node[out_st, below=0.55cm of s5] (s6) {INT8\\TFLite};
  \node[out_st, below=0.55cm of s6] (s7) {EdgeTPU\\TFLite};

  \draw[arr] (s1) -- (s2) node[lbl,above,midway] {to\_coco.py};
  \draw[arr] (s2) -- (s3) node[lbl,above,midway] {tflite\_mm};
  \draw[arr] (s3) -- (s4);
  \draw[arr] (s4) -- (s5) node[lbl,above,midway] {export};
  \draw[arr] (s5) -- (s6) node[lbl,right,midway] {PTQ};
  \draw[arr] (s6) -- (s7) node[lbl,right,midway] {edgetpu\_compiler};

  \end{tikzpicture}
  \caption{Pipeline huấn luyện và xuất mô hình. Từ ảnh thô đến ba file
           mô hình phục vụ các môi trường triển khai khác nhau.
           Bước chuyển đổi sang Pascal VOC XML là yêu cầu bắt buộc của
           \texttt{tflite\_model\_maker}.}
  \label{fig:ai3-pipeline}
\end{figure}

Transfer learning với \texttt{train\_whole\_model=False} --- chỉ fine-tune detection
head, giữ nguyên backbone đã pretrain --- phù hợp với tập dữ liệu tương đối nhỏ
($\sim$374 ảnh train) và giảm nguy cơ overfitting đáng kể so với
huấn luyện từ đầu.


% ----------------------------------------------------------------
\subsection{Quantization INT8 và biên dịch cho EdgeTPU}
\label{subsec:ai3-quant}
% ----------------------------------------------------------------

Mô hình sau training là Float32 --- không chạy được trực tiếp trên Coral EdgeTPU.
EdgeTPU chỉ thực thi mô hình INT8 quantized, trong đó trọng số và activation
được biểu diễn bằng số nguyên 8-bit. Lợi ích: kích thước giảm $\sim4\times$,
và phần cứng EdgeTPU đạt tốc độ inference cực nhanh ($\sim$4--15~ms/frame)
nhờ tận dụng phép tính số nguyên.

Post-Training Quantization (PTQ) sử dụng tập train làm \textit{representative
dataset} để calibrate scale và zero\_point cho từng layer --- đây là bước bắt
buộc để duy trì độ chính xác của INT8 model gần với Float32 baseline.

File \texttt{\_edgetpu.tflite} sau khi biên dịch chứa custom op
\texttt{edgetpu-custom-op} mà TFLite runtime thông thường
\textit{không thể thực thi trên CPU}, gây lỗi:

\begin{lstlisting}[language={}]
RuntimeError: Encountered unresolved custom op: edgetpu-custom-op
\end{lstlisting}

Do đó hệ thống \textbf{duy trì hai file riêng biệt} và lựa chọn tự động:

\begin{center}
\small
\begin{tabular}{lll}
  \toprule
  File & Runtime & Deployment \\
  \midrule
  \texttt{cylinder\_detector\_int8\_edgetpu.tflite} & pycoral   & Coral EdgeTPU \\
  \texttt{cylinder\_detector\_int8.tflite}          & TFLite/ai-edge-litert & CPU \\
  \bottomrule
\end{tabular}
\end{center}


% ----------------------------------------------------------------
\subsection{Kiến trúc tích hợp: C++ ViSP ↔ Python subprocess}
\label{subsec:ai3-arch}
% ----------------------------------------------------------------

\subsubsection*{Quyết định thiết kế}

Thay vì link trực tiếp thư viện TFLite C++ vào build system của ViSP
(đòi hỏi thay đổi lớn vào CMakeLists và phụ thuộc version build cụ thể),
hệ thống áp dụng kiến trúc \textbf{subprocess model}: C++ gọi Python script
như một tiến trình con qua \texttt{popen()}, giao tiếp một chiều qua stdout.
Cách này giúp module AI hoàn toàn tách biệt --- cập nhật mô hình không cần
recompile code robot.

\begin{figure}[htbp]
  \centering
  \begin{tikzpicture}[
    node distance=0.38cm and 0.5cm,
    comp/.style={draw, rounded corners=4pt, text width=3.1cm, align=center,
                 minimum height=0.82cm, font=\small},
    cpp/.style={comp, fill=cppblue!10, draw=cppblue!55},
    py/.style={comp,  fill=pygreen!10, draw=pygreen!60},
    tmp/.style={comp,  fill=orng!18, draw=orng!65, text width=2.5cm,
                minimum height=0.7cm, font=\scriptsize},
    arr/.style={->, >=Stealth, thick},
    ]

  %% ---- C++ side ----
  \node[cpp] (cap)   {Camera capture\\{\scriptsize vpImage{$<$}uchar{$>$}}};
  \node[cpp, below=of cap]  (save)  {Save frame\\{\scriptsize cv::imwrite()}};
  \node[cpp, below=of save] (pop)   {popen() subprocess\\{\scriptsize build cmd}};
  \node[cpp, below=of pop]  (sel)   {select() timeout 3s\\{\scriptsize POSIX wait}};
  \node[cpp, below=of sel]  (parse) {Parse stdout\\{\scriptsize sscanf SUCCESS}};
  \node[cpp, below=of parse](feat)  {vpFeatureBuilder\\{\scriptsize create feature}};

  \draw[arr](cap)--(save); \draw[arr](save)--(pop);
  \draw[arr](pop)--(sel);  \draw[arr](sel)--(parse); \draw[arr](parse)--(feat);

  \node[draw=cppblue!50, dashed, rounded corners=6pt,
        fit=(cap)(save)(pop)(sel)(parse)(feat),
        inner sep=7pt] (cbox) {};
  \node[above=2pt of cbox, font=\bfseries\small, text=cppblue]
       {C++ (ViSP main loop)};

  %% ---- Temp file ----
  \node[tmp, right=2.6cm of save] (jpg)
       {/tmp/visp\_ai\_frame.jpg};

  %% ---- Python side ----
  \node[py, right=2.6cm of pop]   (load) {Load model\\{\scriptsize fallback chain}};
  \node[py, below=of load]        (pre)  {Preprocess\\{\scriptsize resize 320$\times$320}};
  \node[py, below=of pre]         (inf)  {Inference\\{\scriptsize invoke()}};
  \node[py, below=of inf]         (post) {Post-process\\{\scriptsize dequant+filter}};
  \node[py, below=of post]        (out)  {Print stdout\\{\scriptsize SUCCESS u v conf}};

  \draw[arr](load)--(pre); \draw[arr](pre)--(inf);
  \draw[arr](inf)--(post); \draw[arr](post)--(out);

  \node[draw=pygreen!55, dashed, rounded corners=6pt,
        fit=(load)(pre)(inf)(post)(out),
        inner sep=7pt] (pbox) {};
  \node[above=2pt of pbox, font=\bfseries\small, text=pygreen!80!black]
       {Python subprocess};

  %% ---- Data flow arrows ----
  \draw[arr, orng!80]  (save.east) -- (jpg.west)  node[midway,above,font=\scriptsize]{write};
  \draw[arr, orng!80]  (jpg.east)  -- (load.west) node[midway,above,font=\scriptsize]{read};
  \draw[arr, cppblue!60](pop.east)  -- (load.west) node[midway,above,font=\scriptsize,yshift=2pt]{spawn};
  \draw[arr, pygreen!70, bend right=25](out.west)  to (parse.east)
        node[midway,below,font=\scriptsize]{stdout pipe};

  \end{tikzpicture}
  \caption{Kiến trúc tích hợp C++ ViSP $\leftrightarrow$ Python subprocess.
           C++ lưu frame hiện tại ra file tạm, spawn Python làm child process
           qua \texttt{popen()}, đọc kết quả từ stdout sau khi qua timeout guard.}
  \label{fig:ai3-arch}
\end{figure}

\subsubsection*{Protocol trao đổi dữ liệu}

Output của Python script là một dòng duy nhất trên stdout:

\begin{lstlisting}[language={}]
SUCCESS 272.86 190.74 0.8362 238.42 166.89 68.88 47.68
%        u      v      conf   x_min  y_min  width  height
\end{lstlisting}

hoặc \texttt{FAILURE~<lý~do>} nếu không phát hiện được. Tất cả debug log
đi ra \texttt{stderr} --- C++ chỉ đọc stdout --- tránh nhiễu kết quả parse.

Đáng chú ý: trong ViSP, \texttt{vpImagePoint} được định nghĩa là \texttt{(row,~col)}
tương ứng \texttt{(v,~u)}, không phải \texttt{(u,~v)} như thông thường:

\begin{lstlisting}[language=C++]
detected_center = vpImagePoint(v, u);  // row=v, col=u
\end{lstlisting}

Tham số \texttt{hint} cho phép Python chọn cylinder \textit{gần nhất với vị
trí đã biết} thay vì cylinder có confidence cao nhất --- cơ chế này tránh
hiện tượng servo ``nhảy'' sang cylinder khác khi scores thay đổi giữa các frame.


% ----------------------------------------------------------------
\subsection{State machine: PREINIT \texorpdfstring{$\to$}{→} INIT (AI) \texorpdfstring{$\to$}{→} JOINT \texorpdfstring{$\to$}{→} ...}
\label{subsec:ai3-sm}
% ----------------------------------------------------------------

Hệ thống vận hành theo vòng lặp chính liên tục với 6 trạng thái.
Mỗi iteration đọc một frame camera và thực thi logic trạng thái hiện tại.

\begin{figure}[htbp]
  \centering
  \begin{tikzpicture}[
    node distance=1.0cm and 2.0cm,
    st/.style={draw, rounded corners=4pt, fill=cppblue!10, draw=cppblue!55,
               text width=2.3cm, align=center, minimum height=1.0cm, font=\small},
    ai/.style={st, fill=pygreen!10, draw=pygreen!60},
    arr/.style={->, >=Stealth, thick},
    lbl/.style={font=\scriptsize, text=gray!70},
  ]

  %% Top row
  \node[st]  (pre)  {PREINIT\\{\scriptsize home pose}};
  \node[ai,  right=of pre]   (ini)  {INIT\\{\scriptsize AI detect}};
  \node[ai,  right=of ini]   (jnt)  {JOINT\\{\scriptsize AI tracking}};
  \node[st,  right=of jnt]   (app)  {APPROACH\\{\scriptsize chờ RC5}};

  %% Bottom row
  \node[st,  below=1.8cm of app]  (grp)  {GRIPPER\\{\scriptsize đóng gripper}};
  \node[st,  left=of grp]         (cls)  {CLASSIFIED\\{\scriptsize vận chuyển}};

  %% Arrows
  \draw[arr](pre)--(ini)  node[lbl,above,midway]{home\\reached};
  \draw[arr](ini)--(jnt)  node[lbl,above,midway]{AI init\\ok};
  \draw[arr](jnt)--(app)  node[lbl,above,midway]{$|e|<5{\times}10^{-3}$};
  \draw[arr](app)--(grp)  node[lbl,right,midway]{RC5\\signal};
  \draw[arr](grp)--(cls)  node[lbl,below,midway]{grasped};
  \draw[arr](cls.west) to[out=180,in=270,looseness=0.8] (pre.south)
        node[lbl,left,pos=0.5]{cycle\\restart};

  %% AI self-recovery
  \draw[arr, pygreen!70, dashed]
        (jnt.north) to[out=100,in=80,looseness=2.5] (jnt.north)
        node[above=0.7cm of jnt, font=\scriptsize, text=pygreen!80!black]
        {AI re-detect};

  %% Fallback from INIT
  \draw[arr, warnred!70, dashed]
        (ini.south) -- ++(0,-0.7)
        node[below, font=\scriptsize, text=warnred]
        {fallback: click thủ công};

  \end{tikzpicture}
  \caption{State machine của hệ thống pick \& place sau cải tiến.
           Trạng thái tô xanh lá là nơi AI được gọi.
           Vòng lặp tự phục hồi ở JOINT (mũi tên đứt nét) cho phép hệ thống
           tiếp tục mà không cần reset khi mất detection.}
  \label{fig:ai3-sm}
\end{figure}

\paragraph{PREINIT.}
Robot gửi lệnh joint position về home pose
$[0^\circ, 0^\circ, 90^\circ, 0^\circ, 90^\circ, 0^\circ]$.
Servo task được cấu hình với interaction matrix Desired, pseudo-inverse,
gain $\lambda = 0.4$.

\paragraph{INIT.}
Chờ robot đạt home pose (sai số từng joint $< 0.01$) và gripper mở xong.
Khi cả hai điều kiện thỏa, \texttt{detectCylinderWithAI()} được gọi
không có hint --- Python chọn cylinder có confidence cao nhất.
Kết quả trực tiếp tạo visual feature \texttt{vpFeaturePoint}
và add vào servo task. Nếu AI thất bại, \texttt{ai\_hint}
giữ giá trị mặc định (fallback thủ công).

\paragraph{JOINT.}
Khác với cách tiếp cận truyền thống dùng blob tracker mỗi frame,
hệ thống hiện tại \textbf{gọi AI trực tiếp ở mỗi iteration}
với \texttt{hint = \&ai\_hint}. Control law:
\begin{equation}
  \mathbf{v} = -\lambda\,\mathbf{L}_{s^*}^{+}\,\mathbf{e},
  \qquad \mathbf{q}_{k+1} = \mathbf{q}_k + \mathbf{v}\,\Delta t
  \label{eq:ibvs}
\end{equation}
với $\Delta t = 0.5\,$s. Điều kiện hội tụ: $|e_x| < 5{\times}10^{-3}$
và $|e_y| < 5{\times}10^{-3}$ (normalized image coordinates).
Khi đạt, gửi \texttt{"OKE\textbackslash r"} qua UART đến RC5 controller.


% ----------------------------------------------------------------
\subsection{Các vấn đề kỹ thuật và giải pháp}
\label{subsec:ai3-issues}
% ----------------------------------------------------------------

Một số vấn đề nảy sinh trong quá trình tích hợp mà khó thấy trước khi thực
sự chạy hệ thống.

\subsubsection*{P1: Xác định output tensor sai lệch}

\texttt{tflite\_model\_maker} 0.4.x đặt tên tensor theo thứ tự riêng của nó:

\begin{center}
\small
\begin{tabular}{ccc}
  \toprule
  Tensor & Shape & Nội dung \\
  \midrule
  \texttt{:0} & \texttt{[1]}     & num\_detections \\
  \texttt{:1} & \texttt{[1,\,N]} & scores \quad \textbf{← cần lấy cái này} \\
  \texttt{:2} & \texttt{[1,\,N]} & class IDs \quad \textbf{← không phải cái này} \\
  \texttt{:3} & \texttt{[1,\,N,\,4]} & boxes \\
  \bottomrule
\end{tabular}
\end{center}

Thứ tự này không khớp với API của pycoral, vốn giả định positional order
\texttt{boxes/classes/scores/count = 0/1/2/3}. Dùng API đó sẽ đọc
num\_detections (raw uint8 = 255) như scores và crash với \texttt{IndexError}.
Giải pháp: xác định tensor \textbf{theo shape động} --- boxes là
\texttt{[1,N,4]}, scores là tensor \texttt{[1,N]} có name suffix nhỏ nhất
(\texttt{:1} chứ không phải \texttt{:2}).

\subsubsection*{P2: Python 3.10 không có \texttt{tflite-runtime}}

Google không publish wheel \texttt{tflite-runtime} cho Python 3.10+.
Giải pháp là xây dựng \textbf{fallback chain 4 tầng}:

\begin{figure}[htbp]
  \centering
  \begin{tikzpicture}[
    node distance=0.4cm,
    fb/.style={draw, rounded corners=3pt, text width=7.5cm, align=left,
               minimum height=0.85cm, font=\small, inner sep=7pt},
    t1/.style={fb, fill=purple!12, draw=purple!55},
    t2/.style={fb, fill=teal!10,   draw=teal!50},
    t3/.style={fb, fill=cppblue!10, draw=cppblue!50},
    t4/.style={fb, fill=gray!12,   draw=gray!50},
    fa/.style={fb, fill=warnred!10, draw=warnred!50},
    arr/.style={->, >=Stealth, thick, dashed, gray!55},
    lbl/.style={font=\scriptsize, text=warnred!80},
  ]
  \node[t1] (a) {\textbf{1.}~\texttt{pycoral.utils.edgetpu}\quad{\scriptsize Coral EdgeTPU, Python 3.9}};
  \node[t2, below=of a] (b) {\textbf{2.}~\texttt{ai\_edge\_litert.interpreter}\quad{\scriptsize Google's new pkg, Python 3.8--3.12}};
  \node[t3, below=of b] (c) {\textbf{3.}~\texttt{tflite\_runtime.interpreter}\quad{\scriptsize Legacy, Python 3.9}};
  \node[t4, below=of c] (d) {\textbf{4.}~\texttt{tensorflow.lite.Interpreter}\quad{\scriptsize Heavy fallback}};
  \node[fa, below=of d] (e) {\textbf{FAILURE}:~\texttt{no\_tflite\_backend\_available}};

  \draw[arr](a)--(b) node[lbl,right,midway]{ImportError / Coral not connected};
  \draw[arr](b)--(c) node[lbl,right,midway]{ImportError};
  \draw[arr](c)--(d) node[lbl,right,midway]{ImportError};
  \draw[arr](d)--(e) node[lbl,right,midway]{ImportError};
  \end{tikzpicture}
  \caption{Python backend fallback chain trong \texttt{detect\_cylinder.py}.
           Script tự động thử từng backend theo thứ tự ưu tiên,
           tương thích với cả môi trường có Coral lẫn CPU-only.}
  \label{fig:ai3-fallback}
\end{figure}

\texttt{ai-edge-litert} là gói chính thức của Google thay thế \texttt{tflite-runtime},
hỗ trợ Python 3.8--3.12 với API giống hệt nhau --- là giải pháp cho vấn đề Python 3.10.

\subsubsection*{P3: JSON comment trong config}

\texttt{config.json} dùng comment \texttt{//} để ghi chú, nhưng cả Python
\texttt{json.load()} và \texttt{nlohmann::json} đều từ chối mặc định.
Phía C++: \texttt{nlohmann::json::parse(file, nullptr, true, true)} với
\texttt{ignore\_comments=true} (có từ nlohmann/json~3.9.0).
Phía Python: strip comment lines thủ công trước khi \texttt{json.loads()}.


% ----------------------------------------------------------------
\subsection{Thông số cấu hình visual servo}
\label{subsec:ai3-servo}
% ----------------------------------------------------------------

Hệ thống sử dụng \textbf{Image-Based Visual Servoing (IBVS)} với feature là
tọa độ 2D của cylinder trong không gian ảnh normalized $\mathbf{s}=(x,y)$,
desired position $\mathbf{s}^*=(0,0,1)$ --- cylinder phải nằm đúng tâm ảnh.

\begin{table}[htbp]
  \centering
  \caption{Thông số điều khiển visual servo}
  \label{tab:ai3-servo}
  \small
  \begin{tabular}{ll}
    \toprule
    \textbf{Tham số}            & \textbf{Giá trị} \\
    \midrule
    Loại điều khiển             & Eye-in-hand, \texttt{EYEINHAND\_L\_cVe\_eJe} \\
    Interaction matrix          & Desired ($\mathbf{L}_{s^*}$), pseudo-inverse \\
    Gain $\lambda$              & 0.4 \\
    Bước thời gian $\Delta t$   & 0.5\,s \\
    Ngưỡng hội tụ               & $|e_x| < 5{\times}10^{-3}$ và $|e_y| < 5{\times}10^{-3}$ \\
    Desired position            & $(x^*, y^*) = (0,\,0)$ --- tâm ảnh, $Z = 1$ \\
    \bottomrule
  \end{tabular}
\end{table}

Dùng $\mathbf{L}_{s^*}$ (tại vị trí mong muốn) thay vì $\mathbf{L}_{s_{cur}}$
giúp tránh singularity khi feature tiến gần desired point và giữ ổn định
theo ổn định hơn trong giai đoạn hội tụ. Velocity được saturate trước khi
gửi xuống robot để đảm bảo không vượt giới hạn vật lý từng khớp.
   %% mục 4.4
%%     ...
%%     \chapter{Kết quả và Đánh giá}
%%     ...
%%     %%=============================================================
%% chapter5_4.tex
%% Phần 5.4 — Kết quả đánh giá: AI #3 Phát hiện vật thể
%%
%% \input file — cần thêm vào preamble:
%%   \usepackage{tikz}
%%   \usetikzlibrary{arrows.meta,positioning,patterns}
%%   \usepackage{pgfplots}
%%   \pgfplotsset{compat=1.18}
%%   \usepackage{booktabs,multirow,subcaption,xcolor}
%%=============================================================

% ================================================================
\section{Kết quả và Đánh giá: AI \#3 --- Phát hiện vật thể}
\label{sec:ai3-results}
% ================================================================

% ----------------------------------------------------------------
\subsection{Kết quả detection model: COCO metrics}
\label{subsec:ai3-coco}
% ----------------------------------------------------------------

Mô hình được đánh giá trên cả validation set và test set độc lập
theo chuẩn \textbf{COCO evaluation protocol} --- tính trung bình mAP qua
các ngưỡng IoU từ 0.50 đến 0.95 với bước 0.05.

\begin{figure}[htbp]
  \centering
  \begin{tikzpicture}
    \begin{axis}[
      ybar, bar width=0.45cm,
      width=10cm, height=7cm,
      symbolic x coords={mAP (0.50:0.95), AP50, AP75, AR},
      xtick=data,
      x tick label style={font=\small, align=center},
      ymin=0, ymax=1.15,
      ylabel={Score},
      ylabel style={font=\small},
      legend style={at={(0.5,1.03)}, anchor=south, legend columns=2, font=\small},
      legend cell align=left,
      nodes near coords,
      every node near coord/.style={font=\scriptsize, /pgf/number format/fixed,
                                    /pgf/number format/precision=3},
      enlarge x limits=0.25,
      grid=major, grid style={dashed, gray!30},
      axis line style={gray!60},
      tick style={gray!60},
    ]
    \addplot[fill=cppblue!70, draw=cppblue!90] coordinates {
      (mAP (0.50:0.95), 0.791)
      (AP50, 0.990)
      (AP75, 0.990)
      (AR,   0.869)
    };
    \addplot[fill=orng!80, draw=orng!90] coordinates {
      (mAP (0.50:0.95), 0.799)
      (AP50, 1.000)
      (AP75, 0.969)
      (AR,   0.856)
    };
    \legend{Validation, Test}
    \end{axis}
  \end{tikzpicture}
  \caption{COCO metrics trên validation set và test set. AP50 trên test set
           đạt 1.000 cho thấy 100\% cylinder được phát hiện đúng ở ngưỡng IoU
           lỏng lẻ (50\%). mAP$\approx$0.80 phản ánh vẫn còn đây đó những trường hợp
           định vị bounding box chưa được chính xác ở ngưỡng IoU cao.}
  \label{fig:ai3-coco}
\end{figure}

Bảng~\ref{tab:ai3-metrics} tổng hợp số liệu chi tiết:

\begin{table}[htbp]
  \centering
  \caption{COCO metrics trên validation và test set}
  \label{tab:ai3-metrics}
  \small
  \begin{tabular}{lcc}
    \toprule
    \textbf{Metric}                    & \textbf{Validation} & \textbf{Test} \\
    \midrule
    mAP (IoU 0.50:0.95)               & 0.791               & \textbf{0.799} \\
    AP50 (IoU $\geq$ 0.50)            & 0.990               & \textbf{1.000} \\
    AP75 (IoU $\geq$ 0.75)            & 0.990               & 0.969 \\
    AR (maxDets=100)                   & 0.869               & 0.856 \\
    \midrule
    Precision (IoU=0.50, conf=0.50)    & ---                 & \textbf{1.000} \\
    Recall    (IoU=0.50, conf=0.50)    & ---                 & 0.990 \\
    F1-Score                           & ---                 & \textbf{0.995} \\
    TP / FP / FN                       & ---                 & 95 / 0 / 1 \\
    \bottomrule
  \end{tabular}
\end{table}

\begin{figure}[htbp]
  \centering
  \begin{tikzpicture}
    \begin{axis}[
      ybar, bar width=0.55cm,
      width=8.5cm, height=6cm,
      symbolic x coords={Precision, Recall, F1-Score},
      xtick=data,
      x tick label style={font=\small},
      ymin=0.96, ymax=1.04,
      ylabel={Score},
      ylabel style={font=\small},
      nodes near coords,
      every node near coord/.style={font=\small, /pgf/number format/fixed,
                                    /pgf/number format/precision=3},
      enlarge x limits=0.3,
      grid=major, grid style={dashed, gray!30},
      axis line style={gray!60},
    ]
    \addplot[fill=pygreen!65, draw=pygreen!80] coordinates {
      (Precision, 1.000)
      (Recall,    0.990)
      (F1-Score,  0.995)
    };
    \end{axis}
  \end{tikzpicture}
  \caption{Precision, Recall và F1-Score trên test set (PASCAL VOC protocol,
           IoU$\geq$0.50, confidence$\geq$0.50, chỉ lớp cylinder).
           Precision = 1.000 và FP = 0 chứng tỏ không có false positive
           nào ở ngưỡng 0.50. FN = 1 (một cylinder bị bỏ sót) giải thích
           Recall = 0.990.}
  \label{fig:ai3-prf}
\end{figure}

Giá trị AP50 = 1.000 trên test set à ngưỡng IoU lỏng lẻ (50\%) phản ánh
mô hình phát hiện được toàn bộ cylinder. mAP = 0.799 (trung bình
qua IoU 0.50--0.95) thấp hơn một chút do mô hình EfficientDet-Lite0
đơn giản vẫn còn giới hạn trong việc định vị bounding box chính xác
ở các ngưỡng IoU cao ($\geq$0.75) --- điều phổ biến với các kiến trúc
nhe hơn.


% ----------------------------------------------------------------
\subsection{Phân tích confidence score: true positive vs false positive}
\label{subsec:ai3-conf}
% ----------------------------------------------------------------

Một trong những điều quan trọng khi triển khai thực tế là khoảng cách
phân tách giữa true positive và false positive --- điều này quyết định
ngưỡng confidence nào thực tế dùng được mà không sai kết quả.

Kiểm tra trên 218 ảnh cylinder (true positives) và 10 ảnh cube
(false positives):

\begin{figure}[htbp]
  \centering
  \begin{tikzpicture}
    \begin{axis}[
      width=11cm, height=5.5cm,
      xlabel={Confidence score},
      ylabel={},
      xmin=0, xmax=1.05,
      ymin=0, ymax=3,
      xtick={0,0.1,...,1.0},
      ytick={0.7, 2.0},
      yticklabels={\small False Positive\\(cube, $n$=10), \small True Positive\\(cylinder, $n$=218)},
      y tick label style={align=right, font=\small},
      x tick label style={font=\small},
      xlabel style={font=\small},
      grid=both, grid style={dashed, gray!25},
      axis line style={gray!60},
    ]

    %% --- True Positive band (y=2) ---
    %% Most scores: 0.85-0.98 (majority)
    \addplot[only marks, mark=|, mark size=5pt, mark options={draw=pygreen!70, line width=1.2pt},
             opacity=0.5]
    coordinates {
      (0.29,2)(0.61,2)(0.74,2)
      (0.82,2)(0.83,2)(0.84,2)(0.85,2)(0.86,2)(0.87,2)(0.88,2)
      (0.89,2)(0.90,2)(0.91,2)(0.92,2)(0.93,2)(0.94,2)(0.94,2)
      (0.95,2)(0.95,2)(0.96,2)(0.96,2)(0.97,2)(0.97,2)(0.98,2)
    };

    %% --- False Positive band (y=0.7) ---
    \addplot[only marks, mark=|, mark size=5pt, mark options={draw=warnred!70, line width=1.2pt},
             opacity=0.6]
    coordinates {
      (0.08,0.7)(0.13,0.7)(0.17,0.7)(0.21,0.7)
      (0.25,0.7)(0.28,0.7)(0.30,0.7)(0.33,0.7)(0.35,0.7)(0.36,0.7)
    };

    %% --- Threshold line at 0.50 ---
    \addplot[thick, dashed, draw=gray!80, domain=0:3]
    coordinates {(0.50, 0) (0.50, 3)};
    \node[font=\scriptsize, text=gray!80, rotate=90]
         at (axis cs: 0.52, 1.5) {threshold = 0.50};

    %% --- Gap annotation ---
    \draw[->, >=Stealth, thick, gray!70]
         (axis cs:0.38, 2.5) -- (axis cs:0.48, 2.5);
    \draw[->, >=Stealth, thick, gray!70]
         (axis cs:0.34, 2.5) -- (axis cs:0.27, 2.5);
    \node[font=\scriptsize, text=gray!80] at (axis cs:0.37, 2.7) {gap = 0.34};

    %% --- Median annotation for TP ---
    \addplot[only marks, mark=diamond*, mark size=4pt,
             mark options={draw=pygreen!80!black, fill=pygreen!60}]
    coordinates {(0.94, 2)};
    \node[font=\scriptsize, text=pygreen!80!black, above]
         at (axis cs:0.94, 2) {median 0.94};

    %% --- Max FP annotation ---
    \addplot[only marks, mark=diamond*, mark size=4pt,
             mark options={draw=warnred!80, fill=warnred!50}]
    coordinates {(0.36, 0.7)};
    \node[font=\scriptsize, text=warnred, below]
         at (axis cs:0.36, 0.7) {max 0.36};

    \end{axis}
  \end{tikzpicture}
  \caption{Phân tách confidence score giữa true positive (cylinder, xanh lá)
           và false positive (cube, đỏ). Mỗi vạch đứng đại diện một detection.
           Ngưỡng 0.50 nằm chính giữa khoảng gap 0.34,
           loại bỏ toàn bộ false positive mà vẫn chấp nhận 99.5\% true positive.}
  \label{fig:ai3-conf}
\end{figure}

Khoảng cách 0.34 giữa điểm false positive cao nhất (0.36) và 99.1\%
true positive ($\geq$0.70) xác nhận mô hình đã học được đặc trưng
phân biệt mạnh giữa cylinder và cube. Ngưỡng hiện tại 0.50 nằm an toàn
giữa hai vùng phân bố này.

Ba ảnh outlier có score thấp (0.29, 0.61, 0.74) tương ứng với cylinder
bị nhìn từ góc bất thường hoặc bị che khuất một phần. Ỡ runtime,
các trường hợp này fallback về click thủ công --- hành vi an toàn.

\begin{table}[htbp]
  \centering
  \caption{Phân tích confidence score trên 218 ảnh cylinder + 10 ảnh cube}
  \label{tab:ai3-conf}
  \small
  \begin{tabular}{lcc}
    \toprule
    & \textbf{True Positive (cylinder)} & \textbf{False Positive (cube)} \\
    \midrule
    Score range      & 0.29 -- 0.98          & 0.08 -- 0.36 \\
    Median           & \textbf{0.94}         & --- \\
    $\geq$ 0.50      & 99.5\% (217/218)      & \textbf{0/10 (0\%)} \\
    $\geq$ 0.70      & 99.1\% (216/218)      & 0/10 (0\%) \\
    $\geq$ 0.80      & 95.9\% (209/218)      & 0/10 (0\%) \\
    \midrule
    Gap ph\^an t\'ach & \multicolumn{2}{c}{$0.70 - 0.36 = 0.34$} \\
    \bottomrule
  \end{tabular}
\end{table}


% ----------------------------------------------------------------
\subsection{Kết quả tích hợp hệ thống}
\label{subsec:ai3-integration}
% ----------------------------------------------------------------

\begin{table}[htbp]
  \centering
  \caption{So sánh khả năng trước và sau khi tích hợp AI}
  \label{tab:ai3-comparison}
  \small
  \begin{tabular}{lll}
    \toprule
    \textbf{Tiêu chí}                    & \textbf{Hệ thống gốc} & \textbf{Sau cải tiến} \\
    \midrule
    Khởi tạo mỗi chu kỳ               & Click thủ công (bắt buộc) & Tự động (AI) \\
    Tỷ lệ init thành công             & 100\% (do manual)  & $\sim$99.5\% (AI) \\
    Fallback khi AI thất bại           & Không có              & Click thủ công \\
    Phục hồi khi mất tracking          & Reset toàn bộ        & AI re-detect tự động \\
    Chạy không giám sát                & Không thể              & Có (trường hợp bình thường) \\
    \bottomrule
  \end{tabular}
\end{table}

Thiết kế hiện tại gọi AI trực tiếp \textbf{mỗi iteration} của servo
loop thay vì dùng blob tracker. Điều này có một hàm ý đáng chú ý:
\textit{không có tracking failure} theo nghĩa truyền thống --- mỗi frame
là một lần detection độc lập. Tính ổn định của servo phụ thuộc vào
sự nhất quán của AI detection qua các frame liên tiếp. Cơ chế \texttt{--hint}
(truyền vị trí frame trước) giúp được phần nào khi có nhiều cylinder
trong cảnh.


% ----------------------------------------------------------------
\subsection{Phân tích latency: subprocess CPU vs EdgeTPU}
\label{subsec:ai3-latency}
% ----------------------------------------------------------------

Mỗi lần gọi \texttt{detectCylinderWithAI()} bao gồm nhiều thành phần
latency khác nhau. Điều ít người để ý là: phần overhead lớn nhất
không phải là inference mà là \textbf{Python startup time} --- chi phí
mở một process mới mỗi lần gọi \texttt{popen()}.

\begin{figure}[htbp]
  \centering
  \begin{tikzpicture}
    \begin{axis}[
      xbar stacked,
      bar width=0.55cm,
      width=11cm, height=6.5cm,
      symbolic y coords={
        {EdgeTPU\\(persistent, tương lai)},
        {CPU\\(persistent, tương lai)},
        {EdgeTPU\\(subprocess hiện tại)},
        {CPU\\(subprocess hiện tại)}
      },
      ytick=data,
      y tick label style={align=right, font=\small},
      xlabel={Thời gian (ms)},
      xlabel style={font=\small},
      xmin=0, xmax=780,
      legend style={at={(0.5,-0.18)}, anchor=north, legend columns=4,
                    font=\scriptsize},
      legend cell align=left,
      grid=major, grid style={dashed, gray!25},
      axis line style={gray!60},
      enlarge y limits=0.18,
      nodes near coords, nodes near coords align={horizontal},
      point meta=explicit symbolic,
      every node near coord/.style={font=\scriptsize},
    ]
    %%  stacks: I/O | Inference | Model load | Python startup
    %% CPU subprocess: 10 + 150 + 100 + 350 = 610
    \addplot[fill=gray!40, draw=gray!60] coordinates {
      (10, {CPU\\(subprocess hiện tại)})
      (10, {EdgeTPU\\(subprocess hiện tại)})
      (10, {CPU\\(persistent, tương lai)})
      (5,  {EdgeTPU\\(persistent, tương lai)})
    };  %% I/O
    \addplot[fill=warnred!55, draw=warnred!70] coordinates {
      (150, {CPU\\(subprocess hiện tại)})
      (10,  {EdgeTPU\\(subprocess hiện tại)})
      (150, {CPU\\(persistent, tương lai)})
      (10,  {EdgeTPU\\(persistent, tương lai)})
    };  %% Inference
    \addplot[fill=orng!65, draw=orng!80] coordinates {
      (100, {CPU\\(subprocess hiện tại)})
      (50,  {EdgeTPU\\(subprocess hiện tại)})
      (0,   {CPU\\(persistent, tương lai)})
      (0,   {EdgeTPU\\(persistent, tương lai)})
    };  %% Model load
    \addplot[fill=cppblue!55, draw=cppblue!70] coordinates {
      (350, {CPU\\(subprocess hiện tại)})
      (350, {EdgeTPU\\(subprocess hiện tại)})
      (0,   {CPU\\(persistent, tương lai)})
      (0,   {EdgeTPU\\(persistent, tương lai)})
    };  %% Python startup
    \legend{I/O (JPEG), Inference, Model load, Python startup}
    \end{axis}
  \end{tikzpicture}
  \caption{Phân tích latency mỗi lần gọi AI detection.
           Python startup time ($\sim$350~ms) chiếm phần lẫn tổng overhead
           trong kiến trúc subprocess hiện tại --- lớn hơn nhiều so với
           thời gian inference thực tế. Kiến trúc persistent process (tương lai)
           loại bỏ hoàn toàn chi phí này.}
  \label{fig:ai3-latency}
\end{figure}

\begin{table}[htbp]
  \centering
  \caption{So sánh tổng latency giữa các cấu hình triển khai}
  \label{tab:ai3-latency}
  \small
  \begin{tabular}{lcc}
    \toprule
    \textbf{Cấu hình}                    & \textbf{Inference} & \textbf{Tổng / lần gọi} \\
    \midrule
    CPU subprocess (hiện tại)            & $\sim$100--200~ms  & $\sim$600--700~ms \\
    EdgeTPU subprocess (hiện tại)       & $\sim$4--15~ms     & $\sim$420--430~ms \\
    CPU persistent (tương lai)           & $\sim$100--200~ms  & $\sim$160--210~ms \\
    EdgeTPU persistent (tương lai)      & $\sim$4--15~ms     & $\sim$15--20~ms  \\
    \bottomrule
  \end{tabular}
\end{table}

Kết quả trên cho thấy rõ điểm nghẻn hiện tại:
EdgeTPU inference chỉ mất 4--15~ms nhưng toàn bộ lá gọi
vẫn tốn 420~ms do phải khởi động Python mỗi lần. Với
kiến trúc \textbf{persistent Python process} --- một process
duy nhất khởi động một lần, load mô hình vào bộ nhớ, sau
đó tiếp nhận request từ C++ qua stdin/stdout pipe --- latency
mỗi lần gọi giảm xuống còn $\sim$15--20~ms với EdgeTPU,
mở ra khả năng chạy AI detection \textit{mỗi frame} camera
thay vì chỉ khi cần.

% Chú thích về nguồn số liệu
\smallskip
\noindent\textit{\small
Lưu ý: COCO metrics trong Bảng~\ref{tab:ai3-metrics} và Hình~\ref{fig:ai3-coco}
lấy từ kết quả đánh giá hai lớp (cylinder + cube) trên repository huấn luyện.
Phân tích confidence score trong Bảng~\ref{tab:ai3-conf} lấy từ kiểm tra
độc lập trên 218 ảnh cylinder và 10 ảnh cube (CPU inference via ai-edge-litert).
}
   %% mục 5.4
%%     ...
%%   \end{document}
%%
%% Compile: pdflatex -> pdflatex (chạy 2 lần để có label/ref đúng)
%%=============================================================

%%   \begin{document}
%%     ...
%%     \chapter{Thiết kế và Hiện thực}
%%     ...
%%     %%=============================================================
%% chapter4_4.tex
%% Phần 4.4 — AI #3: Phát hiện vật thể cho pick & place
%%
%% Đây là \input file — thêm vào preamble của main.tex:
%%   \usepackage{tikz}
%%   \usetikzlibrary{shapes.geometric,arrows.meta,positioning,
%%     fit,backgrounds,calc,decorations.pathreplacing}
%%   \usepackage{pgfplots}
%%   \pgfplotsset{compat=1.18}
%%   \usepackage{listings,xcolor,booktabs,multirow,subcaption}
%%=============================================================

\definecolor{cppblue}{RGB}{30,100,200}
\definecolor{pygreen}{RGB}{40,150,80}
\definecolor{warnred}{RGB}{200,60,60}
\definecolor{orng}{RGB}{220,130,30}

\lstset{
  basicstyle=\ttfamily\small,
  keywordstyle=\color{cppblue},
  commentstyle=\color{gray!70},
  stringstyle=\color{orng},
  breaklines=true,
  frame=single,
  rulecolor=\color{gray!40},
  framexleftmargin=4pt,
  xleftmargin=8pt,
  numbers=none,
}

% ================================================================
\section{AI \#3 --- Tự động hoá khâu khởi tạo trong hệ thống pick \& place}
\label{sec:ai3}
% ================================================================

% ----------------------------------------------------------------
\subsection{Bối cảnh: vấn đề của hệ thống visual servoing ban đầu}
\label{subsec:ai3-context}
% ----------------------------------------------------------------

Hệ thống pick \& place chạy trên cánh tay robot Denso VS-6577E 6 bậc tự do,
sử dụng thư viện ViSP (Visual Servoing Platform) để điều khiển bằng thị giác.
Camera USB được gắn trực tiếp lên đầu công tác theo cấu hình \textit{eye-in-hand} ---
camera di chuyển cùng robot và nhìn thẳng về phía vật thể cần gắp trong suốt
quá trình tiếp cận.

Vấn đề cốt lõi của hệ thống ban đầu nằm ở bước khởi tạo: sau khi robot về
home pose cố định, \textbf{người vận hành phải click chuột vào vị trí cylinder}
trên cửa sổ hiển thị camera thì blob tracker (\texttt{vpDot2}) mới bắt đầu
hoạt động được. Visual servo chỉ chạy sau thao tác thủ công đó. Ngoài ra,
nếu tracker mất target trong lúc servo --- điều hoàn toàn có thể xảy ra khi
cylinder bị che khuất hoặc di chuyển ra ngoài vùng sáng --- hệ thống reset
hoàn toàn và lại đòi hỏi click thủ công ở lần tiếp theo. Mỗi chu kỳ
pick \& place đều cần sự can thiệp của người vận hành, loại bỏ khả năng
vận hành tự động không giám sát.

\begin{figure}[htbp]
  \centering
  \begin{tikzpicture}[
    node distance=0.8cm and 0.6cm,
    bx/.style={draw, rounded corners=4pt, text width=3.0cm, align=center,
               minimum height=0.9cm, font=\small},
    robot/.style={bx, draw=cppblue!70, fill=cppblue!8},
    manual/.style={bx, draw=warnred!70, fill=warnred!8},
    aibox/.style={bx, draw=pygreen!70, fill=pygreen!8},
    arr/.style={->, >=Stealth, thick},
    rarr/.style={arr, draw=warnred!70},
    garr/.style={arr, draw=pygreen!70},
  ]

  %% --- CỘT TRÁI: hệ thống gốc ---
  \node[robot]                    (L1) {PREINIT\\{\scriptsize robot $\to$ home pose}};
  \node[manual, below=of L1]      (L2) {Click thủ công\\{\scriptsize chọn cylinder}};
  \node[robot,  below=of L2]      (L3) {JOINT\\{\scriptsize vpDot2 tracking}};
  \node[manual, below=of L3]      (L4) {Mất tracking\\{\scriptsize $\to$ reset toàn bộ}};

  \draw[arr]  (L1) -- (L2);
  \draw[arr]  (L2) -- (L3);
  \draw[rarr] (L3) -- (L4);
  \draw[rarr, dashed] (L4.west)
        to[out=190,in=170,looseness=1.3] (L1.west)
        node[midway, left, font=\scriptsize, text=warnred!70] {reset};

  \node[font=\bfseries\small, text=warnred, above=0.5cm of L1]
       {\textit{Trước cải tiến}};

  %% --- CỘT PHẢI: hệ thống mới ---
  \node[robot,  right=3.8cm of L1] (R1) {PREINIT\\{\scriptsize robot $\to$ home pose}};
  \node[aibox, below=of R1]        (R2) {AI detect\\{\scriptsize EfficientDet-Lite0}};
  \node[robot,  below=of R2]       (R3) {JOINT\\{\scriptsize AI tracking mỗi frame}};
  \node[aibox, below=of R3]        (R4) {AI re-detect\\{\scriptsize tự phục hồi}};

  \draw[arr]  (R1) -- (R2);
  \draw[garr] (R2) -- (R3);
  \draw[garr] (R3) -- (R4);
  \draw[garr, dashed] (R4.east)
        to[out=0,in=0,looseness=1.5] (R3.east)
        node[midway, right, font=\scriptsize, text=pygreen!70] {tiếp tục};

  %% fallback arrow từ AI detect
  \node[manual, right=1.2cm of R3] (FB)
       {{\scriptsize fallback:}\\click thủ công};
  \draw[rarr, dashed] (R2.east) -| (FB.north)
        node[near start, above, font=\scriptsize, text=warnred] {AI fail};

  \node[font=\bfseries\small, text=pygreen!80!black, above=0.5cm of R1]
       {\textit{Sau cải tiến}};

  %% Đường phân cách
  \draw[thick, dashed, gray!50]
      ($(L1.east)!0.5!(R1.west)+(0,0.8)$) -- ++(0,-6.5);

  \end{tikzpicture}
  \caption{So sánh luồng vận hành trước và sau khi tích hợp module AI.
           Bước yêu cầu thao tác thủ công (đỏ) được thay thế bằng AI tự động (xanh lá).
           Fallback về click thủ công vẫn được giữ lại cho trường hợp AI thất bại.}
  \label{fig:ai3-before-after}
\end{figure}

Hình~\ref{fig:ai3-before-after} minh hoạ sự thay đổi này. Module AI đảm nhiệm
toàn bộ việc phát hiện cylinder ở bước khởi tạo và tự phục hồi khi mất target
trong quá trình servo. Hệ thống vẫn giữ nguyên fallback về click thủ công để
đảm bảo tương thích ngược.


% ----------------------------------------------------------------
\subsection{Xây dựng và huấn luyện mô hình EfficientDet-Lite0}
\label{subsec:ai3-model}
% ----------------------------------------------------------------

\subsubsection*{Lựa chọn kiến trúc}

Yêu cầu đặt ra khá cụ thể: mô hình phải đủ nhỏ để triển khai được trên
Google Coral EdgeTPU, phải hỗ trợ INT8 quantization, và tương thích với
\texttt{tflite\_model\_maker} --- framework đã có sẵn trong môi trường phát triển.
\textbf{EfficientDet-Lite0} đáp ứng tất cả: đây là biến thể nhỏ nhất trong
họ EfficientDet được Google tối ưu cho thiết bị edge, input cố định
$320 \times 320$ pixel, backbone là EfficientNet-B0 đã được simplify.

\begin{table}[htbp]
  \centering
  \caption{Thông số mô hình EfficientDet-Lite0}
  \label{tab:ai3-model}
  \small
  \begin{tabular}{ll}
    \toprule
    \textbf{Thuộc tính}       & \textbf{Giá trị} \\
    \midrule
    Kiến trúc                 & EfficientDet-Lite0 \\
    Framework training        & TensorFlow~2.8.4 / \texttt{tflite\_model\_maker}~0.4.3 \\
    Input size                & $320 \times 320$~px, uint8 \\
    Số lớp đối tượng          & 2 (cylinder~--~lớp~0, cube~--~lớp~1) \\
    Số epochs / batch size    & 50 / 8 \\
    Transfer learning         & Frozen backbone, fine-tune detection head \\
    Target hardware           & Google Coral USB Accelerator (EdgeTPU) \\
    \bottomrule
  \end{tabular}
\end{table}

Mô hình được huấn luyện để phát hiện \textbf{hai lớp}: cylinder và cube nhựa PLA
in 3D. Việc đưa thêm cube vào tập huấn luyện là có chủ ý --- giúp mô hình học
cách phân biệt hai loại vật thể, từ đó giảm false positive khi cả hai cùng
xuất hiện trong cảnh. Ở phía tích hợp với ViSP, chỉ kết quả cylinder
(\texttt{class\_id~==~0}) được dùng; dự đoán cube bị lọc bỏ.

\subsubsection*{Tập dữ liệu}

Toàn bộ ảnh được thu trực tiếp từ camera eye-in-hand gắn trên robot Denso,
đảm bảo điều kiện thực tế sát với lúc triển khai. Ảnh bao gồm nhiều cấu hình:
góc chụp khác nhau, khoảng cách thay đổi, và số lượng vật thể từ 1 đến 5 trong
một cảnh.

\begin{figure}[htbp]
  \centering
  \begin{tikzpicture}
    \begin{axis}[
      xbar, bar width=0.42cm,
      xlabel={Số lượng ảnh (xấp xỉ)},
      ylabel={Danh mục},
      symbolic y coords={
        Background,
        2 cube + 3 cyl,
        2 cylinders,
        3 cylinders,
        1 cyl + 1 cube,
        1 cube,
        1 cylinder,
        Cylinder (gốc)
      },
      ytick=data,
      xmin=0, xmax=175,
      width=9.5cm, height=6.8cm,
      nodes near coords,
      nodes near coords align={horizontal},
      every node near coord/.style={font=\scriptsize},
      tick label style={font=\small},
      xlabel style={font=\small},
      ylabel style={font=\small},
      enlarge y limits=0.08,
    ]
    \addplot[fill=cppblue!55, draw=cppblue!80] coordinates {
      (38,  Background)
      (28,  2 cube + 3 cyl)
      (52,  2 cylinders)
      (43,  3 cylinders)
      (68,  1 cyl + 1 cube)
      (38,  1 cube)
      (108, 1 cylinder)
      (144, Cylinder (gốc))
    };
    \end{axis}
  \end{tikzpicture}
  \caption{Phân bố tập dữ liệu theo danh mục ảnh (~534 ảnh tổng,
           633 annotation cylinder, 170 annotation cube).
           Các cảnh đa vật thể giúp mô hình học phát hiện từng đối tượng
           riêng lẻ trong môi trường phức tạp.}
  \label{fig:ai3-dataset}
\end{figure}

Annotation được thực hiện bằng \texttt{labelme} theo format bounding box hình chữ
nhật. Workflow:
\begin{enumerate}
  \item Vẽ bounding box cho từng cylinder và cube bằng \texttt{labelme}
  \item Chuyển đổi labelme JSON $\to$ COCO JSON thống nhất
        (script \texttt{labelme\_to\_coco.py})
  \item Tập dữ liệu được split theo tỉ lệ $70/15/15$ (train/val/test),
        seed cố định để đảm bảo tính tái lập
\end{enumerate}

Phân chia ba chiều thay vì hai chiều đảm bảo tập test hoàn toàn độc lập ---
không được nhìn thấy trong suốt quá trình training hoặc điều chỉnh siêu tham số.

\subsubsection*{Pipeline huấn luyện}

\begin{figure}[htbp]
  \centering
  \begin{tikzpicture}[
    node distance=0.45cm and 0.15cm,
    st/.style={draw, rounded corners=3pt, fill=gray!10, draw=gray!55,
               text width=2.0cm, align=center, minimum height=0.85cm, font=\small},
    blue_st/.style={st, fill=cppblue!10, draw=cppblue!55},
    out_st/.style={st, fill=orng!20, draw=orng!70},
    arr/.style={->, >=Stealth, thick, gray!60},
    lbl/.style={font=\scriptsize, text=gray!80},
  ]

  \node[blue_st] (s1) {Annotation\\{\scriptsize labelme}};
  \node[blue_st, right=of s1] (s2) {COCO\\JSON};
  \node[blue_st, right=of s2] (s3) {Pascal\\VOC XML};
  \node[blue_st, right=of s3] (s4) {Training\\{\scriptsize 50 epochs}};
  \node[out_st, right=of s4] (s5) {SavedModel\\{\scriptsize Float32}};
  \node[out_st, below=0.55cm of s5] (s6) {INT8\\TFLite};
  \node[out_st, below=0.55cm of s6] (s7) {EdgeTPU\\TFLite};

  \draw[arr] (s1) -- (s2) node[lbl,above,midway] {to\_coco.py};
  \draw[arr] (s2) -- (s3) node[lbl,above,midway] {tflite\_mm};
  \draw[arr] (s3) -- (s4);
  \draw[arr] (s4) -- (s5) node[lbl,above,midway] {export};
  \draw[arr] (s5) -- (s6) node[lbl,right,midway] {PTQ};
  \draw[arr] (s6) -- (s7) node[lbl,right,midway] {edgetpu\_compiler};

  \end{tikzpicture}
  \caption{Pipeline huấn luyện và xuất mô hình. Từ ảnh thô đến ba file
           mô hình phục vụ các môi trường triển khai khác nhau.
           Bước chuyển đổi sang Pascal VOC XML là yêu cầu bắt buộc của
           \texttt{tflite\_model\_maker}.}
  \label{fig:ai3-pipeline}
\end{figure}

Transfer learning với \texttt{train\_whole\_model=False} --- chỉ fine-tune detection
head, giữ nguyên backbone đã pretrain --- phù hợp với tập dữ liệu tương đối nhỏ
($\sim$374 ảnh train) và giảm nguy cơ overfitting đáng kể so với
huấn luyện từ đầu.


% ----------------------------------------------------------------
\subsection{Quantization INT8 và biên dịch cho EdgeTPU}
\label{subsec:ai3-quant}
% ----------------------------------------------------------------

Mô hình sau training là Float32 --- không chạy được trực tiếp trên Coral EdgeTPU.
EdgeTPU chỉ thực thi mô hình INT8 quantized, trong đó trọng số và activation
được biểu diễn bằng số nguyên 8-bit. Lợi ích: kích thước giảm $\sim4\times$,
và phần cứng EdgeTPU đạt tốc độ inference cực nhanh ($\sim$4--15~ms/frame)
nhờ tận dụng phép tính số nguyên.

Post-Training Quantization (PTQ) sử dụng tập train làm \textit{representative
dataset} để calibrate scale và zero\_point cho từng layer --- đây là bước bắt
buộc để duy trì độ chính xác của INT8 model gần với Float32 baseline.

File \texttt{\_edgetpu.tflite} sau khi biên dịch chứa custom op
\texttt{edgetpu-custom-op} mà TFLite runtime thông thường
\textit{không thể thực thi trên CPU}, gây lỗi:

\begin{lstlisting}[language={}]
RuntimeError: Encountered unresolved custom op: edgetpu-custom-op
\end{lstlisting}

Do đó hệ thống \textbf{duy trì hai file riêng biệt} và lựa chọn tự động:

\begin{center}
\small
\begin{tabular}{lll}
  \toprule
  File & Runtime & Deployment \\
  \midrule
  \texttt{cylinder\_detector\_int8\_edgetpu.tflite} & pycoral   & Coral EdgeTPU \\
  \texttt{cylinder\_detector\_int8.tflite}          & TFLite/ai-edge-litert & CPU \\
  \bottomrule
\end{tabular}
\end{center}


% ----------------------------------------------------------------
\subsection{Kiến trúc tích hợp: C++ ViSP ↔ Python subprocess}
\label{subsec:ai3-arch}
% ----------------------------------------------------------------

\subsubsection*{Quyết định thiết kế}

Thay vì link trực tiếp thư viện TFLite C++ vào build system của ViSP
(đòi hỏi thay đổi lớn vào CMakeLists và phụ thuộc version build cụ thể),
hệ thống áp dụng kiến trúc \textbf{subprocess model}: C++ gọi Python script
như một tiến trình con qua \texttt{popen()}, giao tiếp một chiều qua stdout.
Cách này giúp module AI hoàn toàn tách biệt --- cập nhật mô hình không cần
recompile code robot.

\begin{figure}[htbp]
  \centering
  \begin{tikzpicture}[
    node distance=0.38cm and 0.5cm,
    comp/.style={draw, rounded corners=4pt, text width=3.1cm, align=center,
                 minimum height=0.82cm, font=\small},
    cpp/.style={comp, fill=cppblue!10, draw=cppblue!55},
    py/.style={comp,  fill=pygreen!10, draw=pygreen!60},
    tmp/.style={comp,  fill=orng!18, draw=orng!65, text width=2.5cm,
                minimum height=0.7cm, font=\scriptsize},
    arr/.style={->, >=Stealth, thick},
    ]

  %% ---- C++ side ----
  \node[cpp] (cap)   {Camera capture\\{\scriptsize vpImage{$<$}uchar{$>$}}};
  \node[cpp, below=of cap]  (save)  {Save frame\\{\scriptsize cv::imwrite()}};
  \node[cpp, below=of save] (pop)   {popen() subprocess\\{\scriptsize build cmd}};
  \node[cpp, below=of pop]  (sel)   {select() timeout 3s\\{\scriptsize POSIX wait}};
  \node[cpp, below=of sel]  (parse) {Parse stdout\\{\scriptsize sscanf SUCCESS}};
  \node[cpp, below=of parse](feat)  {vpFeatureBuilder\\{\scriptsize create feature}};

  \draw[arr](cap)--(save); \draw[arr](save)--(pop);
  \draw[arr](pop)--(sel);  \draw[arr](sel)--(parse); \draw[arr](parse)--(feat);

  \node[draw=cppblue!50, dashed, rounded corners=6pt,
        fit=(cap)(save)(pop)(sel)(parse)(feat),
        inner sep=7pt] (cbox) {};
  \node[above=2pt of cbox, font=\bfseries\small, text=cppblue]
       {C++ (ViSP main loop)};

  %% ---- Temp file ----
  \node[tmp, right=2.6cm of save] (jpg)
       {/tmp/visp\_ai\_frame.jpg};

  %% ---- Python side ----
  \node[py, right=2.6cm of pop]   (load) {Load model\\{\scriptsize fallback chain}};
  \node[py, below=of load]        (pre)  {Preprocess\\{\scriptsize resize 320$\times$320}};
  \node[py, below=of pre]         (inf)  {Inference\\{\scriptsize invoke()}};
  \node[py, below=of inf]         (post) {Post-process\\{\scriptsize dequant+filter}};
  \node[py, below=of post]        (out)  {Print stdout\\{\scriptsize SUCCESS u v conf}};

  \draw[arr](load)--(pre); \draw[arr](pre)--(inf);
  \draw[arr](inf)--(post); \draw[arr](post)--(out);

  \node[draw=pygreen!55, dashed, rounded corners=6pt,
        fit=(load)(pre)(inf)(post)(out),
        inner sep=7pt] (pbox) {};
  \node[above=2pt of pbox, font=\bfseries\small, text=pygreen!80!black]
       {Python subprocess};

  %% ---- Data flow arrows ----
  \draw[arr, orng!80]  (save.east) -- (jpg.west)  node[midway,above,font=\scriptsize]{write};
  \draw[arr, orng!80]  (jpg.east)  -- (load.west) node[midway,above,font=\scriptsize]{read};
  \draw[arr, cppblue!60](pop.east)  -- (load.west) node[midway,above,font=\scriptsize,yshift=2pt]{spawn};
  \draw[arr, pygreen!70, bend right=25](out.west)  to (parse.east)
        node[midway,below,font=\scriptsize]{stdout pipe};

  \end{tikzpicture}
  \caption{Kiến trúc tích hợp C++ ViSP $\leftrightarrow$ Python subprocess.
           C++ lưu frame hiện tại ra file tạm, spawn Python làm child process
           qua \texttt{popen()}, đọc kết quả từ stdout sau khi qua timeout guard.}
  \label{fig:ai3-arch}
\end{figure}

\subsubsection*{Protocol trao đổi dữ liệu}

Output của Python script là một dòng duy nhất trên stdout:

\begin{lstlisting}[language={}]
SUCCESS 272.86 190.74 0.8362 238.42 166.89 68.88 47.68
%        u      v      conf   x_min  y_min  width  height
\end{lstlisting}

hoặc \texttt{FAILURE~<lý~do>} nếu không phát hiện được. Tất cả debug log
đi ra \texttt{stderr} --- C++ chỉ đọc stdout --- tránh nhiễu kết quả parse.

Đáng chú ý: trong ViSP, \texttt{vpImagePoint} được định nghĩa là \texttt{(row,~col)}
tương ứng \texttt{(v,~u)}, không phải \texttt{(u,~v)} như thông thường:

\begin{lstlisting}[language=C++]
detected_center = vpImagePoint(v, u);  // row=v, col=u
\end{lstlisting}

Tham số \texttt{hint} cho phép Python chọn cylinder \textit{gần nhất với vị
trí đã biết} thay vì cylinder có confidence cao nhất --- cơ chế này tránh
hiện tượng servo ``nhảy'' sang cylinder khác khi scores thay đổi giữa các frame.


% ----------------------------------------------------------------
\subsection{State machine: PREINIT \texorpdfstring{$\to$}{→} INIT (AI) \texorpdfstring{$\to$}{→} JOINT \texorpdfstring{$\to$}{→} ...}
\label{subsec:ai3-sm}
% ----------------------------------------------------------------

Hệ thống vận hành theo vòng lặp chính liên tục với 6 trạng thái.
Mỗi iteration đọc một frame camera và thực thi logic trạng thái hiện tại.

\begin{figure}[htbp]
  \centering
  \begin{tikzpicture}[
    node distance=1.0cm and 2.0cm,
    st/.style={draw, rounded corners=4pt, fill=cppblue!10, draw=cppblue!55,
               text width=2.3cm, align=center, minimum height=1.0cm, font=\small},
    ai/.style={st, fill=pygreen!10, draw=pygreen!60},
    arr/.style={->, >=Stealth, thick},
    lbl/.style={font=\scriptsize, text=gray!70},
  ]

  %% Top row
  \node[st]  (pre)  {PREINIT\\{\scriptsize home pose}};
  \node[ai,  right=of pre]   (ini)  {INIT\\{\scriptsize AI detect}};
  \node[ai,  right=of ini]   (jnt)  {JOINT\\{\scriptsize AI tracking}};
  \node[st,  right=of jnt]   (app)  {APPROACH\\{\scriptsize chờ RC5}};

  %% Bottom row
  \node[st,  below=1.8cm of app]  (grp)  {GRIPPER\\{\scriptsize đóng gripper}};
  \node[st,  left=of grp]         (cls)  {CLASSIFIED\\{\scriptsize vận chuyển}};

  %% Arrows
  \draw[arr](pre)--(ini)  node[lbl,above,midway]{home\\reached};
  \draw[arr](ini)--(jnt)  node[lbl,above,midway]{AI init\\ok};
  \draw[arr](jnt)--(app)  node[lbl,above,midway]{$|e|<5{\times}10^{-3}$};
  \draw[arr](app)--(grp)  node[lbl,right,midway]{RC5\\signal};
  \draw[arr](grp)--(cls)  node[lbl,below,midway]{grasped};
  \draw[arr](cls.west) to[out=180,in=270,looseness=0.8] (pre.south)
        node[lbl,left,pos=0.5]{cycle\\restart};

  %% AI self-recovery
  \draw[arr, pygreen!70, dashed]
        (jnt.north) to[out=100,in=80,looseness=2.5] (jnt.north)
        node[above=0.7cm of jnt, font=\scriptsize, text=pygreen!80!black]
        {AI re-detect};

  %% Fallback from INIT
  \draw[arr, warnred!70, dashed]
        (ini.south) -- ++(0,-0.7)
        node[below, font=\scriptsize, text=warnred]
        {fallback: click thủ công};

  \end{tikzpicture}
  \caption{State machine của hệ thống pick \& place sau cải tiến.
           Trạng thái tô xanh lá là nơi AI được gọi.
           Vòng lặp tự phục hồi ở JOINT (mũi tên đứt nét) cho phép hệ thống
           tiếp tục mà không cần reset khi mất detection.}
  \label{fig:ai3-sm}
\end{figure}

\paragraph{PREINIT.}
Robot gửi lệnh joint position về home pose
$[0^\circ, 0^\circ, 90^\circ, 0^\circ, 90^\circ, 0^\circ]$.
Servo task được cấu hình với interaction matrix Desired, pseudo-inverse,
gain $\lambda = 0.4$.

\paragraph{INIT.}
Chờ robot đạt home pose (sai số từng joint $< 0.01$) và gripper mở xong.
Khi cả hai điều kiện thỏa, \texttt{detectCylinderWithAI()} được gọi
không có hint --- Python chọn cylinder có confidence cao nhất.
Kết quả trực tiếp tạo visual feature \texttt{vpFeaturePoint}
và add vào servo task. Nếu AI thất bại, \texttt{ai\_hint}
giữ giá trị mặc định (fallback thủ công).

\paragraph{JOINT.}
Khác với cách tiếp cận truyền thống dùng blob tracker mỗi frame,
hệ thống hiện tại \textbf{gọi AI trực tiếp ở mỗi iteration}
với \texttt{hint = \&ai\_hint}. Control law:
\begin{equation}
  \mathbf{v} = -\lambda\,\mathbf{L}_{s^*}^{+}\,\mathbf{e},
  \qquad \mathbf{q}_{k+1} = \mathbf{q}_k + \mathbf{v}\,\Delta t
  \label{eq:ibvs}
\end{equation}
với $\Delta t = 0.5\,$s. Điều kiện hội tụ: $|e_x| < 5{\times}10^{-3}$
và $|e_y| < 5{\times}10^{-3}$ (normalized image coordinates).
Khi đạt, gửi \texttt{"OKE\textbackslash r"} qua UART đến RC5 controller.


% ----------------------------------------------------------------
\subsection{Các vấn đề kỹ thuật và giải pháp}
\label{subsec:ai3-issues}
% ----------------------------------------------------------------

Một số vấn đề nảy sinh trong quá trình tích hợp mà khó thấy trước khi thực
sự chạy hệ thống.

\subsubsection*{P1: Xác định output tensor sai lệch}

\texttt{tflite\_model\_maker} 0.4.x đặt tên tensor theo thứ tự riêng của nó:

\begin{center}
\small
\begin{tabular}{ccc}
  \toprule
  Tensor & Shape & Nội dung \\
  \midrule
  \texttt{:0} & \texttt{[1]}     & num\_detections \\
  \texttt{:1} & \texttt{[1,\,N]} & scores \quad \textbf{← cần lấy cái này} \\
  \texttt{:2} & \texttt{[1,\,N]} & class IDs \quad \textbf{← không phải cái này} \\
  \texttt{:3} & \texttt{[1,\,N,\,4]} & boxes \\
  \bottomrule
\end{tabular}
\end{center}

Thứ tự này không khớp với API của pycoral, vốn giả định positional order
\texttt{boxes/classes/scores/count = 0/1/2/3}. Dùng API đó sẽ đọc
num\_detections (raw uint8 = 255) như scores và crash với \texttt{IndexError}.
Giải pháp: xác định tensor \textbf{theo shape động} --- boxes là
\texttt{[1,N,4]}, scores là tensor \texttt{[1,N]} có name suffix nhỏ nhất
(\texttt{:1} chứ không phải \texttt{:2}).

\subsubsection*{P2: Python 3.10 không có \texttt{tflite-runtime}}

Google không publish wheel \texttt{tflite-runtime} cho Python 3.10+.
Giải pháp là xây dựng \textbf{fallback chain 4 tầng}:

\begin{figure}[htbp]
  \centering
  \begin{tikzpicture}[
    node distance=0.4cm,
    fb/.style={draw, rounded corners=3pt, text width=7.5cm, align=left,
               minimum height=0.85cm, font=\small, inner sep=7pt},
    t1/.style={fb, fill=purple!12, draw=purple!55},
    t2/.style={fb, fill=teal!10,   draw=teal!50},
    t3/.style={fb, fill=cppblue!10, draw=cppblue!50},
    t4/.style={fb, fill=gray!12,   draw=gray!50},
    fa/.style={fb, fill=warnred!10, draw=warnred!50},
    arr/.style={->, >=Stealth, thick, dashed, gray!55},
    lbl/.style={font=\scriptsize, text=warnred!80},
  ]
  \node[t1] (a) {\textbf{1.}~\texttt{pycoral.utils.edgetpu}\quad{\scriptsize Coral EdgeTPU, Python 3.9}};
  \node[t2, below=of a] (b) {\textbf{2.}~\texttt{ai\_edge\_litert.interpreter}\quad{\scriptsize Google's new pkg, Python 3.8--3.12}};
  \node[t3, below=of b] (c) {\textbf{3.}~\texttt{tflite\_runtime.interpreter}\quad{\scriptsize Legacy, Python 3.9}};
  \node[t4, below=of c] (d) {\textbf{4.}~\texttt{tensorflow.lite.Interpreter}\quad{\scriptsize Heavy fallback}};
  \node[fa, below=of d] (e) {\textbf{FAILURE}:~\texttt{no\_tflite\_backend\_available}};

  \draw[arr](a)--(b) node[lbl,right,midway]{ImportError / Coral not connected};
  \draw[arr](b)--(c) node[lbl,right,midway]{ImportError};
  \draw[arr](c)--(d) node[lbl,right,midway]{ImportError};
  \draw[arr](d)--(e) node[lbl,right,midway]{ImportError};
  \end{tikzpicture}
  \caption{Python backend fallback chain trong \texttt{detect\_cylinder.py}.
           Script tự động thử từng backend theo thứ tự ưu tiên,
           tương thích với cả môi trường có Coral lẫn CPU-only.}
  \label{fig:ai3-fallback}
\end{figure}

\texttt{ai-edge-litert} là gói chính thức của Google thay thế \texttt{tflite-runtime},
hỗ trợ Python 3.8--3.12 với API giống hệt nhau --- là giải pháp cho vấn đề Python 3.10.

\subsubsection*{P3: JSON comment trong config}

\texttt{config.json} dùng comment \texttt{//} để ghi chú, nhưng cả Python
\texttt{json.load()} và \texttt{nlohmann::json} đều từ chối mặc định.
Phía C++: \texttt{nlohmann::json::parse(file, nullptr, true, true)} với
\texttt{ignore\_comments=true} (có từ nlohmann/json~3.9.0).
Phía Python: strip comment lines thủ công trước khi \texttt{json.loads()}.


% ----------------------------------------------------------------
\subsection{Thông số cấu hình visual servo}
\label{subsec:ai3-servo}
% ----------------------------------------------------------------

Hệ thống sử dụng \textbf{Image-Based Visual Servoing (IBVS)} với feature là
tọa độ 2D của cylinder trong không gian ảnh normalized $\mathbf{s}=(x,y)$,
desired position $\mathbf{s}^*=(0,0,1)$ --- cylinder phải nằm đúng tâm ảnh.

\begin{table}[htbp]
  \centering
  \caption{Thông số điều khiển visual servo}
  \label{tab:ai3-servo}
  \small
  \begin{tabular}{ll}
    \toprule
    \textbf{Tham số}            & \textbf{Giá trị} \\
    \midrule
    Loại điều khiển             & Eye-in-hand, \texttt{EYEINHAND\_L\_cVe\_eJe} \\
    Interaction matrix          & Desired ($\mathbf{L}_{s^*}$), pseudo-inverse \\
    Gain $\lambda$              & 0.4 \\
    Bước thời gian $\Delta t$   & 0.5\,s \\
    Ngưỡng hội tụ               & $|e_x| < 5{\times}10^{-3}$ và $|e_y| < 5{\times}10^{-3}$ \\
    Desired position            & $(x^*, y^*) = (0,\,0)$ --- tâm ảnh, $Z = 1$ \\
    \bottomrule
  \end{tabular}
\end{table}

Dùng $\mathbf{L}_{s^*}$ (tại vị trí mong muốn) thay vì $\mathbf{L}_{s_{cur}}$
giúp tránh singularity khi feature tiến gần desired point và giữ ổn định
theo ổn định hơn trong giai đoạn hội tụ. Velocity được saturate trước khi
gửi xuống robot để đảm bảo không vượt giới hạn vật lý từng khớp.
   %% mục 4.4
%%     ...
%%     \chapter{Kết quả và Đánh giá}
%%     ...
%%     %%=============================================================
%% chapter5_4.tex
%% Phần 5.4 — Kết quả đánh giá: AI #3 Phát hiện vật thể
%%
%% \input file — cần thêm vào preamble:
%%   \usepackage{tikz}
%%   \usetikzlibrary{arrows.meta,positioning,patterns}
%%   \usepackage{pgfplots}
%%   \pgfplotsset{compat=1.18}
%%   \usepackage{booktabs,multirow,subcaption,xcolor}
%%=============================================================

% ================================================================
\section{Kết quả và Đánh giá: AI \#3 --- Phát hiện vật thể}
\label{sec:ai3-results}
% ================================================================

% ----------------------------------------------------------------
\subsection{Kết quả detection model: COCO metrics}
\label{subsec:ai3-coco}
% ----------------------------------------------------------------

Mô hình được đánh giá trên cả validation set và test set độc lập
theo chuẩn \textbf{COCO evaluation protocol} --- tính trung bình mAP qua
các ngưỡng IoU từ 0.50 đến 0.95 với bước 0.05.

\begin{figure}[htbp]
  \centering
  \begin{tikzpicture}
    \begin{axis}[
      ybar, bar width=0.45cm,
      width=10cm, height=7cm,
      symbolic x coords={mAP (0.50:0.95), AP50, AP75, AR},
      xtick=data,
      x tick label style={font=\small, align=center},
      ymin=0, ymax=1.15,
      ylabel={Score},
      ylabel style={font=\small},
      legend style={at={(0.5,1.03)}, anchor=south, legend columns=2, font=\small},
      legend cell align=left,
      nodes near coords,
      every node near coord/.style={font=\scriptsize, /pgf/number format/fixed,
                                    /pgf/number format/precision=3},
      enlarge x limits=0.25,
      grid=major, grid style={dashed, gray!30},
      axis line style={gray!60},
      tick style={gray!60},
    ]
    \addplot[fill=cppblue!70, draw=cppblue!90] coordinates {
      (mAP (0.50:0.95), 0.791)
      (AP50, 0.990)
      (AP75, 0.990)
      (AR,   0.869)
    };
    \addplot[fill=orng!80, draw=orng!90] coordinates {
      (mAP (0.50:0.95), 0.799)
      (AP50, 1.000)
      (AP75, 0.969)
      (AR,   0.856)
    };
    \legend{Validation, Test}
    \end{axis}
  \end{tikzpicture}
  \caption{COCO metrics trên validation set và test set. AP50 trên test set
           đạt 1.000 cho thấy 100\% cylinder được phát hiện đúng ở ngưỡng IoU
           lỏng lẻ (50\%). mAP$\approx$0.80 phản ánh vẫn còn đây đó những trường hợp
           định vị bounding box chưa được chính xác ở ngưỡng IoU cao.}
  \label{fig:ai3-coco}
\end{figure}

Bảng~\ref{tab:ai3-metrics} tổng hợp số liệu chi tiết:

\begin{table}[htbp]
  \centering
  \caption{COCO metrics trên validation và test set}
  \label{tab:ai3-metrics}
  \small
  \begin{tabular}{lcc}
    \toprule
    \textbf{Metric}                    & \textbf{Validation} & \textbf{Test} \\
    \midrule
    mAP (IoU 0.50:0.95)               & 0.791               & \textbf{0.799} \\
    AP50 (IoU $\geq$ 0.50)            & 0.990               & \textbf{1.000} \\
    AP75 (IoU $\geq$ 0.75)            & 0.990               & 0.969 \\
    AR (maxDets=100)                   & 0.869               & 0.856 \\
    \midrule
    Precision (IoU=0.50, conf=0.50)    & ---                 & \textbf{1.000} \\
    Recall    (IoU=0.50, conf=0.50)    & ---                 & 0.990 \\
    F1-Score                           & ---                 & \textbf{0.995} \\
    TP / FP / FN                       & ---                 & 95 / 0 / 1 \\
    \bottomrule
  \end{tabular}
\end{table}

\begin{figure}[htbp]
  \centering
  \begin{tikzpicture}
    \begin{axis}[
      ybar, bar width=0.55cm,
      width=8.5cm, height=6cm,
      symbolic x coords={Precision, Recall, F1-Score},
      xtick=data,
      x tick label style={font=\small},
      ymin=0.96, ymax=1.04,
      ylabel={Score},
      ylabel style={font=\small},
      nodes near coords,
      every node near coord/.style={font=\small, /pgf/number format/fixed,
                                    /pgf/number format/precision=3},
      enlarge x limits=0.3,
      grid=major, grid style={dashed, gray!30},
      axis line style={gray!60},
    ]
    \addplot[fill=pygreen!65, draw=pygreen!80] coordinates {
      (Precision, 1.000)
      (Recall,    0.990)
      (F1-Score,  0.995)
    };
    \end{axis}
  \end{tikzpicture}
  \caption{Precision, Recall và F1-Score trên test set (PASCAL VOC protocol,
           IoU$\geq$0.50, confidence$\geq$0.50, chỉ lớp cylinder).
           Precision = 1.000 và FP = 0 chứng tỏ không có false positive
           nào ở ngưỡng 0.50. FN = 1 (một cylinder bị bỏ sót) giải thích
           Recall = 0.990.}
  \label{fig:ai3-prf}
\end{figure}

Giá trị AP50 = 1.000 trên test set à ngưỡng IoU lỏng lẻ (50\%) phản ánh
mô hình phát hiện được toàn bộ cylinder. mAP = 0.799 (trung bình
qua IoU 0.50--0.95) thấp hơn một chút do mô hình EfficientDet-Lite0
đơn giản vẫn còn giới hạn trong việc định vị bounding box chính xác
ở các ngưỡng IoU cao ($\geq$0.75) --- điều phổ biến với các kiến trúc
nhe hơn.


% ----------------------------------------------------------------
\subsection{Phân tích confidence score: true positive vs false positive}
\label{subsec:ai3-conf}
% ----------------------------------------------------------------

Một trong những điều quan trọng khi triển khai thực tế là khoảng cách
phân tách giữa true positive và false positive --- điều này quyết định
ngưỡng confidence nào thực tế dùng được mà không sai kết quả.

Kiểm tra trên 218 ảnh cylinder (true positives) và 10 ảnh cube
(false positives):

\begin{figure}[htbp]
  \centering
  \begin{tikzpicture}
    \begin{axis}[
      width=11cm, height=5.5cm,
      xlabel={Confidence score},
      ylabel={},
      xmin=0, xmax=1.05,
      ymin=0, ymax=3,
      xtick={0,0.1,...,1.0},
      ytick={0.7, 2.0},
      yticklabels={\small False Positive\\(cube, $n$=10), \small True Positive\\(cylinder, $n$=218)},
      y tick label style={align=right, font=\small},
      x tick label style={font=\small},
      xlabel style={font=\small},
      grid=both, grid style={dashed, gray!25},
      axis line style={gray!60},
    ]

    %% --- True Positive band (y=2) ---
    %% Most scores: 0.85-0.98 (majority)
    \addplot[only marks, mark=|, mark size=5pt, mark options={draw=pygreen!70, line width=1.2pt},
             opacity=0.5]
    coordinates {
      (0.29,2)(0.61,2)(0.74,2)
      (0.82,2)(0.83,2)(0.84,2)(0.85,2)(0.86,2)(0.87,2)(0.88,2)
      (0.89,2)(0.90,2)(0.91,2)(0.92,2)(0.93,2)(0.94,2)(0.94,2)
      (0.95,2)(0.95,2)(0.96,2)(0.96,2)(0.97,2)(0.97,2)(0.98,2)
    };

    %% --- False Positive band (y=0.7) ---
    \addplot[only marks, mark=|, mark size=5pt, mark options={draw=warnred!70, line width=1.2pt},
             opacity=0.6]
    coordinates {
      (0.08,0.7)(0.13,0.7)(0.17,0.7)(0.21,0.7)
      (0.25,0.7)(0.28,0.7)(0.30,0.7)(0.33,0.7)(0.35,0.7)(0.36,0.7)
    };

    %% --- Threshold line at 0.50 ---
    \addplot[thick, dashed, draw=gray!80, domain=0:3]
    coordinates {(0.50, 0) (0.50, 3)};
    \node[font=\scriptsize, text=gray!80, rotate=90]
         at (axis cs: 0.52, 1.5) {threshold = 0.50};

    %% --- Gap annotation ---
    \draw[->, >=Stealth, thick, gray!70]
         (axis cs:0.38, 2.5) -- (axis cs:0.48, 2.5);
    \draw[->, >=Stealth, thick, gray!70]
         (axis cs:0.34, 2.5) -- (axis cs:0.27, 2.5);
    \node[font=\scriptsize, text=gray!80] at (axis cs:0.37, 2.7) {gap = 0.34};

    %% --- Median annotation for TP ---
    \addplot[only marks, mark=diamond*, mark size=4pt,
             mark options={draw=pygreen!80!black, fill=pygreen!60}]
    coordinates {(0.94, 2)};
    \node[font=\scriptsize, text=pygreen!80!black, above]
         at (axis cs:0.94, 2) {median 0.94};

    %% --- Max FP annotation ---
    \addplot[only marks, mark=diamond*, mark size=4pt,
             mark options={draw=warnred!80, fill=warnred!50}]
    coordinates {(0.36, 0.7)};
    \node[font=\scriptsize, text=warnred, below]
         at (axis cs:0.36, 0.7) {max 0.36};

    \end{axis}
  \end{tikzpicture}
  \caption{Phân tách confidence score giữa true positive (cylinder, xanh lá)
           và false positive (cube, đỏ). Mỗi vạch đứng đại diện một detection.
           Ngưỡng 0.50 nằm chính giữa khoảng gap 0.34,
           loại bỏ toàn bộ false positive mà vẫn chấp nhận 99.5\% true positive.}
  \label{fig:ai3-conf}
\end{figure}

Khoảng cách 0.34 giữa điểm false positive cao nhất (0.36) và 99.1\%
true positive ($\geq$0.70) xác nhận mô hình đã học được đặc trưng
phân biệt mạnh giữa cylinder và cube. Ngưỡng hiện tại 0.50 nằm an toàn
giữa hai vùng phân bố này.

Ba ảnh outlier có score thấp (0.29, 0.61, 0.74) tương ứng với cylinder
bị nhìn từ góc bất thường hoặc bị che khuất một phần. Ỡ runtime,
các trường hợp này fallback về click thủ công --- hành vi an toàn.

\begin{table}[htbp]
  \centering
  \caption{Phân tích confidence score trên 218 ảnh cylinder + 10 ảnh cube}
  \label{tab:ai3-conf}
  \small
  \begin{tabular}{lcc}
    \toprule
    & \textbf{True Positive (cylinder)} & \textbf{False Positive (cube)} \\
    \midrule
    Score range      & 0.29 -- 0.98          & 0.08 -- 0.36 \\
    Median           & \textbf{0.94}         & --- \\
    $\geq$ 0.50      & 99.5\% (217/218)      & \textbf{0/10 (0\%)} \\
    $\geq$ 0.70      & 99.1\% (216/218)      & 0/10 (0\%) \\
    $\geq$ 0.80      & 95.9\% (209/218)      & 0/10 (0\%) \\
    \midrule
    Gap ph\^an t\'ach & \multicolumn{2}{c}{$0.70 - 0.36 = 0.34$} \\
    \bottomrule
  \end{tabular}
\end{table}


% ----------------------------------------------------------------
\subsection{Kết quả tích hợp hệ thống}
\label{subsec:ai3-integration}
% ----------------------------------------------------------------

\begin{table}[htbp]
  \centering
  \caption{So sánh khả năng trước và sau khi tích hợp AI}
  \label{tab:ai3-comparison}
  \small
  \begin{tabular}{lll}
    \toprule
    \textbf{Tiêu chí}                    & \textbf{Hệ thống gốc} & \textbf{Sau cải tiến} \\
    \midrule
    Khởi tạo mỗi chu kỳ               & Click thủ công (bắt buộc) & Tự động (AI) \\
    Tỷ lệ init thành công             & 100\% (do manual)  & $\sim$99.5\% (AI) \\
    Fallback khi AI thất bại           & Không có              & Click thủ công \\
    Phục hồi khi mất tracking          & Reset toàn bộ        & AI re-detect tự động \\
    Chạy không giám sát                & Không thể              & Có (trường hợp bình thường) \\
    \bottomrule
  \end{tabular}
\end{table}

Thiết kế hiện tại gọi AI trực tiếp \textbf{mỗi iteration} của servo
loop thay vì dùng blob tracker. Điều này có một hàm ý đáng chú ý:
\textit{không có tracking failure} theo nghĩa truyền thống --- mỗi frame
là một lần detection độc lập. Tính ổn định của servo phụ thuộc vào
sự nhất quán của AI detection qua các frame liên tiếp. Cơ chế \texttt{--hint}
(truyền vị trí frame trước) giúp được phần nào khi có nhiều cylinder
trong cảnh.


% ----------------------------------------------------------------
\subsection{Phân tích latency: subprocess CPU vs EdgeTPU}
\label{subsec:ai3-latency}
% ----------------------------------------------------------------

Mỗi lần gọi \texttt{detectCylinderWithAI()} bao gồm nhiều thành phần
latency khác nhau. Điều ít người để ý là: phần overhead lớn nhất
không phải là inference mà là \textbf{Python startup time} --- chi phí
mở một process mới mỗi lần gọi \texttt{popen()}.

\begin{figure}[htbp]
  \centering
  \begin{tikzpicture}
    \begin{axis}[
      xbar stacked,
      bar width=0.55cm,
      width=11cm, height=6.5cm,
      symbolic y coords={
        {EdgeTPU\\(persistent, tương lai)},
        {CPU\\(persistent, tương lai)},
        {EdgeTPU\\(subprocess hiện tại)},
        {CPU\\(subprocess hiện tại)}
      },
      ytick=data,
      y tick label style={align=right, font=\small},
      xlabel={Thời gian (ms)},
      xlabel style={font=\small},
      xmin=0, xmax=780,
      legend style={at={(0.5,-0.18)}, anchor=north, legend columns=4,
                    font=\scriptsize},
      legend cell align=left,
      grid=major, grid style={dashed, gray!25},
      axis line style={gray!60},
      enlarge y limits=0.18,
      nodes near coords, nodes near coords align={horizontal},
      point meta=explicit symbolic,
      every node near coord/.style={font=\scriptsize},
    ]
    %%  stacks: I/O | Inference | Model load | Python startup
    %% CPU subprocess: 10 + 150 + 100 + 350 = 610
    \addplot[fill=gray!40, draw=gray!60] coordinates {
      (10, {CPU\\(subprocess hiện tại)})
      (10, {EdgeTPU\\(subprocess hiện tại)})
      (10, {CPU\\(persistent, tương lai)})
      (5,  {EdgeTPU\\(persistent, tương lai)})
    };  %% I/O
    \addplot[fill=warnred!55, draw=warnred!70] coordinates {
      (150, {CPU\\(subprocess hiện tại)})
      (10,  {EdgeTPU\\(subprocess hiện tại)})
      (150, {CPU\\(persistent, tương lai)})
      (10,  {EdgeTPU\\(persistent, tương lai)})
    };  %% Inference
    \addplot[fill=orng!65, draw=orng!80] coordinates {
      (100, {CPU\\(subprocess hiện tại)})
      (50,  {EdgeTPU\\(subprocess hiện tại)})
      (0,   {CPU\\(persistent, tương lai)})
      (0,   {EdgeTPU\\(persistent, tương lai)})
    };  %% Model load
    \addplot[fill=cppblue!55, draw=cppblue!70] coordinates {
      (350, {CPU\\(subprocess hiện tại)})
      (350, {EdgeTPU\\(subprocess hiện tại)})
      (0,   {CPU\\(persistent, tương lai)})
      (0,   {EdgeTPU\\(persistent, tương lai)})
    };  %% Python startup
    \legend{I/O (JPEG), Inference, Model load, Python startup}
    \end{axis}
  \end{tikzpicture}
  \caption{Phân tích latency mỗi lần gọi AI detection.
           Python startup time ($\sim$350~ms) chiếm phần lẫn tổng overhead
           trong kiến trúc subprocess hiện tại --- lớn hơn nhiều so với
           thời gian inference thực tế. Kiến trúc persistent process (tương lai)
           loại bỏ hoàn toàn chi phí này.}
  \label{fig:ai3-latency}
\end{figure}

\begin{table}[htbp]
  \centering
  \caption{So sánh tổng latency giữa các cấu hình triển khai}
  \label{tab:ai3-latency}
  \small
  \begin{tabular}{lcc}
    \toprule
    \textbf{Cấu hình}                    & \textbf{Inference} & \textbf{Tổng / lần gọi} \\
    \midrule
    CPU subprocess (hiện tại)            & $\sim$100--200~ms  & $\sim$600--700~ms \\
    EdgeTPU subprocess (hiện tại)       & $\sim$4--15~ms     & $\sim$420--430~ms \\
    CPU persistent (tương lai)           & $\sim$100--200~ms  & $\sim$160--210~ms \\
    EdgeTPU persistent (tương lai)      & $\sim$4--15~ms     & $\sim$15--20~ms  \\
    \bottomrule
  \end{tabular}
\end{table}

Kết quả trên cho thấy rõ điểm nghẻn hiện tại:
EdgeTPU inference chỉ mất 4--15~ms nhưng toàn bộ lá gọi
vẫn tốn 420~ms do phải khởi động Python mỗi lần. Với
kiến trúc \textbf{persistent Python process} --- một process
duy nhất khởi động một lần, load mô hình vào bộ nhớ, sau
đó tiếp nhận request từ C++ qua stdin/stdout pipe --- latency
mỗi lần gọi giảm xuống còn $\sim$15--20~ms với EdgeTPU,
mở ra khả năng chạy AI detection \textit{mỗi frame} camera
thay vì chỉ khi cần.

% Chú thích về nguồn số liệu
\smallskip
\noindent\textit{\small
Lưu ý: COCO metrics trong Bảng~\ref{tab:ai3-metrics} và Hình~\ref{fig:ai3-coco}
lấy từ kết quả đánh giá hai lớp (cylinder + cube) trên repository huấn luyện.
Phân tích confidence score trong Bảng~\ref{tab:ai3-conf} lấy từ kiểm tra
độc lập trên 218 ảnh cylinder và 10 ảnh cube (CPU inference via ai-edge-litert).
}
   %% mục 5.4
%%     ...
%%   \end{document}
%%
%% Compile: pdflatex -> pdflatex (chạy 2 lần để có label/ref đúng)
%%=============================================================

%% Hoặc copy thẳng các \usepackage cần thiết vào preamble của main.
%%=============================================================

%% Tiếng Việt
\usepackage[T5]{fontenc}
\usepackage[utf8]{inputenc}
\usepackage[vietnamese]{babel}

%% TikZ và các thư viện vẻ hình
\usepackage{tikz}
\usetikzlibrary{
  shapes.geometric,
  arrows.meta,
  positioning,
  fit,
  backgrounds,
  calc,
  decorations.pathreplacing,
  patterns,
}

%% PGFPlots cho biểu đồ
\usepackage{pgfplots}
\pgfplotsset{compat=1.18}
\usepackage{pgfplotstable}

%% Code listings
\usepackage{listings}
\usepackage{xcolor}

%% Màu dùng chung (cần định nghĩa trước khi \input chapter)
\definecolor{cppblue}{RGB}{30,100,200}
\definecolor{pygreen}{RGB}{40,150,80}
\definecolor{warnred}{RGB}{200,60,60}
\definecolor{orng}{RGB}{220,130,30}

%% Cài đặt listings mặc định
\lstset{
  basicstyle=\ttfamily\small,
  keywordstyle=\color{cppblue},
  commentstyle=\color{gray!70},
  stringstyle=\color{orng},
  breaklines=true,
  frame=single,
  rulecolor=\color{gray!40},
  framexleftmargin=4pt,
  xleftmargin=8pt,
  numbers=none,
}

%% Hình ảnh
\usepackage{graphicx}

%% Bảng đẹp hơn
\usepackage{booktabs}
\usepackage{multirow}
\usepackage{subcaption}

%% Đảm bảo float không trọn ra ngoài section
\usepackage[section]{placeins}

%%=============================================================
%% Cách sử dụng trong main.tex:
%%
%%   \documentclass[12pt,a4paper]{report}
%%   %%=============================================================
%% preamble_ai3.tex
%% Preamble dùng chung cho chapter4_4.tex và chapter5_4.tex
%%
%% Đưa vào main.tex bằng:  %%=============================================================
%% preamble_ai3.tex
%% Preamble dùng chung cho chapter4_4.tex và chapter5_4.tex
%%
%% Đưa vào main.tex bằng:  %%=============================================================
%% preamble_ai3.tex
%% Preamble dùng chung cho chapter4_4.tex và chapter5_4.tex
%%
%% Đưa vào main.tex bằng:  \input{preamble_ai3}
%% Hoặc copy thẳng các \usepackage cần thiết vào preamble của main.
%%=============================================================

%% Tiếng Việt
\usepackage[T5]{fontenc}
\usepackage[utf8]{inputenc}
\usepackage[vietnamese]{babel}

%% TikZ và các thư viện vẻ hình
\usepackage{tikz}
\usetikzlibrary{
  shapes.geometric,
  arrows.meta,
  positioning,
  fit,
  backgrounds,
  calc,
  decorations.pathreplacing,
  patterns,
}

%% PGFPlots cho biểu đồ
\usepackage{pgfplots}
\pgfplotsset{compat=1.18}
\usepackage{pgfplotstable}

%% Code listings
\usepackage{listings}
\usepackage{xcolor}

%% Màu dùng chung (cần định nghĩa trước khi \input chapter)
\definecolor{cppblue}{RGB}{30,100,200}
\definecolor{pygreen}{RGB}{40,150,80}
\definecolor{warnred}{RGB}{200,60,60}
\definecolor{orng}{RGB}{220,130,30}

%% Cài đặt listings mặc định
\lstset{
  basicstyle=\ttfamily\small,
  keywordstyle=\color{cppblue},
  commentstyle=\color{gray!70},
  stringstyle=\color{orng},
  breaklines=true,
  frame=single,
  rulecolor=\color{gray!40},
  framexleftmargin=4pt,
  xleftmargin=8pt,
  numbers=none,
}

%% Hình ảnh
\usepackage{graphicx}

%% Bảng đẹp hơn
\usepackage{booktabs}
\usepackage{multirow}
\usepackage{subcaption}

%% Đảm bảo float không trọn ra ngoài section
\usepackage[section]{placeins}

%%=============================================================
%% Cách sử dụng trong main.tex:
%%
%%   \documentclass[12pt,a4paper]{report}
%%   \input{preamble_ai3}
%%   \begin{document}
%%     ...
%%     \chapter{Thiết kế và Hiện thực}
%%     ...
%%     \input{chapter4_4}   %% mục 4.4
%%     ...
%%     \chapter{Kết quả và Đánh giá}
%%     ...
%%     \input{chapter5_4}   %% mục 5.4
%%     ...
%%   \end{document}
%%
%% Compile: pdflatex -> pdflatex (chạy 2 lần để có label/ref đúng)
%%=============================================================

%% Hoặc copy thẳng các \usepackage cần thiết vào preamble của main.
%%=============================================================

%% Tiếng Việt
\usepackage[T5]{fontenc}
\usepackage[utf8]{inputenc}
\usepackage[vietnamese]{babel}

%% TikZ và các thư viện vẻ hình
\usepackage{tikz}
\usetikzlibrary{
  shapes.geometric,
  arrows.meta,
  positioning,
  fit,
  backgrounds,
  calc,
  decorations.pathreplacing,
  patterns,
}

%% PGFPlots cho biểu đồ
\usepackage{pgfplots}
\pgfplotsset{compat=1.18}
\usepackage{pgfplotstable}

%% Code listings
\usepackage{listings}
\usepackage{xcolor}

%% Màu dùng chung (cần định nghĩa trước khi \input chapter)
\definecolor{cppblue}{RGB}{30,100,200}
\definecolor{pygreen}{RGB}{40,150,80}
\definecolor{warnred}{RGB}{200,60,60}
\definecolor{orng}{RGB}{220,130,30}

%% Cài đặt listings mặc định
\lstset{
  basicstyle=\ttfamily\small,
  keywordstyle=\color{cppblue},
  commentstyle=\color{gray!70},
  stringstyle=\color{orng},
  breaklines=true,
  frame=single,
  rulecolor=\color{gray!40},
  framexleftmargin=4pt,
  xleftmargin=8pt,
  numbers=none,
}

%% Hình ảnh
\usepackage{graphicx}

%% Bảng đẹp hơn
\usepackage{booktabs}
\usepackage{multirow}
\usepackage{subcaption}

%% Đảm bảo float không trọn ra ngoài section
\usepackage[section]{placeins}

%%=============================================================
%% Cách sử dụng trong main.tex:
%%
%%   \documentclass[12pt,a4paper]{report}
%%   %%=============================================================
%% preamble_ai3.tex
%% Preamble dùng chung cho chapter4_4.tex và chapter5_4.tex
%%
%% Đưa vào main.tex bằng:  \input{preamble_ai3}
%% Hoặc copy thẳng các \usepackage cần thiết vào preamble của main.
%%=============================================================

%% Tiếng Việt
\usepackage[T5]{fontenc}
\usepackage[utf8]{inputenc}
\usepackage[vietnamese]{babel}

%% TikZ và các thư viện vẻ hình
\usepackage{tikz}
\usetikzlibrary{
  shapes.geometric,
  arrows.meta,
  positioning,
  fit,
  backgrounds,
  calc,
  decorations.pathreplacing,
  patterns,
}

%% PGFPlots cho biểu đồ
\usepackage{pgfplots}
\pgfplotsset{compat=1.18}
\usepackage{pgfplotstable}

%% Code listings
\usepackage{listings}
\usepackage{xcolor}

%% Màu dùng chung (cần định nghĩa trước khi \input chapter)
\definecolor{cppblue}{RGB}{30,100,200}
\definecolor{pygreen}{RGB}{40,150,80}
\definecolor{warnred}{RGB}{200,60,60}
\definecolor{orng}{RGB}{220,130,30}

%% Cài đặt listings mặc định
\lstset{
  basicstyle=\ttfamily\small,
  keywordstyle=\color{cppblue},
  commentstyle=\color{gray!70},
  stringstyle=\color{orng},
  breaklines=true,
  frame=single,
  rulecolor=\color{gray!40},
  framexleftmargin=4pt,
  xleftmargin=8pt,
  numbers=none,
}

%% Hình ảnh
\usepackage{graphicx}

%% Bảng đẹp hơn
\usepackage{booktabs}
\usepackage{multirow}
\usepackage{subcaption}

%% Đảm bảo float không trọn ra ngoài section
\usepackage[section]{placeins}

%%=============================================================
%% Cách sử dụng trong main.tex:
%%
%%   \documentclass[12pt,a4paper]{report}
%%   \input{preamble_ai3}
%%   \begin{document}
%%     ...
%%     \chapter{Thiết kế và Hiện thực}
%%     ...
%%     \input{chapter4_4}   %% mục 4.4
%%     ...
%%     \chapter{Kết quả và Đánh giá}
%%     ...
%%     \input{chapter5_4}   %% mục 5.4
%%     ...
%%   \end{document}
%%
%% Compile: pdflatex -> pdflatex (chạy 2 lần để có label/ref đúng)
%%=============================================================

%%   \begin{document}
%%     ...
%%     \chapter{Thiết kế và Hiện thực}
%%     ...
%%     %%=============================================================
%% chapter4_4.tex
%% Phần 4.4 — AI #3: Phát hiện vật thể cho pick & place
%%
%% Đây là \input file — thêm vào preamble của main.tex:
%%   \usepackage{tikz}
%%   \usetikzlibrary{shapes.geometric,arrows.meta,positioning,
%%     fit,backgrounds,calc,decorations.pathreplacing}
%%   \usepackage{pgfplots}
%%   \pgfplotsset{compat=1.18}
%%   \usepackage{listings,xcolor,booktabs,multirow,subcaption}
%%=============================================================

\definecolor{cppblue}{RGB}{30,100,200}
\definecolor{pygreen}{RGB}{40,150,80}
\definecolor{warnred}{RGB}{200,60,60}
\definecolor{orng}{RGB}{220,130,30}

\lstset{
  basicstyle=\ttfamily\small,
  keywordstyle=\color{cppblue},
  commentstyle=\color{gray!70},
  stringstyle=\color{orng},
  breaklines=true,
  frame=single,
  rulecolor=\color{gray!40},
  framexleftmargin=4pt,
  xleftmargin=8pt,
  numbers=none,
}

% ================================================================
\section{AI \#3 --- Tự động hoá khâu khởi tạo trong hệ thống pick \& place}
\label{sec:ai3}
% ================================================================

% ----------------------------------------------------------------
\subsection{Bối cảnh: vấn đề của hệ thống visual servoing ban đầu}
\label{subsec:ai3-context}
% ----------------------------------------------------------------

Hệ thống pick \& place chạy trên cánh tay robot Denso VS-6577E 6 bậc tự do,
sử dụng thư viện ViSP (Visual Servoing Platform) để điều khiển bằng thị giác.
Camera USB được gắn trực tiếp lên đầu công tác theo cấu hình \textit{eye-in-hand} ---
camera di chuyển cùng robot và nhìn thẳng về phía vật thể cần gắp trong suốt
quá trình tiếp cận.

Vấn đề cốt lõi của hệ thống ban đầu nằm ở bước khởi tạo: sau khi robot về
home pose cố định, \textbf{người vận hành phải click chuột vào vị trí cylinder}
trên cửa sổ hiển thị camera thì blob tracker (\texttt{vpDot2}) mới bắt đầu
hoạt động được. Visual servo chỉ chạy sau thao tác thủ công đó. Ngoài ra,
nếu tracker mất target trong lúc servo --- điều hoàn toàn có thể xảy ra khi
cylinder bị che khuất hoặc di chuyển ra ngoài vùng sáng --- hệ thống reset
hoàn toàn và lại đòi hỏi click thủ công ở lần tiếp theo. Mỗi chu kỳ
pick \& place đều cần sự can thiệp của người vận hành, loại bỏ khả năng
vận hành tự động không giám sát.

\begin{figure}[htbp]
  \centering
  \begin{tikzpicture}[
    node distance=0.8cm and 0.6cm,
    bx/.style={draw, rounded corners=4pt, text width=3.0cm, align=center,
               minimum height=0.9cm, font=\small},
    robot/.style={bx, draw=cppblue!70, fill=cppblue!8},
    manual/.style={bx, draw=warnred!70, fill=warnred!8},
    aibox/.style={bx, draw=pygreen!70, fill=pygreen!8},
    arr/.style={->, >=Stealth, thick},
    rarr/.style={arr, draw=warnred!70},
    garr/.style={arr, draw=pygreen!70},
  ]

  %% --- CỘT TRÁI: hệ thống gốc ---
  \node[robot]                    (L1) {PREINIT\\{\scriptsize robot $\to$ home pose}};
  \node[manual, below=of L1]      (L2) {Click thủ công\\{\scriptsize chọn cylinder}};
  \node[robot,  below=of L2]      (L3) {JOINT\\{\scriptsize vpDot2 tracking}};
  \node[manual, below=of L3]      (L4) {Mất tracking\\{\scriptsize $\to$ reset toàn bộ}};

  \draw[arr]  (L1) -- (L2);
  \draw[arr]  (L2) -- (L3);
  \draw[rarr] (L3) -- (L4);
  \draw[rarr, dashed] (L4.west)
        to[out=190,in=170,looseness=1.3] (L1.west)
        node[midway, left, font=\scriptsize, text=warnred!70] {reset};

  \node[font=\bfseries\small, text=warnred, above=0.5cm of L1]
       {\textit{Trước cải tiến}};

  %% --- CỘT PHẢI: hệ thống mới ---
  \node[robot,  right=3.8cm of L1] (R1) {PREINIT\\{\scriptsize robot $\to$ home pose}};
  \node[aibox, below=of R1]        (R2) {AI detect\\{\scriptsize EfficientDet-Lite0}};
  \node[robot,  below=of R2]       (R3) {JOINT\\{\scriptsize AI tracking mỗi frame}};
  \node[aibox, below=of R3]        (R4) {AI re-detect\\{\scriptsize tự phục hồi}};

  \draw[arr]  (R1) -- (R2);
  \draw[garr] (R2) -- (R3);
  \draw[garr] (R3) -- (R4);
  \draw[garr, dashed] (R4.east)
        to[out=0,in=0,looseness=1.5] (R3.east)
        node[midway, right, font=\scriptsize, text=pygreen!70] {tiếp tục};

  %% fallback arrow từ AI detect
  \node[manual, right=1.2cm of R3] (FB)
       {{\scriptsize fallback:}\\click thủ công};
  \draw[rarr, dashed] (R2.east) -| (FB.north)
        node[near start, above, font=\scriptsize, text=warnred] {AI fail};

  \node[font=\bfseries\small, text=pygreen!80!black, above=0.5cm of R1]
       {\textit{Sau cải tiến}};

  %% Đường phân cách
  \draw[thick, dashed, gray!50]
      ($(L1.east)!0.5!(R1.west)+(0,0.8)$) -- ++(0,-6.5);

  \end{tikzpicture}
  \caption{So sánh luồng vận hành trước và sau khi tích hợp module AI.
           Bước yêu cầu thao tác thủ công (đỏ) được thay thế bằng AI tự động (xanh lá).
           Fallback về click thủ công vẫn được giữ lại cho trường hợp AI thất bại.}
  \label{fig:ai3-before-after}
\end{figure}

Hình~\ref{fig:ai3-before-after} minh hoạ sự thay đổi này. Module AI đảm nhiệm
toàn bộ việc phát hiện cylinder ở bước khởi tạo và tự phục hồi khi mất target
trong quá trình servo. Hệ thống vẫn giữ nguyên fallback về click thủ công để
đảm bảo tương thích ngược.


% ----------------------------------------------------------------
\subsection{Xây dựng và huấn luyện mô hình EfficientDet-Lite0}
\label{subsec:ai3-model}
% ----------------------------------------------------------------

\subsubsection*{Lựa chọn kiến trúc}

Yêu cầu đặt ra khá cụ thể: mô hình phải đủ nhỏ để triển khai được trên
Google Coral EdgeTPU, phải hỗ trợ INT8 quantization, và tương thích với
\texttt{tflite\_model\_maker} --- framework đã có sẵn trong môi trường phát triển.
\textbf{EfficientDet-Lite0} đáp ứng tất cả: đây là biến thể nhỏ nhất trong
họ EfficientDet được Google tối ưu cho thiết bị edge, input cố định
$320 \times 320$ pixel, backbone là EfficientNet-B0 đã được simplify.

\begin{table}[htbp]
  \centering
  \caption{Thông số mô hình EfficientDet-Lite0}
  \label{tab:ai3-model}
  \small
  \begin{tabular}{ll}
    \toprule
    \textbf{Thuộc tính}       & \textbf{Giá trị} \\
    \midrule
    Kiến trúc                 & EfficientDet-Lite0 \\
    Framework training        & TensorFlow~2.8.4 / \texttt{tflite\_model\_maker}~0.4.3 \\
    Input size                & $320 \times 320$~px, uint8 \\
    Số lớp đối tượng          & 2 (cylinder~--~lớp~0, cube~--~lớp~1) \\
    Số epochs / batch size    & 50 / 8 \\
    Transfer learning         & Frozen backbone, fine-tune detection head \\
    Target hardware           & Google Coral USB Accelerator (EdgeTPU) \\
    \bottomrule
  \end{tabular}
\end{table}

Mô hình được huấn luyện để phát hiện \textbf{hai lớp}: cylinder và cube nhựa PLA
in 3D. Việc đưa thêm cube vào tập huấn luyện là có chủ ý --- giúp mô hình học
cách phân biệt hai loại vật thể, từ đó giảm false positive khi cả hai cùng
xuất hiện trong cảnh. Ở phía tích hợp với ViSP, chỉ kết quả cylinder
(\texttt{class\_id~==~0}) được dùng; dự đoán cube bị lọc bỏ.

\subsubsection*{Tập dữ liệu}

Toàn bộ ảnh được thu trực tiếp từ camera eye-in-hand gắn trên robot Denso,
đảm bảo điều kiện thực tế sát với lúc triển khai. Ảnh bao gồm nhiều cấu hình:
góc chụp khác nhau, khoảng cách thay đổi, và số lượng vật thể từ 1 đến 5 trong
một cảnh.

\begin{figure}[htbp]
  \centering
  \begin{tikzpicture}
    \begin{axis}[
      xbar, bar width=0.42cm,
      xlabel={Số lượng ảnh (xấp xỉ)},
      ylabel={Danh mục},
      symbolic y coords={
        Background,
        2 cube + 3 cyl,
        2 cylinders,
        3 cylinders,
        1 cyl + 1 cube,
        1 cube,
        1 cylinder,
        Cylinder (gốc)
      },
      ytick=data,
      xmin=0, xmax=175,
      width=9.5cm, height=6.8cm,
      nodes near coords,
      nodes near coords align={horizontal},
      every node near coord/.style={font=\scriptsize},
      tick label style={font=\small},
      xlabel style={font=\small},
      ylabel style={font=\small},
      enlarge y limits=0.08,
    ]
    \addplot[fill=cppblue!55, draw=cppblue!80] coordinates {
      (38,  Background)
      (28,  2 cube + 3 cyl)
      (52,  2 cylinders)
      (43,  3 cylinders)
      (68,  1 cyl + 1 cube)
      (38,  1 cube)
      (108, 1 cylinder)
      (144, Cylinder (gốc))
    };
    \end{axis}
  \end{tikzpicture}
  \caption{Phân bố tập dữ liệu theo danh mục ảnh (~534 ảnh tổng,
           633 annotation cylinder, 170 annotation cube).
           Các cảnh đa vật thể giúp mô hình học phát hiện từng đối tượng
           riêng lẻ trong môi trường phức tạp.}
  \label{fig:ai3-dataset}
\end{figure}

Annotation được thực hiện bằng \texttt{labelme} theo format bounding box hình chữ
nhật. Workflow:
\begin{enumerate}
  \item Vẽ bounding box cho từng cylinder và cube bằng \texttt{labelme}
  \item Chuyển đổi labelme JSON $\to$ COCO JSON thống nhất
        (script \texttt{labelme\_to\_coco.py})
  \item Tập dữ liệu được split theo tỉ lệ $70/15/15$ (train/val/test),
        seed cố định để đảm bảo tính tái lập
\end{enumerate}

Phân chia ba chiều thay vì hai chiều đảm bảo tập test hoàn toàn độc lập ---
không được nhìn thấy trong suốt quá trình training hoặc điều chỉnh siêu tham số.

\subsubsection*{Pipeline huấn luyện}

\begin{figure}[htbp]
  \centering
  \begin{tikzpicture}[
    node distance=0.45cm and 0.15cm,
    st/.style={draw, rounded corners=3pt, fill=gray!10, draw=gray!55,
               text width=2.0cm, align=center, minimum height=0.85cm, font=\small},
    blue_st/.style={st, fill=cppblue!10, draw=cppblue!55},
    out_st/.style={st, fill=orng!20, draw=orng!70},
    arr/.style={->, >=Stealth, thick, gray!60},
    lbl/.style={font=\scriptsize, text=gray!80},
  ]

  \node[blue_st] (s1) {Annotation\\{\scriptsize labelme}};
  \node[blue_st, right=of s1] (s2) {COCO\\JSON};
  \node[blue_st, right=of s2] (s3) {Pascal\\VOC XML};
  \node[blue_st, right=of s3] (s4) {Training\\{\scriptsize 50 epochs}};
  \node[out_st, right=of s4] (s5) {SavedModel\\{\scriptsize Float32}};
  \node[out_st, below=0.55cm of s5] (s6) {INT8\\TFLite};
  \node[out_st, below=0.55cm of s6] (s7) {EdgeTPU\\TFLite};

  \draw[arr] (s1) -- (s2) node[lbl,above,midway] {to\_coco.py};
  \draw[arr] (s2) -- (s3) node[lbl,above,midway] {tflite\_mm};
  \draw[arr] (s3) -- (s4);
  \draw[arr] (s4) -- (s5) node[lbl,above,midway] {export};
  \draw[arr] (s5) -- (s6) node[lbl,right,midway] {PTQ};
  \draw[arr] (s6) -- (s7) node[lbl,right,midway] {edgetpu\_compiler};

  \end{tikzpicture}
  \caption{Pipeline huấn luyện và xuất mô hình. Từ ảnh thô đến ba file
           mô hình phục vụ các môi trường triển khai khác nhau.
           Bước chuyển đổi sang Pascal VOC XML là yêu cầu bắt buộc của
           \texttt{tflite\_model\_maker}.}
  \label{fig:ai3-pipeline}
\end{figure}

Transfer learning với \texttt{train\_whole\_model=False} --- chỉ fine-tune detection
head, giữ nguyên backbone đã pretrain --- phù hợp với tập dữ liệu tương đối nhỏ
($\sim$374 ảnh train) và giảm nguy cơ overfitting đáng kể so với
huấn luyện từ đầu.


% ----------------------------------------------------------------
\subsection{Quantization INT8 và biên dịch cho EdgeTPU}
\label{subsec:ai3-quant}
% ----------------------------------------------------------------

Mô hình sau training là Float32 --- không chạy được trực tiếp trên Coral EdgeTPU.
EdgeTPU chỉ thực thi mô hình INT8 quantized, trong đó trọng số và activation
được biểu diễn bằng số nguyên 8-bit. Lợi ích: kích thước giảm $\sim4\times$,
và phần cứng EdgeTPU đạt tốc độ inference cực nhanh ($\sim$4--15~ms/frame)
nhờ tận dụng phép tính số nguyên.

Post-Training Quantization (PTQ) sử dụng tập train làm \textit{representative
dataset} để calibrate scale và zero\_point cho từng layer --- đây là bước bắt
buộc để duy trì độ chính xác của INT8 model gần với Float32 baseline.

File \texttt{\_edgetpu.tflite} sau khi biên dịch chứa custom op
\texttt{edgetpu-custom-op} mà TFLite runtime thông thường
\textit{không thể thực thi trên CPU}, gây lỗi:

\begin{lstlisting}[language={}]
RuntimeError: Encountered unresolved custom op: edgetpu-custom-op
\end{lstlisting}

Do đó hệ thống \textbf{duy trì hai file riêng biệt} và lựa chọn tự động:

\begin{center}
\small
\begin{tabular}{lll}
  \toprule
  File & Runtime & Deployment \\
  \midrule
  \texttt{cylinder\_detector\_int8\_edgetpu.tflite} & pycoral   & Coral EdgeTPU \\
  \texttt{cylinder\_detector\_int8.tflite}          & TFLite/ai-edge-litert & CPU \\
  \bottomrule
\end{tabular}
\end{center}


% ----------------------------------------------------------------
\subsection{Kiến trúc tích hợp: C++ ViSP ↔ Python subprocess}
\label{subsec:ai3-arch}
% ----------------------------------------------------------------

\subsubsection*{Quyết định thiết kế}

Thay vì link trực tiếp thư viện TFLite C++ vào build system của ViSP
(đòi hỏi thay đổi lớn vào CMakeLists và phụ thuộc version build cụ thể),
hệ thống áp dụng kiến trúc \textbf{subprocess model}: C++ gọi Python script
như một tiến trình con qua \texttt{popen()}, giao tiếp một chiều qua stdout.
Cách này giúp module AI hoàn toàn tách biệt --- cập nhật mô hình không cần
recompile code robot.

\begin{figure}[htbp]
  \centering
  \begin{tikzpicture}[
    node distance=0.38cm and 0.5cm,
    comp/.style={draw, rounded corners=4pt, text width=3.1cm, align=center,
                 minimum height=0.82cm, font=\small},
    cpp/.style={comp, fill=cppblue!10, draw=cppblue!55},
    py/.style={comp,  fill=pygreen!10, draw=pygreen!60},
    tmp/.style={comp,  fill=orng!18, draw=orng!65, text width=2.5cm,
                minimum height=0.7cm, font=\scriptsize},
    arr/.style={->, >=Stealth, thick},
    ]

  %% ---- C++ side ----
  \node[cpp] (cap)   {Camera capture\\{\scriptsize vpImage{$<$}uchar{$>$}}};
  \node[cpp, below=of cap]  (save)  {Save frame\\{\scriptsize cv::imwrite()}};
  \node[cpp, below=of save] (pop)   {popen() subprocess\\{\scriptsize build cmd}};
  \node[cpp, below=of pop]  (sel)   {select() timeout 3s\\{\scriptsize POSIX wait}};
  \node[cpp, below=of sel]  (parse) {Parse stdout\\{\scriptsize sscanf SUCCESS}};
  \node[cpp, below=of parse](feat)  {vpFeatureBuilder\\{\scriptsize create feature}};

  \draw[arr](cap)--(save); \draw[arr](save)--(pop);
  \draw[arr](pop)--(sel);  \draw[arr](sel)--(parse); \draw[arr](parse)--(feat);

  \node[draw=cppblue!50, dashed, rounded corners=6pt,
        fit=(cap)(save)(pop)(sel)(parse)(feat),
        inner sep=7pt] (cbox) {};
  \node[above=2pt of cbox, font=\bfseries\small, text=cppblue]
       {C++ (ViSP main loop)};

  %% ---- Temp file ----
  \node[tmp, right=2.6cm of save] (jpg)
       {/tmp/visp\_ai\_frame.jpg};

  %% ---- Python side ----
  \node[py, right=2.6cm of pop]   (load) {Load model\\{\scriptsize fallback chain}};
  \node[py, below=of load]        (pre)  {Preprocess\\{\scriptsize resize 320$\times$320}};
  \node[py, below=of pre]         (inf)  {Inference\\{\scriptsize invoke()}};
  \node[py, below=of inf]         (post) {Post-process\\{\scriptsize dequant+filter}};
  \node[py, below=of post]        (out)  {Print stdout\\{\scriptsize SUCCESS u v conf}};

  \draw[arr](load)--(pre); \draw[arr](pre)--(inf);
  \draw[arr](inf)--(post); \draw[arr](post)--(out);

  \node[draw=pygreen!55, dashed, rounded corners=6pt,
        fit=(load)(pre)(inf)(post)(out),
        inner sep=7pt] (pbox) {};
  \node[above=2pt of pbox, font=\bfseries\small, text=pygreen!80!black]
       {Python subprocess};

  %% ---- Data flow arrows ----
  \draw[arr, orng!80]  (save.east) -- (jpg.west)  node[midway,above,font=\scriptsize]{write};
  \draw[arr, orng!80]  (jpg.east)  -- (load.west) node[midway,above,font=\scriptsize]{read};
  \draw[arr, cppblue!60](pop.east)  -- (load.west) node[midway,above,font=\scriptsize,yshift=2pt]{spawn};
  \draw[arr, pygreen!70, bend right=25](out.west)  to (parse.east)
        node[midway,below,font=\scriptsize]{stdout pipe};

  \end{tikzpicture}
  \caption{Kiến trúc tích hợp C++ ViSP $\leftrightarrow$ Python subprocess.
           C++ lưu frame hiện tại ra file tạm, spawn Python làm child process
           qua \texttt{popen()}, đọc kết quả từ stdout sau khi qua timeout guard.}
  \label{fig:ai3-arch}
\end{figure}

\subsubsection*{Protocol trao đổi dữ liệu}

Output của Python script là một dòng duy nhất trên stdout:

\begin{lstlisting}[language={}]
SUCCESS 272.86 190.74 0.8362 238.42 166.89 68.88 47.68
%        u      v      conf   x_min  y_min  width  height
\end{lstlisting}

hoặc \texttt{FAILURE~<lý~do>} nếu không phát hiện được. Tất cả debug log
đi ra \texttt{stderr} --- C++ chỉ đọc stdout --- tránh nhiễu kết quả parse.

Đáng chú ý: trong ViSP, \texttt{vpImagePoint} được định nghĩa là \texttt{(row,~col)}
tương ứng \texttt{(v,~u)}, không phải \texttt{(u,~v)} như thông thường:

\begin{lstlisting}[language=C++]
detected_center = vpImagePoint(v, u);  // row=v, col=u
\end{lstlisting}

Tham số \texttt{hint} cho phép Python chọn cylinder \textit{gần nhất với vị
trí đã biết} thay vì cylinder có confidence cao nhất --- cơ chế này tránh
hiện tượng servo ``nhảy'' sang cylinder khác khi scores thay đổi giữa các frame.


% ----------------------------------------------------------------
\subsection{State machine: PREINIT \texorpdfstring{$\to$}{→} INIT (AI) \texorpdfstring{$\to$}{→} JOINT \texorpdfstring{$\to$}{→} ...}
\label{subsec:ai3-sm}
% ----------------------------------------------------------------

Hệ thống vận hành theo vòng lặp chính liên tục với 6 trạng thái.
Mỗi iteration đọc một frame camera và thực thi logic trạng thái hiện tại.

\begin{figure}[htbp]
  \centering
  \begin{tikzpicture}[
    node distance=1.0cm and 2.0cm,
    st/.style={draw, rounded corners=4pt, fill=cppblue!10, draw=cppblue!55,
               text width=2.3cm, align=center, minimum height=1.0cm, font=\small},
    ai/.style={st, fill=pygreen!10, draw=pygreen!60},
    arr/.style={->, >=Stealth, thick},
    lbl/.style={font=\scriptsize, text=gray!70},
  ]

  %% Top row
  \node[st]  (pre)  {PREINIT\\{\scriptsize home pose}};
  \node[ai,  right=of pre]   (ini)  {INIT\\{\scriptsize AI detect}};
  \node[ai,  right=of ini]   (jnt)  {JOINT\\{\scriptsize AI tracking}};
  \node[st,  right=of jnt]   (app)  {APPROACH\\{\scriptsize chờ RC5}};

  %% Bottom row
  \node[st,  below=1.8cm of app]  (grp)  {GRIPPER\\{\scriptsize đóng gripper}};
  \node[st,  left=of grp]         (cls)  {CLASSIFIED\\{\scriptsize vận chuyển}};

  %% Arrows
  \draw[arr](pre)--(ini)  node[lbl,above,midway]{home\\reached};
  \draw[arr](ini)--(jnt)  node[lbl,above,midway]{AI init\\ok};
  \draw[arr](jnt)--(app)  node[lbl,above,midway]{$|e|<5{\times}10^{-3}$};
  \draw[arr](app)--(grp)  node[lbl,right,midway]{RC5\\signal};
  \draw[arr](grp)--(cls)  node[lbl,below,midway]{grasped};
  \draw[arr](cls.west) to[out=180,in=270,looseness=0.8] (pre.south)
        node[lbl,left,pos=0.5]{cycle\\restart};

  %% AI self-recovery
  \draw[arr, pygreen!70, dashed]
        (jnt.north) to[out=100,in=80,looseness=2.5] (jnt.north)
        node[above=0.7cm of jnt, font=\scriptsize, text=pygreen!80!black]
        {AI re-detect};

  %% Fallback from INIT
  \draw[arr, warnred!70, dashed]
        (ini.south) -- ++(0,-0.7)
        node[below, font=\scriptsize, text=warnred]
        {fallback: click thủ công};

  \end{tikzpicture}
  \caption{State machine của hệ thống pick \& place sau cải tiến.
           Trạng thái tô xanh lá là nơi AI được gọi.
           Vòng lặp tự phục hồi ở JOINT (mũi tên đứt nét) cho phép hệ thống
           tiếp tục mà không cần reset khi mất detection.}
  \label{fig:ai3-sm}
\end{figure}

\paragraph{PREINIT.}
Robot gửi lệnh joint position về home pose
$[0^\circ, 0^\circ, 90^\circ, 0^\circ, 90^\circ, 0^\circ]$.
Servo task được cấu hình với interaction matrix Desired, pseudo-inverse,
gain $\lambda = 0.4$.

\paragraph{INIT.}
Chờ robot đạt home pose (sai số từng joint $< 0.01$) và gripper mở xong.
Khi cả hai điều kiện thỏa, \texttt{detectCylinderWithAI()} được gọi
không có hint --- Python chọn cylinder có confidence cao nhất.
Kết quả trực tiếp tạo visual feature \texttt{vpFeaturePoint}
và add vào servo task. Nếu AI thất bại, \texttt{ai\_hint}
giữ giá trị mặc định (fallback thủ công).

\paragraph{JOINT.}
Khác với cách tiếp cận truyền thống dùng blob tracker mỗi frame,
hệ thống hiện tại \textbf{gọi AI trực tiếp ở mỗi iteration}
với \texttt{hint = \&ai\_hint}. Control law:
\begin{equation}
  \mathbf{v} = -\lambda\,\mathbf{L}_{s^*}^{+}\,\mathbf{e},
  \qquad \mathbf{q}_{k+1} = \mathbf{q}_k + \mathbf{v}\,\Delta t
  \label{eq:ibvs}
\end{equation}
với $\Delta t = 0.5\,$s. Điều kiện hội tụ: $|e_x| < 5{\times}10^{-3}$
và $|e_y| < 5{\times}10^{-3}$ (normalized image coordinates).
Khi đạt, gửi \texttt{"OKE\textbackslash r"} qua UART đến RC5 controller.


% ----------------------------------------------------------------
\subsection{Các vấn đề kỹ thuật và giải pháp}
\label{subsec:ai3-issues}
% ----------------------------------------------------------------

Một số vấn đề nảy sinh trong quá trình tích hợp mà khó thấy trước khi thực
sự chạy hệ thống.

\subsubsection*{P1: Xác định output tensor sai lệch}

\texttt{tflite\_model\_maker} 0.4.x đặt tên tensor theo thứ tự riêng của nó:

\begin{center}
\small
\begin{tabular}{ccc}
  \toprule
  Tensor & Shape & Nội dung \\
  \midrule
  \texttt{:0} & \texttt{[1]}     & num\_detections \\
  \texttt{:1} & \texttt{[1,\,N]} & scores \quad \textbf{← cần lấy cái này} \\
  \texttt{:2} & \texttt{[1,\,N]} & class IDs \quad \textbf{← không phải cái này} \\
  \texttt{:3} & \texttt{[1,\,N,\,4]} & boxes \\
  \bottomrule
\end{tabular}
\end{center}

Thứ tự này không khớp với API của pycoral, vốn giả định positional order
\texttt{boxes/classes/scores/count = 0/1/2/3}. Dùng API đó sẽ đọc
num\_detections (raw uint8 = 255) như scores và crash với \texttt{IndexError}.
Giải pháp: xác định tensor \textbf{theo shape động} --- boxes là
\texttt{[1,N,4]}, scores là tensor \texttt{[1,N]} có name suffix nhỏ nhất
(\texttt{:1} chứ không phải \texttt{:2}).

\subsubsection*{P2: Python 3.10 không có \texttt{tflite-runtime}}

Google không publish wheel \texttt{tflite-runtime} cho Python 3.10+.
Giải pháp là xây dựng \textbf{fallback chain 4 tầng}:

\begin{figure}[htbp]
  \centering
  \begin{tikzpicture}[
    node distance=0.4cm,
    fb/.style={draw, rounded corners=3pt, text width=7.5cm, align=left,
               minimum height=0.85cm, font=\small, inner sep=7pt},
    t1/.style={fb, fill=purple!12, draw=purple!55},
    t2/.style={fb, fill=teal!10,   draw=teal!50},
    t3/.style={fb, fill=cppblue!10, draw=cppblue!50},
    t4/.style={fb, fill=gray!12,   draw=gray!50},
    fa/.style={fb, fill=warnred!10, draw=warnred!50},
    arr/.style={->, >=Stealth, thick, dashed, gray!55},
    lbl/.style={font=\scriptsize, text=warnred!80},
  ]
  \node[t1] (a) {\textbf{1.}~\texttt{pycoral.utils.edgetpu}\quad{\scriptsize Coral EdgeTPU, Python 3.9}};
  \node[t2, below=of a] (b) {\textbf{2.}~\texttt{ai\_edge\_litert.interpreter}\quad{\scriptsize Google's new pkg, Python 3.8--3.12}};
  \node[t3, below=of b] (c) {\textbf{3.}~\texttt{tflite\_runtime.interpreter}\quad{\scriptsize Legacy, Python 3.9}};
  \node[t4, below=of c] (d) {\textbf{4.}~\texttt{tensorflow.lite.Interpreter}\quad{\scriptsize Heavy fallback}};
  \node[fa, below=of d] (e) {\textbf{FAILURE}:~\texttt{no\_tflite\_backend\_available}};

  \draw[arr](a)--(b) node[lbl,right,midway]{ImportError / Coral not connected};
  \draw[arr](b)--(c) node[lbl,right,midway]{ImportError};
  \draw[arr](c)--(d) node[lbl,right,midway]{ImportError};
  \draw[arr](d)--(e) node[lbl,right,midway]{ImportError};
  \end{tikzpicture}
  \caption{Python backend fallback chain trong \texttt{detect\_cylinder.py}.
           Script tự động thử từng backend theo thứ tự ưu tiên,
           tương thích với cả môi trường có Coral lẫn CPU-only.}
  \label{fig:ai3-fallback}
\end{figure}

\texttt{ai-edge-litert} là gói chính thức của Google thay thế \texttt{tflite-runtime},
hỗ trợ Python 3.8--3.12 với API giống hệt nhau --- là giải pháp cho vấn đề Python 3.10.

\subsubsection*{P3: JSON comment trong config}

\texttt{config.json} dùng comment \texttt{//} để ghi chú, nhưng cả Python
\texttt{json.load()} và \texttt{nlohmann::json} đều từ chối mặc định.
Phía C++: \texttt{nlohmann::json::parse(file, nullptr, true, true)} với
\texttt{ignore\_comments=true} (có từ nlohmann/json~3.9.0).
Phía Python: strip comment lines thủ công trước khi \texttt{json.loads()}.


% ----------------------------------------------------------------
\subsection{Thông số cấu hình visual servo}
\label{subsec:ai3-servo}
% ----------------------------------------------------------------

Hệ thống sử dụng \textbf{Image-Based Visual Servoing (IBVS)} với feature là
tọa độ 2D của cylinder trong không gian ảnh normalized $\mathbf{s}=(x,y)$,
desired position $\mathbf{s}^*=(0,0,1)$ --- cylinder phải nằm đúng tâm ảnh.

\begin{table}[htbp]
  \centering
  \caption{Thông số điều khiển visual servo}
  \label{tab:ai3-servo}
  \small
  \begin{tabular}{ll}
    \toprule
    \textbf{Tham số}            & \textbf{Giá trị} \\
    \midrule
    Loại điều khiển             & Eye-in-hand, \texttt{EYEINHAND\_L\_cVe\_eJe} \\
    Interaction matrix          & Desired ($\mathbf{L}_{s^*}$), pseudo-inverse \\
    Gain $\lambda$              & 0.4 \\
    Bước thời gian $\Delta t$   & 0.5\,s \\
    Ngưỡng hội tụ               & $|e_x| < 5{\times}10^{-3}$ và $|e_y| < 5{\times}10^{-3}$ \\
    Desired position            & $(x^*, y^*) = (0,\,0)$ --- tâm ảnh, $Z = 1$ \\
    \bottomrule
  \end{tabular}
\end{table}

Dùng $\mathbf{L}_{s^*}$ (tại vị trí mong muốn) thay vì $\mathbf{L}_{s_{cur}}$
giúp tránh singularity khi feature tiến gần desired point và giữ ổn định
theo ổn định hơn trong giai đoạn hội tụ. Velocity được saturate trước khi
gửi xuống robot để đảm bảo không vượt giới hạn vật lý từng khớp.
   %% mục 4.4
%%     ...
%%     \chapter{Kết quả và Đánh giá}
%%     ...
%%     %%=============================================================
%% chapter5_4.tex
%% Phần 5.4 — Kết quả đánh giá: AI #3 Phát hiện vật thể
%%
%% \input file — cần thêm vào preamble:
%%   \usepackage{tikz}
%%   \usetikzlibrary{arrows.meta,positioning,patterns}
%%   \usepackage{pgfplots}
%%   \pgfplotsset{compat=1.18}
%%   \usepackage{booktabs,multirow,subcaption,xcolor}
%%=============================================================

% ================================================================
\section{Kết quả và Đánh giá: AI \#3 --- Phát hiện vật thể}
\label{sec:ai3-results}
% ================================================================

% ----------------------------------------------------------------
\subsection{Kết quả detection model: COCO metrics}
\label{subsec:ai3-coco}
% ----------------------------------------------------------------

Mô hình được đánh giá trên cả validation set và test set độc lập
theo chuẩn \textbf{COCO evaluation protocol} --- tính trung bình mAP qua
các ngưỡng IoU từ 0.50 đến 0.95 với bước 0.05.

\begin{figure}[htbp]
  \centering
  \begin{tikzpicture}
    \begin{axis}[
      ybar, bar width=0.45cm,
      width=10cm, height=7cm,
      symbolic x coords={mAP (0.50:0.95), AP50, AP75, AR},
      xtick=data,
      x tick label style={font=\small, align=center},
      ymin=0, ymax=1.15,
      ylabel={Score},
      ylabel style={font=\small},
      legend style={at={(0.5,1.03)}, anchor=south, legend columns=2, font=\small},
      legend cell align=left,
      nodes near coords,
      every node near coord/.style={font=\scriptsize, /pgf/number format/fixed,
                                    /pgf/number format/precision=3},
      enlarge x limits=0.25,
      grid=major, grid style={dashed, gray!30},
      axis line style={gray!60},
      tick style={gray!60},
    ]
    \addplot[fill=cppblue!70, draw=cppblue!90] coordinates {
      (mAP (0.50:0.95), 0.791)
      (AP50, 0.990)
      (AP75, 0.990)
      (AR,   0.869)
    };
    \addplot[fill=orng!80, draw=orng!90] coordinates {
      (mAP (0.50:0.95), 0.799)
      (AP50, 1.000)
      (AP75, 0.969)
      (AR,   0.856)
    };
    \legend{Validation, Test}
    \end{axis}
  \end{tikzpicture}
  \caption{COCO metrics trên validation set và test set. AP50 trên test set
           đạt 1.000 cho thấy 100\% cylinder được phát hiện đúng ở ngưỡng IoU
           lỏng lẻ (50\%). mAP$\approx$0.80 phản ánh vẫn còn đây đó những trường hợp
           định vị bounding box chưa được chính xác ở ngưỡng IoU cao.}
  \label{fig:ai3-coco}
\end{figure}

Bảng~\ref{tab:ai3-metrics} tổng hợp số liệu chi tiết:

\begin{table}[htbp]
  \centering
  \caption{COCO metrics trên validation và test set}
  \label{tab:ai3-metrics}
  \small
  \begin{tabular}{lcc}
    \toprule
    \textbf{Metric}                    & \textbf{Validation} & \textbf{Test} \\
    \midrule
    mAP (IoU 0.50:0.95)               & 0.791               & \textbf{0.799} \\
    AP50 (IoU $\geq$ 0.50)            & 0.990               & \textbf{1.000} \\
    AP75 (IoU $\geq$ 0.75)            & 0.990               & 0.969 \\
    AR (maxDets=100)                   & 0.869               & 0.856 \\
    \midrule
    Precision (IoU=0.50, conf=0.50)    & ---                 & \textbf{1.000} \\
    Recall    (IoU=0.50, conf=0.50)    & ---                 & 0.990 \\
    F1-Score                           & ---                 & \textbf{0.995} \\
    TP / FP / FN                       & ---                 & 95 / 0 / 1 \\
    \bottomrule
  \end{tabular}
\end{table}

\begin{figure}[htbp]
  \centering
  \begin{tikzpicture}
    \begin{axis}[
      ybar, bar width=0.55cm,
      width=8.5cm, height=6cm,
      symbolic x coords={Precision, Recall, F1-Score},
      xtick=data,
      x tick label style={font=\small},
      ymin=0.96, ymax=1.04,
      ylabel={Score},
      ylabel style={font=\small},
      nodes near coords,
      every node near coord/.style={font=\small, /pgf/number format/fixed,
                                    /pgf/number format/precision=3},
      enlarge x limits=0.3,
      grid=major, grid style={dashed, gray!30},
      axis line style={gray!60},
    ]
    \addplot[fill=pygreen!65, draw=pygreen!80] coordinates {
      (Precision, 1.000)
      (Recall,    0.990)
      (F1-Score,  0.995)
    };
    \end{axis}
  \end{tikzpicture}
  \caption{Precision, Recall và F1-Score trên test set (PASCAL VOC protocol,
           IoU$\geq$0.50, confidence$\geq$0.50, chỉ lớp cylinder).
           Precision = 1.000 và FP = 0 chứng tỏ không có false positive
           nào ở ngưỡng 0.50. FN = 1 (một cylinder bị bỏ sót) giải thích
           Recall = 0.990.}
  \label{fig:ai3-prf}
\end{figure}

Giá trị AP50 = 1.000 trên test set à ngưỡng IoU lỏng lẻ (50\%) phản ánh
mô hình phát hiện được toàn bộ cylinder. mAP = 0.799 (trung bình
qua IoU 0.50--0.95) thấp hơn một chút do mô hình EfficientDet-Lite0
đơn giản vẫn còn giới hạn trong việc định vị bounding box chính xác
ở các ngưỡng IoU cao ($\geq$0.75) --- điều phổ biến với các kiến trúc
nhe hơn.


% ----------------------------------------------------------------
\subsection{Phân tích confidence score: true positive vs false positive}
\label{subsec:ai3-conf}
% ----------------------------------------------------------------

Một trong những điều quan trọng khi triển khai thực tế là khoảng cách
phân tách giữa true positive và false positive --- điều này quyết định
ngưỡng confidence nào thực tế dùng được mà không sai kết quả.

Kiểm tra trên 218 ảnh cylinder (true positives) và 10 ảnh cube
(false positives):

\begin{figure}[htbp]
  \centering
  \begin{tikzpicture}
    \begin{axis}[
      width=11cm, height=5.5cm,
      xlabel={Confidence score},
      ylabel={},
      xmin=0, xmax=1.05,
      ymin=0, ymax=3,
      xtick={0,0.1,...,1.0},
      ytick={0.7, 2.0},
      yticklabels={\small False Positive\\(cube, $n$=10), \small True Positive\\(cylinder, $n$=218)},
      y tick label style={align=right, font=\small},
      x tick label style={font=\small},
      xlabel style={font=\small},
      grid=both, grid style={dashed, gray!25},
      axis line style={gray!60},
    ]

    %% --- True Positive band (y=2) ---
    %% Most scores: 0.85-0.98 (majority)
    \addplot[only marks, mark=|, mark size=5pt, mark options={draw=pygreen!70, line width=1.2pt},
             opacity=0.5]
    coordinates {
      (0.29,2)(0.61,2)(0.74,2)
      (0.82,2)(0.83,2)(0.84,2)(0.85,2)(0.86,2)(0.87,2)(0.88,2)
      (0.89,2)(0.90,2)(0.91,2)(0.92,2)(0.93,2)(0.94,2)(0.94,2)
      (0.95,2)(0.95,2)(0.96,2)(0.96,2)(0.97,2)(0.97,2)(0.98,2)
    };

    %% --- False Positive band (y=0.7) ---
    \addplot[only marks, mark=|, mark size=5pt, mark options={draw=warnred!70, line width=1.2pt},
             opacity=0.6]
    coordinates {
      (0.08,0.7)(0.13,0.7)(0.17,0.7)(0.21,0.7)
      (0.25,0.7)(0.28,0.7)(0.30,0.7)(0.33,0.7)(0.35,0.7)(0.36,0.7)
    };

    %% --- Threshold line at 0.50 ---
    \addplot[thick, dashed, draw=gray!80, domain=0:3]
    coordinates {(0.50, 0) (0.50, 3)};
    \node[font=\scriptsize, text=gray!80, rotate=90]
         at (axis cs: 0.52, 1.5) {threshold = 0.50};

    %% --- Gap annotation ---
    \draw[->, >=Stealth, thick, gray!70]
         (axis cs:0.38, 2.5) -- (axis cs:0.48, 2.5);
    \draw[->, >=Stealth, thick, gray!70]
         (axis cs:0.34, 2.5) -- (axis cs:0.27, 2.5);
    \node[font=\scriptsize, text=gray!80] at (axis cs:0.37, 2.7) {gap = 0.34};

    %% --- Median annotation for TP ---
    \addplot[only marks, mark=diamond*, mark size=4pt,
             mark options={draw=pygreen!80!black, fill=pygreen!60}]
    coordinates {(0.94, 2)};
    \node[font=\scriptsize, text=pygreen!80!black, above]
         at (axis cs:0.94, 2) {median 0.94};

    %% --- Max FP annotation ---
    \addplot[only marks, mark=diamond*, mark size=4pt,
             mark options={draw=warnred!80, fill=warnred!50}]
    coordinates {(0.36, 0.7)};
    \node[font=\scriptsize, text=warnred, below]
         at (axis cs:0.36, 0.7) {max 0.36};

    \end{axis}
  \end{tikzpicture}
  \caption{Phân tách confidence score giữa true positive (cylinder, xanh lá)
           và false positive (cube, đỏ). Mỗi vạch đứng đại diện một detection.
           Ngưỡng 0.50 nằm chính giữa khoảng gap 0.34,
           loại bỏ toàn bộ false positive mà vẫn chấp nhận 99.5\% true positive.}
  \label{fig:ai3-conf}
\end{figure}

Khoảng cách 0.34 giữa điểm false positive cao nhất (0.36) và 99.1\%
true positive ($\geq$0.70) xác nhận mô hình đã học được đặc trưng
phân biệt mạnh giữa cylinder và cube. Ngưỡng hiện tại 0.50 nằm an toàn
giữa hai vùng phân bố này.

Ba ảnh outlier có score thấp (0.29, 0.61, 0.74) tương ứng với cylinder
bị nhìn từ góc bất thường hoặc bị che khuất một phần. Ỡ runtime,
các trường hợp này fallback về click thủ công --- hành vi an toàn.

\begin{table}[htbp]
  \centering
  \caption{Phân tích confidence score trên 218 ảnh cylinder + 10 ảnh cube}
  \label{tab:ai3-conf}
  \small
  \begin{tabular}{lcc}
    \toprule
    & \textbf{True Positive (cylinder)} & \textbf{False Positive (cube)} \\
    \midrule
    Score range      & 0.29 -- 0.98          & 0.08 -- 0.36 \\
    Median           & \textbf{0.94}         & --- \\
    $\geq$ 0.50      & 99.5\% (217/218)      & \textbf{0/10 (0\%)} \\
    $\geq$ 0.70      & 99.1\% (216/218)      & 0/10 (0\%) \\
    $\geq$ 0.80      & 95.9\% (209/218)      & 0/10 (0\%) \\
    \midrule
    Gap ph\^an t\'ach & \multicolumn{2}{c}{$0.70 - 0.36 = 0.34$} \\
    \bottomrule
  \end{tabular}
\end{table}


% ----------------------------------------------------------------
\subsection{Kết quả tích hợp hệ thống}
\label{subsec:ai3-integration}
% ----------------------------------------------------------------

\begin{table}[htbp]
  \centering
  \caption{So sánh khả năng trước và sau khi tích hợp AI}
  \label{tab:ai3-comparison}
  \small
  \begin{tabular}{lll}
    \toprule
    \textbf{Tiêu chí}                    & \textbf{Hệ thống gốc} & \textbf{Sau cải tiến} \\
    \midrule
    Khởi tạo mỗi chu kỳ               & Click thủ công (bắt buộc) & Tự động (AI) \\
    Tỷ lệ init thành công             & 100\% (do manual)  & $\sim$99.5\% (AI) \\
    Fallback khi AI thất bại           & Không có              & Click thủ công \\
    Phục hồi khi mất tracking          & Reset toàn bộ        & AI re-detect tự động \\
    Chạy không giám sát                & Không thể              & Có (trường hợp bình thường) \\
    \bottomrule
  \end{tabular}
\end{table}

Thiết kế hiện tại gọi AI trực tiếp \textbf{mỗi iteration} của servo
loop thay vì dùng blob tracker. Điều này có một hàm ý đáng chú ý:
\textit{không có tracking failure} theo nghĩa truyền thống --- mỗi frame
là một lần detection độc lập. Tính ổn định của servo phụ thuộc vào
sự nhất quán của AI detection qua các frame liên tiếp. Cơ chế \texttt{--hint}
(truyền vị trí frame trước) giúp được phần nào khi có nhiều cylinder
trong cảnh.


% ----------------------------------------------------------------
\subsection{Phân tích latency: subprocess CPU vs EdgeTPU}
\label{subsec:ai3-latency}
% ----------------------------------------------------------------

Mỗi lần gọi \texttt{detectCylinderWithAI()} bao gồm nhiều thành phần
latency khác nhau. Điều ít người để ý là: phần overhead lớn nhất
không phải là inference mà là \textbf{Python startup time} --- chi phí
mở một process mới mỗi lần gọi \texttt{popen()}.

\begin{figure}[htbp]
  \centering
  \begin{tikzpicture}
    \begin{axis}[
      xbar stacked,
      bar width=0.55cm,
      width=11cm, height=6.5cm,
      symbolic y coords={
        {EdgeTPU\\(persistent, tương lai)},
        {CPU\\(persistent, tương lai)},
        {EdgeTPU\\(subprocess hiện tại)},
        {CPU\\(subprocess hiện tại)}
      },
      ytick=data,
      y tick label style={align=right, font=\small},
      xlabel={Thời gian (ms)},
      xlabel style={font=\small},
      xmin=0, xmax=780,
      legend style={at={(0.5,-0.18)}, anchor=north, legend columns=4,
                    font=\scriptsize},
      legend cell align=left,
      grid=major, grid style={dashed, gray!25},
      axis line style={gray!60},
      enlarge y limits=0.18,
      nodes near coords, nodes near coords align={horizontal},
      point meta=explicit symbolic,
      every node near coord/.style={font=\scriptsize},
    ]
    %%  stacks: I/O | Inference | Model load | Python startup
    %% CPU subprocess: 10 + 150 + 100 + 350 = 610
    \addplot[fill=gray!40, draw=gray!60] coordinates {
      (10, {CPU\\(subprocess hiện tại)})
      (10, {EdgeTPU\\(subprocess hiện tại)})
      (10, {CPU\\(persistent, tương lai)})
      (5,  {EdgeTPU\\(persistent, tương lai)})
    };  %% I/O
    \addplot[fill=warnred!55, draw=warnred!70] coordinates {
      (150, {CPU\\(subprocess hiện tại)})
      (10,  {EdgeTPU\\(subprocess hiện tại)})
      (150, {CPU\\(persistent, tương lai)})
      (10,  {EdgeTPU\\(persistent, tương lai)})
    };  %% Inference
    \addplot[fill=orng!65, draw=orng!80] coordinates {
      (100, {CPU\\(subprocess hiện tại)})
      (50,  {EdgeTPU\\(subprocess hiện tại)})
      (0,   {CPU\\(persistent, tương lai)})
      (0,   {EdgeTPU\\(persistent, tương lai)})
    };  %% Model load
    \addplot[fill=cppblue!55, draw=cppblue!70] coordinates {
      (350, {CPU\\(subprocess hiện tại)})
      (350, {EdgeTPU\\(subprocess hiện tại)})
      (0,   {CPU\\(persistent, tương lai)})
      (0,   {EdgeTPU\\(persistent, tương lai)})
    };  %% Python startup
    \legend{I/O (JPEG), Inference, Model load, Python startup}
    \end{axis}
  \end{tikzpicture}
  \caption{Phân tích latency mỗi lần gọi AI detection.
           Python startup time ($\sim$350~ms) chiếm phần lẫn tổng overhead
           trong kiến trúc subprocess hiện tại --- lớn hơn nhiều so với
           thời gian inference thực tế. Kiến trúc persistent process (tương lai)
           loại bỏ hoàn toàn chi phí này.}
  \label{fig:ai3-latency}
\end{figure}

\begin{table}[htbp]
  \centering
  \caption{So sánh tổng latency giữa các cấu hình triển khai}
  \label{tab:ai3-latency}
  \small
  \begin{tabular}{lcc}
    \toprule
    \textbf{Cấu hình}                    & \textbf{Inference} & \textbf{Tổng / lần gọi} \\
    \midrule
    CPU subprocess (hiện tại)            & $\sim$100--200~ms  & $\sim$600--700~ms \\
    EdgeTPU subprocess (hiện tại)       & $\sim$4--15~ms     & $\sim$420--430~ms \\
    CPU persistent (tương lai)           & $\sim$100--200~ms  & $\sim$160--210~ms \\
    EdgeTPU persistent (tương lai)      & $\sim$4--15~ms     & $\sim$15--20~ms  \\
    \bottomrule
  \end{tabular}
\end{table}

Kết quả trên cho thấy rõ điểm nghẻn hiện tại:
EdgeTPU inference chỉ mất 4--15~ms nhưng toàn bộ lá gọi
vẫn tốn 420~ms do phải khởi động Python mỗi lần. Với
kiến trúc \textbf{persistent Python process} --- một process
duy nhất khởi động một lần, load mô hình vào bộ nhớ, sau
đó tiếp nhận request từ C++ qua stdin/stdout pipe --- latency
mỗi lần gọi giảm xuống còn $\sim$15--20~ms với EdgeTPU,
mở ra khả năng chạy AI detection \textit{mỗi frame} camera
thay vì chỉ khi cần.

% Chú thích về nguồn số liệu
\smallskip
\noindent\textit{\small
Lưu ý: COCO metrics trong Bảng~\ref{tab:ai3-metrics} và Hình~\ref{fig:ai3-coco}
lấy từ kết quả đánh giá hai lớp (cylinder + cube) trên repository huấn luyện.
Phân tích confidence score trong Bảng~\ref{tab:ai3-conf} lấy từ kiểm tra
độc lập trên 218 ảnh cylinder và 10 ảnh cube (CPU inference via ai-edge-litert).
}
   %% mục 5.4
%%     ...
%%   \end{document}
%%
%% Compile: pdflatex -> pdflatex (chạy 2 lần để có label/ref đúng)
%%=============================================================

%% Hoặc copy thẳng các \usepackage cần thiết vào preamble của main.
%%=============================================================

%% Tiếng Việt
\usepackage[T5]{fontenc}
\usepackage[utf8]{inputenc}
\usepackage[vietnamese]{babel}

%% TikZ và các thư viện vẻ hình
\usepackage{tikz}
\usetikzlibrary{
  shapes.geometric,
  arrows.meta,
  positioning,
  fit,
  backgrounds,
  calc,
  decorations.pathreplacing,
  patterns,
}

%% PGFPlots cho biểu đồ
\usepackage{pgfplots}
\pgfplotsset{compat=1.18}
\usepackage{pgfplotstable}

%% Code listings
\usepackage{listings}
\usepackage{xcolor}

%% Màu dùng chung (cần định nghĩa trước khi \input chapter)
\definecolor{cppblue}{RGB}{30,100,200}
\definecolor{pygreen}{RGB}{40,150,80}
\definecolor{warnred}{RGB}{200,60,60}
\definecolor{orng}{RGB}{220,130,30}

%% Cài đặt listings mặc định
\lstset{
  basicstyle=\ttfamily\small,
  keywordstyle=\color{cppblue},
  commentstyle=\color{gray!70},
  stringstyle=\color{orng},
  breaklines=true,
  frame=single,
  rulecolor=\color{gray!40},
  framexleftmargin=4pt,
  xleftmargin=8pt,
  numbers=none,
}

%% Hình ảnh
\usepackage{graphicx}

%% Bảng đẹp hơn
\usepackage{booktabs}
\usepackage{multirow}
\usepackage{subcaption}

%% Đảm bảo float không trọn ra ngoài section
\usepackage[section]{placeins}

%%=============================================================
%% Cách sử dụng trong main.tex:
%%
%%   \documentclass[12pt,a4paper]{report}
%%   %%=============================================================
%% preamble_ai3.tex
%% Preamble dùng chung cho chapter4_4.tex và chapter5_4.tex
%%
%% Đưa vào main.tex bằng:  %%=============================================================
%% preamble_ai3.tex
%% Preamble dùng chung cho chapter4_4.tex và chapter5_4.tex
%%
%% Đưa vào main.tex bằng:  \input{preamble_ai3}
%% Hoặc copy thẳng các \usepackage cần thiết vào preamble của main.
%%=============================================================

%% Tiếng Việt
\usepackage[T5]{fontenc}
\usepackage[utf8]{inputenc}
\usepackage[vietnamese]{babel}

%% TikZ và các thư viện vẻ hình
\usepackage{tikz}
\usetikzlibrary{
  shapes.geometric,
  arrows.meta,
  positioning,
  fit,
  backgrounds,
  calc,
  decorations.pathreplacing,
  patterns,
}

%% PGFPlots cho biểu đồ
\usepackage{pgfplots}
\pgfplotsset{compat=1.18}
\usepackage{pgfplotstable}

%% Code listings
\usepackage{listings}
\usepackage{xcolor}

%% Màu dùng chung (cần định nghĩa trước khi \input chapter)
\definecolor{cppblue}{RGB}{30,100,200}
\definecolor{pygreen}{RGB}{40,150,80}
\definecolor{warnred}{RGB}{200,60,60}
\definecolor{orng}{RGB}{220,130,30}

%% Cài đặt listings mặc định
\lstset{
  basicstyle=\ttfamily\small,
  keywordstyle=\color{cppblue},
  commentstyle=\color{gray!70},
  stringstyle=\color{orng},
  breaklines=true,
  frame=single,
  rulecolor=\color{gray!40},
  framexleftmargin=4pt,
  xleftmargin=8pt,
  numbers=none,
}

%% Hình ảnh
\usepackage{graphicx}

%% Bảng đẹp hơn
\usepackage{booktabs}
\usepackage{multirow}
\usepackage{subcaption}

%% Đảm bảo float không trọn ra ngoài section
\usepackage[section]{placeins}

%%=============================================================
%% Cách sử dụng trong main.tex:
%%
%%   \documentclass[12pt,a4paper]{report}
%%   \input{preamble_ai3}
%%   \begin{document}
%%     ...
%%     \chapter{Thiết kế và Hiện thực}
%%     ...
%%     \input{chapter4_4}   %% mục 4.4
%%     ...
%%     \chapter{Kết quả và Đánh giá}
%%     ...
%%     \input{chapter5_4}   %% mục 5.4
%%     ...
%%   \end{document}
%%
%% Compile: pdflatex -> pdflatex (chạy 2 lần để có label/ref đúng)
%%=============================================================

%% Hoặc copy thẳng các \usepackage cần thiết vào preamble của main.
%%=============================================================

%% Tiếng Việt
\usepackage[T5]{fontenc}
\usepackage[utf8]{inputenc}
\usepackage[vietnamese]{babel}

%% TikZ và các thư viện vẻ hình
\usepackage{tikz}
\usetikzlibrary{
  shapes.geometric,
  arrows.meta,
  positioning,
  fit,
  backgrounds,
  calc,
  decorations.pathreplacing,
  patterns,
}

%% PGFPlots cho biểu đồ
\usepackage{pgfplots}
\pgfplotsset{compat=1.18}
\usepackage{pgfplotstable}

%% Code listings
\usepackage{listings}
\usepackage{xcolor}

%% Màu dùng chung (cần định nghĩa trước khi \input chapter)
\definecolor{cppblue}{RGB}{30,100,200}
\definecolor{pygreen}{RGB}{40,150,80}
\definecolor{warnred}{RGB}{200,60,60}
\definecolor{orng}{RGB}{220,130,30}

%% Cài đặt listings mặc định
\lstset{
  basicstyle=\ttfamily\small,
  keywordstyle=\color{cppblue},
  commentstyle=\color{gray!70},
  stringstyle=\color{orng},
  breaklines=true,
  frame=single,
  rulecolor=\color{gray!40},
  framexleftmargin=4pt,
  xleftmargin=8pt,
  numbers=none,
}

%% Hình ảnh
\usepackage{graphicx}

%% Bảng đẹp hơn
\usepackage{booktabs}
\usepackage{multirow}
\usepackage{subcaption}

%% Đảm bảo float không trọn ra ngoài section
\usepackage[section]{placeins}

%%=============================================================
%% Cách sử dụng trong main.tex:
%%
%%   \documentclass[12pt,a4paper]{report}
%%   %%=============================================================
%% preamble_ai3.tex
%% Preamble dùng chung cho chapter4_4.tex và chapter5_4.tex
%%
%% Đưa vào main.tex bằng:  \input{preamble_ai3}
%% Hoặc copy thẳng các \usepackage cần thiết vào preamble của main.
%%=============================================================

%% Tiếng Việt
\usepackage[T5]{fontenc}
\usepackage[utf8]{inputenc}
\usepackage[vietnamese]{babel}

%% TikZ và các thư viện vẻ hình
\usepackage{tikz}
\usetikzlibrary{
  shapes.geometric,
  arrows.meta,
  positioning,
  fit,
  backgrounds,
  calc,
  decorations.pathreplacing,
  patterns,
}

%% PGFPlots cho biểu đồ
\usepackage{pgfplots}
\pgfplotsset{compat=1.18}
\usepackage{pgfplotstable}

%% Code listings
\usepackage{listings}
\usepackage{xcolor}

%% Màu dùng chung (cần định nghĩa trước khi \input chapter)
\definecolor{cppblue}{RGB}{30,100,200}
\definecolor{pygreen}{RGB}{40,150,80}
\definecolor{warnred}{RGB}{200,60,60}
\definecolor{orng}{RGB}{220,130,30}

%% Cài đặt listings mặc định
\lstset{
  basicstyle=\ttfamily\small,
  keywordstyle=\color{cppblue},
  commentstyle=\color{gray!70},
  stringstyle=\color{orng},
  breaklines=true,
  frame=single,
  rulecolor=\color{gray!40},
  framexleftmargin=4pt,
  xleftmargin=8pt,
  numbers=none,
}

%% Hình ảnh
\usepackage{graphicx}

%% Bảng đẹp hơn
\usepackage{booktabs}
\usepackage{multirow}
\usepackage{subcaption}

%% Đảm bảo float không trọn ra ngoài section
\usepackage[section]{placeins}

%%=============================================================
%% Cách sử dụng trong main.tex:
%%
%%   \documentclass[12pt,a4paper]{report}
%%   \input{preamble_ai3}
%%   \begin{document}
%%     ...
%%     \chapter{Thiết kế và Hiện thực}
%%     ...
%%     \input{chapter4_4}   %% mục 4.4
%%     ...
%%     \chapter{Kết quả và Đánh giá}
%%     ...
%%     \input{chapter5_4}   %% mục 5.4
%%     ...
%%   \end{document}
%%
%% Compile: pdflatex -> pdflatex (chạy 2 lần để có label/ref đúng)
%%=============================================================

%%   \begin{document}
%%     ...
%%     \chapter{Thiết kế và Hiện thực}
%%     ...
%%     %%=============================================================
%% chapter4_4.tex
%% Phần 4.4 — AI #3: Phát hiện vật thể cho pick & place
%%
%% Đây là \input file — thêm vào preamble của main.tex:
%%   \usepackage{tikz}
%%   \usetikzlibrary{shapes.geometric,arrows.meta,positioning,
%%     fit,backgrounds,calc,decorations.pathreplacing}
%%   \usepackage{pgfplots}
%%   \pgfplotsset{compat=1.18}
%%   \usepackage{listings,xcolor,booktabs,multirow,subcaption}
%%=============================================================

\definecolor{cppblue}{RGB}{30,100,200}
\definecolor{pygreen}{RGB}{40,150,80}
\definecolor{warnred}{RGB}{200,60,60}
\definecolor{orng}{RGB}{220,130,30}

\lstset{
  basicstyle=\ttfamily\small,
  keywordstyle=\color{cppblue},
  commentstyle=\color{gray!70},
  stringstyle=\color{orng},
  breaklines=true,
  frame=single,
  rulecolor=\color{gray!40},
  framexleftmargin=4pt,
  xleftmargin=8pt,
  numbers=none,
}

% ================================================================
\section{AI \#3 --- Tự động hoá khâu khởi tạo trong hệ thống pick \& place}
\label{sec:ai3}
% ================================================================

% ----------------------------------------------------------------
\subsection{Bối cảnh: vấn đề của hệ thống visual servoing ban đầu}
\label{subsec:ai3-context}
% ----------------------------------------------------------------

Hệ thống pick \& place chạy trên cánh tay robot Denso VS-6577E 6 bậc tự do,
sử dụng thư viện ViSP (Visual Servoing Platform) để điều khiển bằng thị giác.
Camera USB được gắn trực tiếp lên đầu công tác theo cấu hình \textit{eye-in-hand} ---
camera di chuyển cùng robot và nhìn thẳng về phía vật thể cần gắp trong suốt
quá trình tiếp cận.

Vấn đề cốt lõi của hệ thống ban đầu nằm ở bước khởi tạo: sau khi robot về
home pose cố định, \textbf{người vận hành phải click chuột vào vị trí cylinder}
trên cửa sổ hiển thị camera thì blob tracker (\texttt{vpDot2}) mới bắt đầu
hoạt động được. Visual servo chỉ chạy sau thao tác thủ công đó. Ngoài ra,
nếu tracker mất target trong lúc servo --- điều hoàn toàn có thể xảy ra khi
cylinder bị che khuất hoặc di chuyển ra ngoài vùng sáng --- hệ thống reset
hoàn toàn và lại đòi hỏi click thủ công ở lần tiếp theo. Mỗi chu kỳ
pick \& place đều cần sự can thiệp của người vận hành, loại bỏ khả năng
vận hành tự động không giám sát.

\begin{figure}[htbp]
  \centering
  \begin{tikzpicture}[
    node distance=0.8cm and 0.6cm,
    bx/.style={draw, rounded corners=4pt, text width=3.0cm, align=center,
               minimum height=0.9cm, font=\small},
    robot/.style={bx, draw=cppblue!70, fill=cppblue!8},
    manual/.style={bx, draw=warnred!70, fill=warnred!8},
    aibox/.style={bx, draw=pygreen!70, fill=pygreen!8},
    arr/.style={->, >=Stealth, thick},
    rarr/.style={arr, draw=warnred!70},
    garr/.style={arr, draw=pygreen!70},
  ]

  %% --- CỘT TRÁI: hệ thống gốc ---
  \node[robot]                    (L1) {PREINIT\\{\scriptsize robot $\to$ home pose}};
  \node[manual, below=of L1]      (L2) {Click thủ công\\{\scriptsize chọn cylinder}};
  \node[robot,  below=of L2]      (L3) {JOINT\\{\scriptsize vpDot2 tracking}};
  \node[manual, below=of L3]      (L4) {Mất tracking\\{\scriptsize $\to$ reset toàn bộ}};

  \draw[arr]  (L1) -- (L2);
  \draw[arr]  (L2) -- (L3);
  \draw[rarr] (L3) -- (L4);
  \draw[rarr, dashed] (L4.west)
        to[out=190,in=170,looseness=1.3] (L1.west)
        node[midway, left, font=\scriptsize, text=warnred!70] {reset};

  \node[font=\bfseries\small, text=warnred, above=0.5cm of L1]
       {\textit{Trước cải tiến}};

  %% --- CỘT PHẢI: hệ thống mới ---
  \node[robot,  right=3.8cm of L1] (R1) {PREINIT\\{\scriptsize robot $\to$ home pose}};
  \node[aibox, below=of R1]        (R2) {AI detect\\{\scriptsize EfficientDet-Lite0}};
  \node[robot,  below=of R2]       (R3) {JOINT\\{\scriptsize AI tracking mỗi frame}};
  \node[aibox, below=of R3]        (R4) {AI re-detect\\{\scriptsize tự phục hồi}};

  \draw[arr]  (R1) -- (R2);
  \draw[garr] (R2) -- (R3);
  \draw[garr] (R3) -- (R4);
  \draw[garr, dashed] (R4.east)
        to[out=0,in=0,looseness=1.5] (R3.east)
        node[midway, right, font=\scriptsize, text=pygreen!70] {tiếp tục};

  %% fallback arrow từ AI detect
  \node[manual, right=1.2cm of R3] (FB)
       {{\scriptsize fallback:}\\click thủ công};
  \draw[rarr, dashed] (R2.east) -| (FB.north)
        node[near start, above, font=\scriptsize, text=warnred] {AI fail};

  \node[font=\bfseries\small, text=pygreen!80!black, above=0.5cm of R1]
       {\textit{Sau cải tiến}};

  %% Đường phân cách
  \draw[thick, dashed, gray!50]
      ($(L1.east)!0.5!(R1.west)+(0,0.8)$) -- ++(0,-6.5);

  \end{tikzpicture}
  \caption{So sánh luồng vận hành trước và sau khi tích hợp module AI.
           Bước yêu cầu thao tác thủ công (đỏ) được thay thế bằng AI tự động (xanh lá).
           Fallback về click thủ công vẫn được giữ lại cho trường hợp AI thất bại.}
  \label{fig:ai3-before-after}
\end{figure}

Hình~\ref{fig:ai3-before-after} minh hoạ sự thay đổi này. Module AI đảm nhiệm
toàn bộ việc phát hiện cylinder ở bước khởi tạo và tự phục hồi khi mất target
trong quá trình servo. Hệ thống vẫn giữ nguyên fallback về click thủ công để
đảm bảo tương thích ngược.


% ----------------------------------------------------------------
\subsection{Xây dựng và huấn luyện mô hình EfficientDet-Lite0}
\label{subsec:ai3-model}
% ----------------------------------------------------------------

\subsubsection*{Lựa chọn kiến trúc}

Yêu cầu đặt ra khá cụ thể: mô hình phải đủ nhỏ để triển khai được trên
Google Coral EdgeTPU, phải hỗ trợ INT8 quantization, và tương thích với
\texttt{tflite\_model\_maker} --- framework đã có sẵn trong môi trường phát triển.
\textbf{EfficientDet-Lite0} đáp ứng tất cả: đây là biến thể nhỏ nhất trong
họ EfficientDet được Google tối ưu cho thiết bị edge, input cố định
$320 \times 320$ pixel, backbone là EfficientNet-B0 đã được simplify.

\begin{table}[htbp]
  \centering
  \caption{Thông số mô hình EfficientDet-Lite0}
  \label{tab:ai3-model}
  \small
  \begin{tabular}{ll}
    \toprule
    \textbf{Thuộc tính}       & \textbf{Giá trị} \\
    \midrule
    Kiến trúc                 & EfficientDet-Lite0 \\
    Framework training        & TensorFlow~2.8.4 / \texttt{tflite\_model\_maker}~0.4.3 \\
    Input size                & $320 \times 320$~px, uint8 \\
    Số lớp đối tượng          & 2 (cylinder~--~lớp~0, cube~--~lớp~1) \\
    Số epochs / batch size    & 50 / 8 \\
    Transfer learning         & Frozen backbone, fine-tune detection head \\
    Target hardware           & Google Coral USB Accelerator (EdgeTPU) \\
    \bottomrule
  \end{tabular}
\end{table}

Mô hình được huấn luyện để phát hiện \textbf{hai lớp}: cylinder và cube nhựa PLA
in 3D. Việc đưa thêm cube vào tập huấn luyện là có chủ ý --- giúp mô hình học
cách phân biệt hai loại vật thể, từ đó giảm false positive khi cả hai cùng
xuất hiện trong cảnh. Ở phía tích hợp với ViSP, chỉ kết quả cylinder
(\texttt{class\_id~==~0}) được dùng; dự đoán cube bị lọc bỏ.

\subsubsection*{Tập dữ liệu}

Toàn bộ ảnh được thu trực tiếp từ camera eye-in-hand gắn trên robot Denso,
đảm bảo điều kiện thực tế sát với lúc triển khai. Ảnh bao gồm nhiều cấu hình:
góc chụp khác nhau, khoảng cách thay đổi, và số lượng vật thể từ 1 đến 5 trong
một cảnh.

\begin{figure}[htbp]
  \centering
  \begin{tikzpicture}
    \begin{axis}[
      xbar, bar width=0.42cm,
      xlabel={Số lượng ảnh (xấp xỉ)},
      ylabel={Danh mục},
      symbolic y coords={
        Background,
        2 cube + 3 cyl,
        2 cylinders,
        3 cylinders,
        1 cyl + 1 cube,
        1 cube,
        1 cylinder,
        Cylinder (gốc)
      },
      ytick=data,
      xmin=0, xmax=175,
      width=9.5cm, height=6.8cm,
      nodes near coords,
      nodes near coords align={horizontal},
      every node near coord/.style={font=\scriptsize},
      tick label style={font=\small},
      xlabel style={font=\small},
      ylabel style={font=\small},
      enlarge y limits=0.08,
    ]
    \addplot[fill=cppblue!55, draw=cppblue!80] coordinates {
      (38,  Background)
      (28,  2 cube + 3 cyl)
      (52,  2 cylinders)
      (43,  3 cylinders)
      (68,  1 cyl + 1 cube)
      (38,  1 cube)
      (108, 1 cylinder)
      (144, Cylinder (gốc))
    };
    \end{axis}
  \end{tikzpicture}
  \caption{Phân bố tập dữ liệu theo danh mục ảnh (~534 ảnh tổng,
           633 annotation cylinder, 170 annotation cube).
           Các cảnh đa vật thể giúp mô hình học phát hiện từng đối tượng
           riêng lẻ trong môi trường phức tạp.}
  \label{fig:ai3-dataset}
\end{figure}

Annotation được thực hiện bằng \texttt{labelme} theo format bounding box hình chữ
nhật. Workflow:
\begin{enumerate}
  \item Vẽ bounding box cho từng cylinder và cube bằng \texttt{labelme}
  \item Chuyển đổi labelme JSON $\to$ COCO JSON thống nhất
        (script \texttt{labelme\_to\_coco.py})
  \item Tập dữ liệu được split theo tỉ lệ $70/15/15$ (train/val/test),
        seed cố định để đảm bảo tính tái lập
\end{enumerate}

Phân chia ba chiều thay vì hai chiều đảm bảo tập test hoàn toàn độc lập ---
không được nhìn thấy trong suốt quá trình training hoặc điều chỉnh siêu tham số.

\subsubsection*{Pipeline huấn luyện}

\begin{figure}[htbp]
  \centering
  \begin{tikzpicture}[
    node distance=0.45cm and 0.15cm,
    st/.style={draw, rounded corners=3pt, fill=gray!10, draw=gray!55,
               text width=2.0cm, align=center, minimum height=0.85cm, font=\small},
    blue_st/.style={st, fill=cppblue!10, draw=cppblue!55},
    out_st/.style={st, fill=orng!20, draw=orng!70},
    arr/.style={->, >=Stealth, thick, gray!60},
    lbl/.style={font=\scriptsize, text=gray!80},
  ]

  \node[blue_st] (s1) {Annotation\\{\scriptsize labelme}};
  \node[blue_st, right=of s1] (s2) {COCO\\JSON};
  \node[blue_st, right=of s2] (s3) {Pascal\\VOC XML};
  \node[blue_st, right=of s3] (s4) {Training\\{\scriptsize 50 epochs}};
  \node[out_st, right=of s4] (s5) {SavedModel\\{\scriptsize Float32}};
  \node[out_st, below=0.55cm of s5] (s6) {INT8\\TFLite};
  \node[out_st, below=0.55cm of s6] (s7) {EdgeTPU\\TFLite};

  \draw[arr] (s1) -- (s2) node[lbl,above,midway] {to\_coco.py};
  \draw[arr] (s2) -- (s3) node[lbl,above,midway] {tflite\_mm};
  \draw[arr] (s3) -- (s4);
  \draw[arr] (s4) -- (s5) node[lbl,above,midway] {export};
  \draw[arr] (s5) -- (s6) node[lbl,right,midway] {PTQ};
  \draw[arr] (s6) -- (s7) node[lbl,right,midway] {edgetpu\_compiler};

  \end{tikzpicture}
  \caption{Pipeline huấn luyện và xuất mô hình. Từ ảnh thô đến ba file
           mô hình phục vụ các môi trường triển khai khác nhau.
           Bước chuyển đổi sang Pascal VOC XML là yêu cầu bắt buộc của
           \texttt{tflite\_model\_maker}.}
  \label{fig:ai3-pipeline}
\end{figure}

Transfer learning với \texttt{train\_whole\_model=False} --- chỉ fine-tune detection
head, giữ nguyên backbone đã pretrain --- phù hợp với tập dữ liệu tương đối nhỏ
($\sim$374 ảnh train) và giảm nguy cơ overfitting đáng kể so với
huấn luyện từ đầu.


% ----------------------------------------------------------------
\subsection{Quantization INT8 và biên dịch cho EdgeTPU}
\label{subsec:ai3-quant}
% ----------------------------------------------------------------

Mô hình sau training là Float32 --- không chạy được trực tiếp trên Coral EdgeTPU.
EdgeTPU chỉ thực thi mô hình INT8 quantized, trong đó trọng số và activation
được biểu diễn bằng số nguyên 8-bit. Lợi ích: kích thước giảm $\sim4\times$,
và phần cứng EdgeTPU đạt tốc độ inference cực nhanh ($\sim$4--15~ms/frame)
nhờ tận dụng phép tính số nguyên.

Post-Training Quantization (PTQ) sử dụng tập train làm \textit{representative
dataset} để calibrate scale và zero\_point cho từng layer --- đây là bước bắt
buộc để duy trì độ chính xác của INT8 model gần với Float32 baseline.

File \texttt{\_edgetpu.tflite} sau khi biên dịch chứa custom op
\texttt{edgetpu-custom-op} mà TFLite runtime thông thường
\textit{không thể thực thi trên CPU}, gây lỗi:

\begin{lstlisting}[language={}]
RuntimeError: Encountered unresolved custom op: edgetpu-custom-op
\end{lstlisting}

Do đó hệ thống \textbf{duy trì hai file riêng biệt} và lựa chọn tự động:

\begin{center}
\small
\begin{tabular}{lll}
  \toprule
  File & Runtime & Deployment \\
  \midrule
  \texttt{cylinder\_detector\_int8\_edgetpu.tflite} & pycoral   & Coral EdgeTPU \\
  \texttt{cylinder\_detector\_int8.tflite}          & TFLite/ai-edge-litert & CPU \\
  \bottomrule
\end{tabular}
\end{center}


% ----------------------------------------------------------------
\subsection{Kiến trúc tích hợp: C++ ViSP ↔ Python subprocess}
\label{subsec:ai3-arch}
% ----------------------------------------------------------------

\subsubsection*{Quyết định thiết kế}

Thay vì link trực tiếp thư viện TFLite C++ vào build system của ViSP
(đòi hỏi thay đổi lớn vào CMakeLists và phụ thuộc version build cụ thể),
hệ thống áp dụng kiến trúc \textbf{subprocess model}: C++ gọi Python script
như một tiến trình con qua \texttt{popen()}, giao tiếp một chiều qua stdout.
Cách này giúp module AI hoàn toàn tách biệt --- cập nhật mô hình không cần
recompile code robot.

\begin{figure}[htbp]
  \centering
  \begin{tikzpicture}[
    node distance=0.38cm and 0.5cm,
    comp/.style={draw, rounded corners=4pt, text width=3.1cm, align=center,
                 minimum height=0.82cm, font=\small},
    cpp/.style={comp, fill=cppblue!10, draw=cppblue!55},
    py/.style={comp,  fill=pygreen!10, draw=pygreen!60},
    tmp/.style={comp,  fill=orng!18, draw=orng!65, text width=2.5cm,
                minimum height=0.7cm, font=\scriptsize},
    arr/.style={->, >=Stealth, thick},
    ]

  %% ---- C++ side ----
  \node[cpp] (cap)   {Camera capture\\{\scriptsize vpImage{$<$}uchar{$>$}}};
  \node[cpp, below=of cap]  (save)  {Save frame\\{\scriptsize cv::imwrite()}};
  \node[cpp, below=of save] (pop)   {popen() subprocess\\{\scriptsize build cmd}};
  \node[cpp, below=of pop]  (sel)   {select() timeout 3s\\{\scriptsize POSIX wait}};
  \node[cpp, below=of sel]  (parse) {Parse stdout\\{\scriptsize sscanf SUCCESS}};
  \node[cpp, below=of parse](feat)  {vpFeatureBuilder\\{\scriptsize create feature}};

  \draw[arr](cap)--(save); \draw[arr](save)--(pop);
  \draw[arr](pop)--(sel);  \draw[arr](sel)--(parse); \draw[arr](parse)--(feat);

  \node[draw=cppblue!50, dashed, rounded corners=6pt,
        fit=(cap)(save)(pop)(sel)(parse)(feat),
        inner sep=7pt] (cbox) {};
  \node[above=2pt of cbox, font=\bfseries\small, text=cppblue]
       {C++ (ViSP main loop)};

  %% ---- Temp file ----
  \node[tmp, right=2.6cm of save] (jpg)
       {/tmp/visp\_ai\_frame.jpg};

  %% ---- Python side ----
  \node[py, right=2.6cm of pop]   (load) {Load model\\{\scriptsize fallback chain}};
  \node[py, below=of load]        (pre)  {Preprocess\\{\scriptsize resize 320$\times$320}};
  \node[py, below=of pre]         (inf)  {Inference\\{\scriptsize invoke()}};
  \node[py, below=of inf]         (post) {Post-process\\{\scriptsize dequant+filter}};
  \node[py, below=of post]        (out)  {Print stdout\\{\scriptsize SUCCESS u v conf}};

  \draw[arr](load)--(pre); \draw[arr](pre)--(inf);
  \draw[arr](inf)--(post); \draw[arr](post)--(out);

  \node[draw=pygreen!55, dashed, rounded corners=6pt,
        fit=(load)(pre)(inf)(post)(out),
        inner sep=7pt] (pbox) {};
  \node[above=2pt of pbox, font=\bfseries\small, text=pygreen!80!black]
       {Python subprocess};

  %% ---- Data flow arrows ----
  \draw[arr, orng!80]  (save.east) -- (jpg.west)  node[midway,above,font=\scriptsize]{write};
  \draw[arr, orng!80]  (jpg.east)  -- (load.west) node[midway,above,font=\scriptsize]{read};
  \draw[arr, cppblue!60](pop.east)  -- (load.west) node[midway,above,font=\scriptsize,yshift=2pt]{spawn};
  \draw[arr, pygreen!70, bend right=25](out.west)  to (parse.east)
        node[midway,below,font=\scriptsize]{stdout pipe};

  \end{tikzpicture}
  \caption{Kiến trúc tích hợp C++ ViSP $\leftrightarrow$ Python subprocess.
           C++ lưu frame hiện tại ra file tạm, spawn Python làm child process
           qua \texttt{popen()}, đọc kết quả từ stdout sau khi qua timeout guard.}
  \label{fig:ai3-arch}
\end{figure}

\subsubsection*{Protocol trao đổi dữ liệu}

Output của Python script là một dòng duy nhất trên stdout:

\begin{lstlisting}[language={}]
SUCCESS 272.86 190.74 0.8362 238.42 166.89 68.88 47.68
%        u      v      conf   x_min  y_min  width  height
\end{lstlisting}

hoặc \texttt{FAILURE~<lý~do>} nếu không phát hiện được. Tất cả debug log
đi ra \texttt{stderr} --- C++ chỉ đọc stdout --- tránh nhiễu kết quả parse.

Đáng chú ý: trong ViSP, \texttt{vpImagePoint} được định nghĩa là \texttt{(row,~col)}
tương ứng \texttt{(v,~u)}, không phải \texttt{(u,~v)} như thông thường:

\begin{lstlisting}[language=C++]
detected_center = vpImagePoint(v, u);  // row=v, col=u
\end{lstlisting}

Tham số \texttt{hint} cho phép Python chọn cylinder \textit{gần nhất với vị
trí đã biết} thay vì cylinder có confidence cao nhất --- cơ chế này tránh
hiện tượng servo ``nhảy'' sang cylinder khác khi scores thay đổi giữa các frame.


% ----------------------------------------------------------------
\subsection{State machine: PREINIT \texorpdfstring{$\to$}{→} INIT (AI) \texorpdfstring{$\to$}{→} JOINT \texorpdfstring{$\to$}{→} ...}
\label{subsec:ai3-sm}
% ----------------------------------------------------------------

Hệ thống vận hành theo vòng lặp chính liên tục với 6 trạng thái.
Mỗi iteration đọc một frame camera và thực thi logic trạng thái hiện tại.

\begin{figure}[htbp]
  \centering
  \begin{tikzpicture}[
    node distance=1.0cm and 2.0cm,
    st/.style={draw, rounded corners=4pt, fill=cppblue!10, draw=cppblue!55,
               text width=2.3cm, align=center, minimum height=1.0cm, font=\small},
    ai/.style={st, fill=pygreen!10, draw=pygreen!60},
    arr/.style={->, >=Stealth, thick},
    lbl/.style={font=\scriptsize, text=gray!70},
  ]

  %% Top row
  \node[st]  (pre)  {PREINIT\\{\scriptsize home pose}};
  \node[ai,  right=of pre]   (ini)  {INIT\\{\scriptsize AI detect}};
  \node[ai,  right=of ini]   (jnt)  {JOINT\\{\scriptsize AI tracking}};
  \node[st,  right=of jnt]   (app)  {APPROACH\\{\scriptsize chờ RC5}};

  %% Bottom row
  \node[st,  below=1.8cm of app]  (grp)  {GRIPPER\\{\scriptsize đóng gripper}};
  \node[st,  left=of grp]         (cls)  {CLASSIFIED\\{\scriptsize vận chuyển}};

  %% Arrows
  \draw[arr](pre)--(ini)  node[lbl,above,midway]{home\\reached};
  \draw[arr](ini)--(jnt)  node[lbl,above,midway]{AI init\\ok};
  \draw[arr](jnt)--(app)  node[lbl,above,midway]{$|e|<5{\times}10^{-3}$};
  \draw[arr](app)--(grp)  node[lbl,right,midway]{RC5\\signal};
  \draw[arr](grp)--(cls)  node[lbl,below,midway]{grasped};
  \draw[arr](cls.west) to[out=180,in=270,looseness=0.8] (pre.south)
        node[lbl,left,pos=0.5]{cycle\\restart};

  %% AI self-recovery
  \draw[arr, pygreen!70, dashed]
        (jnt.north) to[out=100,in=80,looseness=2.5] (jnt.north)
        node[above=0.7cm of jnt, font=\scriptsize, text=pygreen!80!black]
        {AI re-detect};

  %% Fallback from INIT
  \draw[arr, warnred!70, dashed]
        (ini.south) -- ++(0,-0.7)
        node[below, font=\scriptsize, text=warnred]
        {fallback: click thủ công};

  \end{tikzpicture}
  \caption{State machine của hệ thống pick \& place sau cải tiến.
           Trạng thái tô xanh lá là nơi AI được gọi.
           Vòng lặp tự phục hồi ở JOINT (mũi tên đứt nét) cho phép hệ thống
           tiếp tục mà không cần reset khi mất detection.}
  \label{fig:ai3-sm}
\end{figure}

\paragraph{PREINIT.}
Robot gửi lệnh joint position về home pose
$[0^\circ, 0^\circ, 90^\circ, 0^\circ, 90^\circ, 0^\circ]$.
Servo task được cấu hình với interaction matrix Desired, pseudo-inverse,
gain $\lambda = 0.4$.

\paragraph{INIT.}
Chờ robot đạt home pose (sai số từng joint $< 0.01$) và gripper mở xong.
Khi cả hai điều kiện thỏa, \texttt{detectCylinderWithAI()} được gọi
không có hint --- Python chọn cylinder có confidence cao nhất.
Kết quả trực tiếp tạo visual feature \texttt{vpFeaturePoint}
và add vào servo task. Nếu AI thất bại, \texttt{ai\_hint}
giữ giá trị mặc định (fallback thủ công).

\paragraph{JOINT.}
Khác với cách tiếp cận truyền thống dùng blob tracker mỗi frame,
hệ thống hiện tại \textbf{gọi AI trực tiếp ở mỗi iteration}
với \texttt{hint = \&ai\_hint}. Control law:
\begin{equation}
  \mathbf{v} = -\lambda\,\mathbf{L}_{s^*}^{+}\,\mathbf{e},
  \qquad \mathbf{q}_{k+1} = \mathbf{q}_k + \mathbf{v}\,\Delta t
  \label{eq:ibvs}
\end{equation}
với $\Delta t = 0.5\,$s. Điều kiện hội tụ: $|e_x| < 5{\times}10^{-3}$
và $|e_y| < 5{\times}10^{-3}$ (normalized image coordinates).
Khi đạt, gửi \texttt{"OKE\textbackslash r"} qua UART đến RC5 controller.


% ----------------------------------------------------------------
\subsection{Các vấn đề kỹ thuật và giải pháp}
\label{subsec:ai3-issues}
% ----------------------------------------------------------------

Một số vấn đề nảy sinh trong quá trình tích hợp mà khó thấy trước khi thực
sự chạy hệ thống.

\subsubsection*{P1: Xác định output tensor sai lệch}

\texttt{tflite\_model\_maker} 0.4.x đặt tên tensor theo thứ tự riêng của nó:

\begin{center}
\small
\begin{tabular}{ccc}
  \toprule
  Tensor & Shape & Nội dung \\
  \midrule
  \texttt{:0} & \texttt{[1]}     & num\_detections \\
  \texttt{:1} & \texttt{[1,\,N]} & scores \quad \textbf{← cần lấy cái này} \\
  \texttt{:2} & \texttt{[1,\,N]} & class IDs \quad \textbf{← không phải cái này} \\
  \texttt{:3} & \texttt{[1,\,N,\,4]} & boxes \\
  \bottomrule
\end{tabular}
\end{center}

Thứ tự này không khớp với API của pycoral, vốn giả định positional order
\texttt{boxes/classes/scores/count = 0/1/2/3}. Dùng API đó sẽ đọc
num\_detections (raw uint8 = 255) như scores và crash với \texttt{IndexError}.
Giải pháp: xác định tensor \textbf{theo shape động} --- boxes là
\texttt{[1,N,4]}, scores là tensor \texttt{[1,N]} có name suffix nhỏ nhất
(\texttt{:1} chứ không phải \texttt{:2}).

\subsubsection*{P2: Python 3.10 không có \texttt{tflite-runtime}}

Google không publish wheel \texttt{tflite-runtime} cho Python 3.10+.
Giải pháp là xây dựng \textbf{fallback chain 4 tầng}:

\begin{figure}[htbp]
  \centering
  \begin{tikzpicture}[
    node distance=0.4cm,
    fb/.style={draw, rounded corners=3pt, text width=7.5cm, align=left,
               minimum height=0.85cm, font=\small, inner sep=7pt},
    t1/.style={fb, fill=purple!12, draw=purple!55},
    t2/.style={fb, fill=teal!10,   draw=teal!50},
    t3/.style={fb, fill=cppblue!10, draw=cppblue!50},
    t4/.style={fb, fill=gray!12,   draw=gray!50},
    fa/.style={fb, fill=warnred!10, draw=warnred!50},
    arr/.style={->, >=Stealth, thick, dashed, gray!55},
    lbl/.style={font=\scriptsize, text=warnred!80},
  ]
  \node[t1] (a) {\textbf{1.}~\texttt{pycoral.utils.edgetpu}\quad{\scriptsize Coral EdgeTPU, Python 3.9}};
  \node[t2, below=of a] (b) {\textbf{2.}~\texttt{ai\_edge\_litert.interpreter}\quad{\scriptsize Google's new pkg, Python 3.8--3.12}};
  \node[t3, below=of b] (c) {\textbf{3.}~\texttt{tflite\_runtime.interpreter}\quad{\scriptsize Legacy, Python 3.9}};
  \node[t4, below=of c] (d) {\textbf{4.}~\texttt{tensorflow.lite.Interpreter}\quad{\scriptsize Heavy fallback}};
  \node[fa, below=of d] (e) {\textbf{FAILURE}:~\texttt{no\_tflite\_backend\_available}};

  \draw[arr](a)--(b) node[lbl,right,midway]{ImportError / Coral not connected};
  \draw[arr](b)--(c) node[lbl,right,midway]{ImportError};
  \draw[arr](c)--(d) node[lbl,right,midway]{ImportError};
  \draw[arr](d)--(e) node[lbl,right,midway]{ImportError};
  \end{tikzpicture}
  \caption{Python backend fallback chain trong \texttt{detect\_cylinder.py}.
           Script tự động thử từng backend theo thứ tự ưu tiên,
           tương thích với cả môi trường có Coral lẫn CPU-only.}
  \label{fig:ai3-fallback}
\end{figure}

\texttt{ai-edge-litert} là gói chính thức của Google thay thế \texttt{tflite-runtime},
hỗ trợ Python 3.8--3.12 với API giống hệt nhau --- là giải pháp cho vấn đề Python 3.10.

\subsubsection*{P3: JSON comment trong config}

\texttt{config.json} dùng comment \texttt{//} để ghi chú, nhưng cả Python
\texttt{json.load()} và \texttt{nlohmann::json} đều từ chối mặc định.
Phía C++: \texttt{nlohmann::json::parse(file, nullptr, true, true)} với
\texttt{ignore\_comments=true} (có từ nlohmann/json~3.9.0).
Phía Python: strip comment lines thủ công trước khi \texttt{json.loads()}.


% ----------------------------------------------------------------
\subsection{Thông số cấu hình visual servo}
\label{subsec:ai3-servo}
% ----------------------------------------------------------------

Hệ thống sử dụng \textbf{Image-Based Visual Servoing (IBVS)} với feature là
tọa độ 2D của cylinder trong không gian ảnh normalized $\mathbf{s}=(x,y)$,
desired position $\mathbf{s}^*=(0,0,1)$ --- cylinder phải nằm đúng tâm ảnh.

\begin{table}[htbp]
  \centering
  \caption{Thông số điều khiển visual servo}
  \label{tab:ai3-servo}
  \small
  \begin{tabular}{ll}
    \toprule
    \textbf{Tham số}            & \textbf{Giá trị} \\
    \midrule
    Loại điều khiển             & Eye-in-hand, \texttt{EYEINHAND\_L\_cVe\_eJe} \\
    Interaction matrix          & Desired ($\mathbf{L}_{s^*}$), pseudo-inverse \\
    Gain $\lambda$              & 0.4 \\
    Bước thời gian $\Delta t$   & 0.5\,s \\
    Ngưỡng hội tụ               & $|e_x| < 5{\times}10^{-3}$ và $|e_y| < 5{\times}10^{-3}$ \\
    Desired position            & $(x^*, y^*) = (0,\,0)$ --- tâm ảnh, $Z = 1$ \\
    \bottomrule
  \end{tabular}
\end{table}

Dùng $\mathbf{L}_{s^*}$ (tại vị trí mong muốn) thay vì $\mathbf{L}_{s_{cur}}$
giúp tránh singularity khi feature tiến gần desired point và giữ ổn định
theo ổn định hơn trong giai đoạn hội tụ. Velocity được saturate trước khi
gửi xuống robot để đảm bảo không vượt giới hạn vật lý từng khớp.
   %% mục 4.4
%%     ...
%%     \chapter{Kết quả và Đánh giá}
%%     ...
%%     %%=============================================================
%% chapter5_4.tex
%% Phần 5.4 — Kết quả đánh giá: AI #3 Phát hiện vật thể
%%
%% \input file — cần thêm vào preamble:
%%   \usepackage{tikz}
%%   \usetikzlibrary{arrows.meta,positioning,patterns}
%%   \usepackage{pgfplots}
%%   \pgfplotsset{compat=1.18}
%%   \usepackage{booktabs,multirow,subcaption,xcolor}
%%=============================================================

% ================================================================
\section{Kết quả và Đánh giá: AI \#3 --- Phát hiện vật thể}
\label{sec:ai3-results}
% ================================================================

% ----------------------------------------------------------------
\subsection{Kết quả detection model: COCO metrics}
\label{subsec:ai3-coco}
% ----------------------------------------------------------------

Mô hình được đánh giá trên cả validation set và test set độc lập
theo chuẩn \textbf{COCO evaluation protocol} --- tính trung bình mAP qua
các ngưỡng IoU từ 0.50 đến 0.95 với bước 0.05.

\begin{figure}[htbp]
  \centering
  \begin{tikzpicture}
    \begin{axis}[
      ybar, bar width=0.45cm,
      width=10cm, height=7cm,
      symbolic x coords={mAP (0.50:0.95), AP50, AP75, AR},
      xtick=data,
      x tick label style={font=\small, align=center},
      ymin=0, ymax=1.15,
      ylabel={Score},
      ylabel style={font=\small},
      legend style={at={(0.5,1.03)}, anchor=south, legend columns=2, font=\small},
      legend cell align=left,
      nodes near coords,
      every node near coord/.style={font=\scriptsize, /pgf/number format/fixed,
                                    /pgf/number format/precision=3},
      enlarge x limits=0.25,
      grid=major, grid style={dashed, gray!30},
      axis line style={gray!60},
      tick style={gray!60},
    ]
    \addplot[fill=cppblue!70, draw=cppblue!90] coordinates {
      (mAP (0.50:0.95), 0.791)
      (AP50, 0.990)
      (AP75, 0.990)
      (AR,   0.869)
    };
    \addplot[fill=orng!80, draw=orng!90] coordinates {
      (mAP (0.50:0.95), 0.799)
      (AP50, 1.000)
      (AP75, 0.969)
      (AR,   0.856)
    };
    \legend{Validation, Test}
    \end{axis}
  \end{tikzpicture}
  \caption{COCO metrics trên validation set và test set. AP50 trên test set
           đạt 1.000 cho thấy 100\% cylinder được phát hiện đúng ở ngưỡng IoU
           lỏng lẻ (50\%). mAP$\approx$0.80 phản ánh vẫn còn đây đó những trường hợp
           định vị bounding box chưa được chính xác ở ngưỡng IoU cao.}
  \label{fig:ai3-coco}
\end{figure}

Bảng~\ref{tab:ai3-metrics} tổng hợp số liệu chi tiết:

\begin{table}[htbp]
  \centering
  \caption{COCO metrics trên validation và test set}
  \label{tab:ai3-metrics}
  \small
  \begin{tabular}{lcc}
    \toprule
    \textbf{Metric}                    & \textbf{Validation} & \textbf{Test} \\
    \midrule
    mAP (IoU 0.50:0.95)               & 0.791               & \textbf{0.799} \\
    AP50 (IoU $\geq$ 0.50)            & 0.990               & \textbf{1.000} \\
    AP75 (IoU $\geq$ 0.75)            & 0.990               & 0.969 \\
    AR (maxDets=100)                   & 0.869               & 0.856 \\
    \midrule
    Precision (IoU=0.50, conf=0.50)    & ---                 & \textbf{1.000} \\
    Recall    (IoU=0.50, conf=0.50)    & ---                 & 0.990 \\
    F1-Score                           & ---                 & \textbf{0.995} \\
    TP / FP / FN                       & ---                 & 95 / 0 / 1 \\
    \bottomrule
  \end{tabular}
\end{table}

\begin{figure}[htbp]
  \centering
  \begin{tikzpicture}
    \begin{axis}[
      ybar, bar width=0.55cm,
      width=8.5cm, height=6cm,
      symbolic x coords={Precision, Recall, F1-Score},
      xtick=data,
      x tick label style={font=\small},
      ymin=0.96, ymax=1.04,
      ylabel={Score},
      ylabel style={font=\small},
      nodes near coords,
      every node near coord/.style={font=\small, /pgf/number format/fixed,
                                    /pgf/number format/precision=3},
      enlarge x limits=0.3,
      grid=major, grid style={dashed, gray!30},
      axis line style={gray!60},
    ]
    \addplot[fill=pygreen!65, draw=pygreen!80] coordinates {
      (Precision, 1.000)
      (Recall,    0.990)
      (F1-Score,  0.995)
    };
    \end{axis}
  \end{tikzpicture}
  \caption{Precision, Recall và F1-Score trên test set (PASCAL VOC protocol,
           IoU$\geq$0.50, confidence$\geq$0.50, chỉ lớp cylinder).
           Precision = 1.000 và FP = 0 chứng tỏ không có false positive
           nào ở ngưỡng 0.50. FN = 1 (một cylinder bị bỏ sót) giải thích
           Recall = 0.990.}
  \label{fig:ai3-prf}
\end{figure}

Giá trị AP50 = 1.000 trên test set à ngưỡng IoU lỏng lẻ (50\%) phản ánh
mô hình phát hiện được toàn bộ cylinder. mAP = 0.799 (trung bình
qua IoU 0.50--0.95) thấp hơn một chút do mô hình EfficientDet-Lite0
đơn giản vẫn còn giới hạn trong việc định vị bounding box chính xác
ở các ngưỡng IoU cao ($\geq$0.75) --- điều phổ biến với các kiến trúc
nhe hơn.


% ----------------------------------------------------------------
\subsection{Phân tích confidence score: true positive vs false positive}
\label{subsec:ai3-conf}
% ----------------------------------------------------------------

Một trong những điều quan trọng khi triển khai thực tế là khoảng cách
phân tách giữa true positive và false positive --- điều này quyết định
ngưỡng confidence nào thực tế dùng được mà không sai kết quả.

Kiểm tra trên 218 ảnh cylinder (true positives) và 10 ảnh cube
(false positives):

\begin{figure}[htbp]
  \centering
  \begin{tikzpicture}
    \begin{axis}[
      width=11cm, height=5.5cm,
      xlabel={Confidence score},
      ylabel={},
      xmin=0, xmax=1.05,
      ymin=0, ymax=3,
      xtick={0,0.1,...,1.0},
      ytick={0.7, 2.0},
      yticklabels={\small False Positive\\(cube, $n$=10), \small True Positive\\(cylinder, $n$=218)},
      y tick label style={align=right, font=\small},
      x tick label style={font=\small},
      xlabel style={font=\small},
      grid=both, grid style={dashed, gray!25},
      axis line style={gray!60},
    ]

    %% --- True Positive band (y=2) ---
    %% Most scores: 0.85-0.98 (majority)
    \addplot[only marks, mark=|, mark size=5pt, mark options={draw=pygreen!70, line width=1.2pt},
             opacity=0.5]
    coordinates {
      (0.29,2)(0.61,2)(0.74,2)
      (0.82,2)(0.83,2)(0.84,2)(0.85,2)(0.86,2)(0.87,2)(0.88,2)
      (0.89,2)(0.90,2)(0.91,2)(0.92,2)(0.93,2)(0.94,2)(0.94,2)
      (0.95,2)(0.95,2)(0.96,2)(0.96,2)(0.97,2)(0.97,2)(0.98,2)
    };

    %% --- False Positive band (y=0.7) ---
    \addplot[only marks, mark=|, mark size=5pt, mark options={draw=warnred!70, line width=1.2pt},
             opacity=0.6]
    coordinates {
      (0.08,0.7)(0.13,0.7)(0.17,0.7)(0.21,0.7)
      (0.25,0.7)(0.28,0.7)(0.30,0.7)(0.33,0.7)(0.35,0.7)(0.36,0.7)
    };

    %% --- Threshold line at 0.50 ---
    \addplot[thick, dashed, draw=gray!80, domain=0:3]
    coordinates {(0.50, 0) (0.50, 3)};
    \node[font=\scriptsize, text=gray!80, rotate=90]
         at (axis cs: 0.52, 1.5) {threshold = 0.50};

    %% --- Gap annotation ---
    \draw[->, >=Stealth, thick, gray!70]
         (axis cs:0.38, 2.5) -- (axis cs:0.48, 2.5);
    \draw[->, >=Stealth, thick, gray!70]
         (axis cs:0.34, 2.5) -- (axis cs:0.27, 2.5);
    \node[font=\scriptsize, text=gray!80] at (axis cs:0.37, 2.7) {gap = 0.34};

    %% --- Median annotation for TP ---
    \addplot[only marks, mark=diamond*, mark size=4pt,
             mark options={draw=pygreen!80!black, fill=pygreen!60}]
    coordinates {(0.94, 2)};
    \node[font=\scriptsize, text=pygreen!80!black, above]
         at (axis cs:0.94, 2) {median 0.94};

    %% --- Max FP annotation ---
    \addplot[only marks, mark=diamond*, mark size=4pt,
             mark options={draw=warnred!80, fill=warnred!50}]
    coordinates {(0.36, 0.7)};
    \node[font=\scriptsize, text=warnred, below]
         at (axis cs:0.36, 0.7) {max 0.36};

    \end{axis}
  \end{tikzpicture}
  \caption{Phân tách confidence score giữa true positive (cylinder, xanh lá)
           và false positive (cube, đỏ). Mỗi vạch đứng đại diện một detection.
           Ngưỡng 0.50 nằm chính giữa khoảng gap 0.34,
           loại bỏ toàn bộ false positive mà vẫn chấp nhận 99.5\% true positive.}
  \label{fig:ai3-conf}
\end{figure}

Khoảng cách 0.34 giữa điểm false positive cao nhất (0.36) và 99.1\%
true positive ($\geq$0.70) xác nhận mô hình đã học được đặc trưng
phân biệt mạnh giữa cylinder và cube. Ngưỡng hiện tại 0.50 nằm an toàn
giữa hai vùng phân bố này.

Ba ảnh outlier có score thấp (0.29, 0.61, 0.74) tương ứng với cylinder
bị nhìn từ góc bất thường hoặc bị che khuất một phần. Ỡ runtime,
các trường hợp này fallback về click thủ công --- hành vi an toàn.

\begin{table}[htbp]
  \centering
  \caption{Phân tích confidence score trên 218 ảnh cylinder + 10 ảnh cube}
  \label{tab:ai3-conf}
  \small
  \begin{tabular}{lcc}
    \toprule
    & \textbf{True Positive (cylinder)} & \textbf{False Positive (cube)} \\
    \midrule
    Score range      & 0.29 -- 0.98          & 0.08 -- 0.36 \\
    Median           & \textbf{0.94}         & --- \\
    $\geq$ 0.50      & 99.5\% (217/218)      & \textbf{0/10 (0\%)} \\
    $\geq$ 0.70      & 99.1\% (216/218)      & 0/10 (0\%) \\
    $\geq$ 0.80      & 95.9\% (209/218)      & 0/10 (0\%) \\
    \midrule
    Gap ph\^an t\'ach & \multicolumn{2}{c}{$0.70 - 0.36 = 0.34$} \\
    \bottomrule
  \end{tabular}
\end{table}


% ----------------------------------------------------------------
\subsection{Kết quả tích hợp hệ thống}
\label{subsec:ai3-integration}
% ----------------------------------------------------------------

\begin{table}[htbp]
  \centering
  \caption{So sánh khả năng trước và sau khi tích hợp AI}
  \label{tab:ai3-comparison}
  \small
  \begin{tabular}{lll}
    \toprule
    \textbf{Tiêu chí}                    & \textbf{Hệ thống gốc} & \textbf{Sau cải tiến} \\
    \midrule
    Khởi tạo mỗi chu kỳ               & Click thủ công (bắt buộc) & Tự động (AI) \\
    Tỷ lệ init thành công             & 100\% (do manual)  & $\sim$99.5\% (AI) \\
    Fallback khi AI thất bại           & Không có              & Click thủ công \\
    Phục hồi khi mất tracking          & Reset toàn bộ        & AI re-detect tự động \\
    Chạy không giám sát                & Không thể              & Có (trường hợp bình thường) \\
    \bottomrule
  \end{tabular}
\end{table}

Thiết kế hiện tại gọi AI trực tiếp \textbf{mỗi iteration} của servo
loop thay vì dùng blob tracker. Điều này có một hàm ý đáng chú ý:
\textit{không có tracking failure} theo nghĩa truyền thống --- mỗi frame
là một lần detection độc lập. Tính ổn định của servo phụ thuộc vào
sự nhất quán của AI detection qua các frame liên tiếp. Cơ chế \texttt{--hint}
(truyền vị trí frame trước) giúp được phần nào khi có nhiều cylinder
trong cảnh.


% ----------------------------------------------------------------
\subsection{Phân tích latency: subprocess CPU vs EdgeTPU}
\label{subsec:ai3-latency}
% ----------------------------------------------------------------

Mỗi lần gọi \texttt{detectCylinderWithAI()} bao gồm nhiều thành phần
latency khác nhau. Điều ít người để ý là: phần overhead lớn nhất
không phải là inference mà là \textbf{Python startup time} --- chi phí
mở một process mới mỗi lần gọi \texttt{popen()}.

\begin{figure}[htbp]
  \centering
  \begin{tikzpicture}
    \begin{axis}[
      xbar stacked,
      bar width=0.55cm,
      width=11cm, height=6.5cm,
      symbolic y coords={
        {EdgeTPU\\(persistent, tương lai)},
        {CPU\\(persistent, tương lai)},
        {EdgeTPU\\(subprocess hiện tại)},
        {CPU\\(subprocess hiện tại)}
      },
      ytick=data,
      y tick label style={align=right, font=\small},
      xlabel={Thời gian (ms)},
      xlabel style={font=\small},
      xmin=0, xmax=780,
      legend style={at={(0.5,-0.18)}, anchor=north, legend columns=4,
                    font=\scriptsize},
      legend cell align=left,
      grid=major, grid style={dashed, gray!25},
      axis line style={gray!60},
      enlarge y limits=0.18,
      nodes near coords, nodes near coords align={horizontal},
      point meta=explicit symbolic,
      every node near coord/.style={font=\scriptsize},
    ]
    %%  stacks: I/O | Inference | Model load | Python startup
    %% CPU subprocess: 10 + 150 + 100 + 350 = 610
    \addplot[fill=gray!40, draw=gray!60] coordinates {
      (10, {CPU\\(subprocess hiện tại)})
      (10, {EdgeTPU\\(subprocess hiện tại)})
      (10, {CPU\\(persistent, tương lai)})
      (5,  {EdgeTPU\\(persistent, tương lai)})
    };  %% I/O
    \addplot[fill=warnred!55, draw=warnred!70] coordinates {
      (150, {CPU\\(subprocess hiện tại)})
      (10,  {EdgeTPU\\(subprocess hiện tại)})
      (150, {CPU\\(persistent, tương lai)})
      (10,  {EdgeTPU\\(persistent, tương lai)})
    };  %% Inference
    \addplot[fill=orng!65, draw=orng!80] coordinates {
      (100, {CPU\\(subprocess hiện tại)})
      (50,  {EdgeTPU\\(subprocess hiện tại)})
      (0,   {CPU\\(persistent, tương lai)})
      (0,   {EdgeTPU\\(persistent, tương lai)})
    };  %% Model load
    \addplot[fill=cppblue!55, draw=cppblue!70] coordinates {
      (350, {CPU\\(subprocess hiện tại)})
      (350, {EdgeTPU\\(subprocess hiện tại)})
      (0,   {CPU\\(persistent, tương lai)})
      (0,   {EdgeTPU\\(persistent, tương lai)})
    };  %% Python startup
    \legend{I/O (JPEG), Inference, Model load, Python startup}
    \end{axis}
  \end{tikzpicture}
  \caption{Phân tích latency mỗi lần gọi AI detection.
           Python startup time ($\sim$350~ms) chiếm phần lẫn tổng overhead
           trong kiến trúc subprocess hiện tại --- lớn hơn nhiều so với
           thời gian inference thực tế. Kiến trúc persistent process (tương lai)
           loại bỏ hoàn toàn chi phí này.}
  \label{fig:ai3-latency}
\end{figure}

\begin{table}[htbp]
  \centering
  \caption{So sánh tổng latency giữa các cấu hình triển khai}
  \label{tab:ai3-latency}
  \small
  \begin{tabular}{lcc}
    \toprule
    \textbf{Cấu hình}                    & \textbf{Inference} & \textbf{Tổng / lần gọi} \\
    \midrule
    CPU subprocess (hiện tại)            & $\sim$100--200~ms  & $\sim$600--700~ms \\
    EdgeTPU subprocess (hiện tại)       & $\sim$4--15~ms     & $\sim$420--430~ms \\
    CPU persistent (tương lai)           & $\sim$100--200~ms  & $\sim$160--210~ms \\
    EdgeTPU persistent (tương lai)      & $\sim$4--15~ms     & $\sim$15--20~ms  \\
    \bottomrule
  \end{tabular}
\end{table}

Kết quả trên cho thấy rõ điểm nghẻn hiện tại:
EdgeTPU inference chỉ mất 4--15~ms nhưng toàn bộ lá gọi
vẫn tốn 420~ms do phải khởi động Python mỗi lần. Với
kiến trúc \textbf{persistent Python process} --- một process
duy nhất khởi động một lần, load mô hình vào bộ nhớ, sau
đó tiếp nhận request từ C++ qua stdin/stdout pipe --- latency
mỗi lần gọi giảm xuống còn $\sim$15--20~ms với EdgeTPU,
mở ra khả năng chạy AI detection \textit{mỗi frame} camera
thay vì chỉ khi cần.

% Chú thích về nguồn số liệu
\smallskip
\noindent\textit{\small
Lưu ý: COCO metrics trong Bảng~\ref{tab:ai3-metrics} và Hình~\ref{fig:ai3-coco}
lấy từ kết quả đánh giá hai lớp (cylinder + cube) trên repository huấn luyện.
Phân tích confidence score trong Bảng~\ref{tab:ai3-conf} lấy từ kiểm tra
độc lập trên 218 ảnh cylinder và 10 ảnh cube (CPU inference via ai-edge-litert).
}
   %% mục 5.4
%%     ...
%%   \end{document}
%%
%% Compile: pdflatex -> pdflatex (chạy 2 lần để có label/ref đúng)
%%=============================================================

%%   \begin{document}
%%     ...
%%     \chapter{Thiết kế và Hiện thực}
%%     ...
%%     %%=============================================================
%% chapter4_4.tex
%% Phần 4.4 — AI #3: Phát hiện vật thể cho pick & place
%%
%% Đây là \input file — thêm vào preamble của main.tex:
%%   \usepackage{tikz}
%%   \usetikzlibrary{shapes.geometric,arrows.meta,positioning,
%%     fit,backgrounds,calc,decorations.pathreplacing}
%%   \usepackage{pgfplots}
%%   \pgfplotsset{compat=1.18}
%%   \usepackage{listings,xcolor,booktabs,multirow,subcaption}
%%=============================================================

\definecolor{cppblue}{RGB}{30,100,200}
\definecolor{pygreen}{RGB}{40,150,80}
\definecolor{warnred}{RGB}{200,60,60}
\definecolor{orng}{RGB}{220,130,30}

\lstset{
  basicstyle=\ttfamily\small,
  keywordstyle=\color{cppblue},
  commentstyle=\color{gray!70},
  stringstyle=\color{orng},
  breaklines=true,
  frame=single,
  rulecolor=\color{gray!40},
  framexleftmargin=4pt,
  xleftmargin=8pt,
  numbers=none,
}

% ================================================================
\section{AI \#3 --- Tự động hoá khâu khởi tạo trong hệ thống pick \& place}
\label{sec:ai3}
% ================================================================

% ----------------------------------------------------------------
\subsection{Bối cảnh: vấn đề của hệ thống visual servoing ban đầu}
\label{subsec:ai3-context}
% ----------------------------------------------------------------

Hệ thống pick \& place chạy trên cánh tay robot Denso VS-6577E 6 bậc tự do,
sử dụng thư viện ViSP (Visual Servoing Platform) để điều khiển bằng thị giác.
Camera USB được gắn trực tiếp lên đầu công tác theo cấu hình \textit{eye-in-hand} ---
camera di chuyển cùng robot và nhìn thẳng về phía vật thể cần gắp trong suốt
quá trình tiếp cận.

Vấn đề cốt lõi của hệ thống ban đầu nằm ở bước khởi tạo: sau khi robot về
home pose cố định, \textbf{người vận hành phải click chuột vào vị trí cylinder}
trên cửa sổ hiển thị camera thì blob tracker (\texttt{vpDot2}) mới bắt đầu
hoạt động được. Visual servo chỉ chạy sau thao tác thủ công đó. Ngoài ra,
nếu tracker mất target trong lúc servo --- điều hoàn toàn có thể xảy ra khi
cylinder bị che khuất hoặc di chuyển ra ngoài vùng sáng --- hệ thống reset
hoàn toàn và lại đòi hỏi click thủ công ở lần tiếp theo. Mỗi chu kỳ
pick \& place đều cần sự can thiệp của người vận hành, loại bỏ khả năng
vận hành tự động không giám sát.

\begin{figure}[htbp]
  \centering
  \begin{tikzpicture}[
    node distance=0.8cm and 0.6cm,
    bx/.style={draw, rounded corners=4pt, text width=3.0cm, align=center,
               minimum height=0.9cm, font=\small},
    robot/.style={bx, draw=cppblue!70, fill=cppblue!8},
    manual/.style={bx, draw=warnred!70, fill=warnred!8},
    aibox/.style={bx, draw=pygreen!70, fill=pygreen!8},
    arr/.style={->, >=Stealth, thick},
    rarr/.style={arr, draw=warnred!70},
    garr/.style={arr, draw=pygreen!70},
  ]

  %% --- CỘT TRÁI: hệ thống gốc ---
  \node[robot]                    (L1) {PREINIT\\{\scriptsize robot $\to$ home pose}};
  \node[manual, below=of L1]      (L2) {Click thủ công\\{\scriptsize chọn cylinder}};
  \node[robot,  below=of L2]      (L3) {JOINT\\{\scriptsize vpDot2 tracking}};
  \node[manual, below=of L3]      (L4) {Mất tracking\\{\scriptsize $\to$ reset toàn bộ}};

  \draw[arr]  (L1) -- (L2);
  \draw[arr]  (L2) -- (L3);
  \draw[rarr] (L3) -- (L4);
  \draw[rarr, dashed] (L4.west)
        to[out=190,in=170,looseness=1.3] (L1.west)
        node[midway, left, font=\scriptsize, text=warnred!70] {reset};

  \node[font=\bfseries\small, text=warnred, above=0.5cm of L1]
       {\textit{Trước cải tiến}};

  %% --- CỘT PHẢI: hệ thống mới ---
  \node[robot,  right=3.8cm of L1] (R1) {PREINIT\\{\scriptsize robot $\to$ home pose}};
  \node[aibox, below=of R1]        (R2) {AI detect\\{\scriptsize EfficientDet-Lite0}};
  \node[robot,  below=of R2]       (R3) {JOINT\\{\scriptsize AI tracking mỗi frame}};
  \node[aibox, below=of R3]        (R4) {AI re-detect\\{\scriptsize tự phục hồi}};

  \draw[arr]  (R1) -- (R2);
  \draw[garr] (R2) -- (R3);
  \draw[garr] (R3) -- (R4);
  \draw[garr, dashed] (R4.east)
        to[out=0,in=0,looseness=1.5] (R3.east)
        node[midway, right, font=\scriptsize, text=pygreen!70] {tiếp tục};

  %% fallback arrow từ AI detect
  \node[manual, right=1.2cm of R3] (FB)
       {{\scriptsize fallback:}\\click thủ công};
  \draw[rarr, dashed] (R2.east) -| (FB.north)
        node[near start, above, font=\scriptsize, text=warnred] {AI fail};

  \node[font=\bfseries\small, text=pygreen!80!black, above=0.5cm of R1]
       {\textit{Sau cải tiến}};

  %% Đường phân cách
  \draw[thick, dashed, gray!50]
      ($(L1.east)!0.5!(R1.west)+(0,0.8)$) -- ++(0,-6.5);

  \end{tikzpicture}
  \caption{So sánh luồng vận hành trước và sau khi tích hợp module AI.
           Bước yêu cầu thao tác thủ công (đỏ) được thay thế bằng AI tự động (xanh lá).
           Fallback về click thủ công vẫn được giữ lại cho trường hợp AI thất bại.}
  \label{fig:ai3-before-after}
\end{figure}

Hình~\ref{fig:ai3-before-after} minh hoạ sự thay đổi này. Module AI đảm nhiệm
toàn bộ việc phát hiện cylinder ở bước khởi tạo và tự phục hồi khi mất target
trong quá trình servo. Hệ thống vẫn giữ nguyên fallback về click thủ công để
đảm bảo tương thích ngược.


% ----------------------------------------------------------------
\subsection{Xây dựng và huấn luyện mô hình EfficientDet-Lite0}
\label{subsec:ai3-model}
% ----------------------------------------------------------------

\subsubsection*{Lựa chọn kiến trúc}

Yêu cầu đặt ra khá cụ thể: mô hình phải đủ nhỏ để triển khai được trên
Google Coral EdgeTPU, phải hỗ trợ INT8 quantization, và tương thích với
\texttt{tflite\_model\_maker} --- framework đã có sẵn trong môi trường phát triển.
\textbf{EfficientDet-Lite0} đáp ứng tất cả: đây là biến thể nhỏ nhất trong
họ EfficientDet được Google tối ưu cho thiết bị edge, input cố định
$320 \times 320$ pixel, backbone là EfficientNet-B0 đã được simplify.

\begin{table}[htbp]
  \centering
  \caption{Thông số mô hình EfficientDet-Lite0}
  \label{tab:ai3-model}
  \small
  \begin{tabular}{ll}
    \toprule
    \textbf{Thuộc tính}       & \textbf{Giá trị} \\
    \midrule
    Kiến trúc                 & EfficientDet-Lite0 \\
    Framework training        & TensorFlow~2.8.4 / \texttt{tflite\_model\_maker}~0.4.3 \\
    Input size                & $320 \times 320$~px, uint8 \\
    Số lớp đối tượng          & 2 (cylinder~--~lớp~0, cube~--~lớp~1) \\
    Số epochs / batch size    & 50 / 8 \\
    Transfer learning         & Frozen backbone, fine-tune detection head \\
    Target hardware           & Google Coral USB Accelerator (EdgeTPU) \\
    \bottomrule
  \end{tabular}
\end{table}

Mô hình được huấn luyện để phát hiện \textbf{hai lớp}: cylinder và cube nhựa PLA
in 3D. Việc đưa thêm cube vào tập huấn luyện là có chủ ý --- giúp mô hình học
cách phân biệt hai loại vật thể, từ đó giảm false positive khi cả hai cùng
xuất hiện trong cảnh. Ở phía tích hợp với ViSP, chỉ kết quả cylinder
(\texttt{class\_id~==~0}) được dùng; dự đoán cube bị lọc bỏ.

\subsubsection*{Tập dữ liệu}

Toàn bộ ảnh được thu trực tiếp từ camera eye-in-hand gắn trên robot Denso,
đảm bảo điều kiện thực tế sát với lúc triển khai. Ảnh bao gồm nhiều cấu hình:
góc chụp khác nhau, khoảng cách thay đổi, và số lượng vật thể từ 1 đến 5 trong
một cảnh.

\begin{figure}[htbp]
  \centering
  \begin{tikzpicture}
    \begin{axis}[
      xbar, bar width=0.42cm,
      xlabel={Số lượng ảnh (xấp xỉ)},
      ylabel={Danh mục},
      symbolic y coords={
        Background,
        2 cube + 3 cyl,
        2 cylinders,
        3 cylinders,
        1 cyl + 1 cube,
        1 cube,
        1 cylinder,
        Cylinder (gốc)
      },
      ytick=data,
      xmin=0, xmax=175,
      width=9.5cm, height=6.8cm,
      nodes near coords,
      nodes near coords align={horizontal},
      every node near coord/.style={font=\scriptsize},
      tick label style={font=\small},
      xlabel style={font=\small},
      ylabel style={font=\small},
      enlarge y limits=0.08,
    ]
    \addplot[fill=cppblue!55, draw=cppblue!80] coordinates {
      (38,  Background)
      (28,  2 cube + 3 cyl)
      (52,  2 cylinders)
      (43,  3 cylinders)
      (68,  1 cyl + 1 cube)
      (38,  1 cube)
      (108, 1 cylinder)
      (144, Cylinder (gốc))
    };
    \end{axis}
  \end{tikzpicture}
  \caption{Phân bố tập dữ liệu theo danh mục ảnh (~534 ảnh tổng,
           633 annotation cylinder, 170 annotation cube).
           Các cảnh đa vật thể giúp mô hình học phát hiện từng đối tượng
           riêng lẻ trong môi trường phức tạp.}
  \label{fig:ai3-dataset}
\end{figure}

Annotation được thực hiện bằng \texttt{labelme} theo format bounding box hình chữ
nhật. Workflow:
\begin{enumerate}
  \item Vẽ bounding box cho từng cylinder và cube bằng \texttt{labelme}
  \item Chuyển đổi labelme JSON $\to$ COCO JSON thống nhất
        (script \texttt{labelme\_to\_coco.py})
  \item Tập dữ liệu được split theo tỉ lệ $70/15/15$ (train/val/test),
        seed cố định để đảm bảo tính tái lập
\end{enumerate}

Phân chia ba chiều thay vì hai chiều đảm bảo tập test hoàn toàn độc lập ---
không được nhìn thấy trong suốt quá trình training hoặc điều chỉnh siêu tham số.

\subsubsection*{Pipeline huấn luyện}

\begin{figure}[htbp]
  \centering
  \begin{tikzpicture}[
    node distance=0.45cm and 0.15cm,
    st/.style={draw, rounded corners=3pt, fill=gray!10, draw=gray!55,
               text width=2.0cm, align=center, minimum height=0.85cm, font=\small},
    blue_st/.style={st, fill=cppblue!10, draw=cppblue!55},
    out_st/.style={st, fill=orng!20, draw=orng!70},
    arr/.style={->, >=Stealth, thick, gray!60},
    lbl/.style={font=\scriptsize, text=gray!80},
  ]

  \node[blue_st] (s1) {Annotation\\{\scriptsize labelme}};
  \node[blue_st, right=of s1] (s2) {COCO\\JSON};
  \node[blue_st, right=of s2] (s3) {Pascal\\VOC XML};
  \node[blue_st, right=of s3] (s4) {Training\\{\scriptsize 50 epochs}};
  \node[out_st, right=of s4] (s5) {SavedModel\\{\scriptsize Float32}};
  \node[out_st, below=0.55cm of s5] (s6) {INT8\\TFLite};
  \node[out_st, below=0.55cm of s6] (s7) {EdgeTPU\\TFLite};

  \draw[arr] (s1) -- (s2) node[lbl,above,midway] {to\_coco.py};
  \draw[arr] (s2) -- (s3) node[lbl,above,midway] {tflite\_mm};
  \draw[arr] (s3) -- (s4);
  \draw[arr] (s4) -- (s5) node[lbl,above,midway] {export};
  \draw[arr] (s5) -- (s6) node[lbl,right,midway] {PTQ};
  \draw[arr] (s6) -- (s7) node[lbl,right,midway] {edgetpu\_compiler};

  \end{tikzpicture}
  \caption{Pipeline huấn luyện và xuất mô hình. Từ ảnh thô đến ba file
           mô hình phục vụ các môi trường triển khai khác nhau.
           Bước chuyển đổi sang Pascal VOC XML là yêu cầu bắt buộc của
           \texttt{tflite\_model\_maker}.}
  \label{fig:ai3-pipeline}
\end{figure}

Transfer learning với \texttt{train\_whole\_model=False} --- chỉ fine-tune detection
head, giữ nguyên backbone đã pretrain --- phù hợp với tập dữ liệu tương đối nhỏ
($\sim$374 ảnh train) và giảm nguy cơ overfitting đáng kể so với
huấn luyện từ đầu.


% ----------------------------------------------------------------
\subsection{Quantization INT8 và biên dịch cho EdgeTPU}
\label{subsec:ai3-quant}
% ----------------------------------------------------------------

Mô hình sau training là Float32 --- không chạy được trực tiếp trên Coral EdgeTPU.
EdgeTPU chỉ thực thi mô hình INT8 quantized, trong đó trọng số và activation
được biểu diễn bằng số nguyên 8-bit. Lợi ích: kích thước giảm $\sim4\times$,
và phần cứng EdgeTPU đạt tốc độ inference cực nhanh ($\sim$4--15~ms/frame)
nhờ tận dụng phép tính số nguyên.

Post-Training Quantization (PTQ) sử dụng tập train làm \textit{representative
dataset} để calibrate scale và zero\_point cho từng layer --- đây là bước bắt
buộc để duy trì độ chính xác của INT8 model gần với Float32 baseline.

File \texttt{\_edgetpu.tflite} sau khi biên dịch chứa custom op
\texttt{edgetpu-custom-op} mà TFLite runtime thông thường
\textit{không thể thực thi trên CPU}, gây lỗi:

\begin{lstlisting}[language={}]
RuntimeError: Encountered unresolved custom op: edgetpu-custom-op
\end{lstlisting}

Do đó hệ thống \textbf{duy trì hai file riêng biệt} và lựa chọn tự động:

\begin{center}
\small
\begin{tabular}{lll}
  \toprule
  File & Runtime & Deployment \\
  \midrule
  \texttt{cylinder\_detector\_int8\_edgetpu.tflite} & pycoral   & Coral EdgeTPU \\
  \texttt{cylinder\_detector\_int8.tflite}          & TFLite/ai-edge-litert & CPU \\
  \bottomrule
\end{tabular}
\end{center}


% ----------------------------------------------------------------
\subsection{Kiến trúc tích hợp: C++ ViSP ↔ Python subprocess}
\label{subsec:ai3-arch}
% ----------------------------------------------------------------

\subsubsection*{Quyết định thiết kế}

Thay vì link trực tiếp thư viện TFLite C++ vào build system của ViSP
(đòi hỏi thay đổi lớn vào CMakeLists và phụ thuộc version build cụ thể),
hệ thống áp dụng kiến trúc \textbf{subprocess model}: C++ gọi Python script
như một tiến trình con qua \texttt{popen()}, giao tiếp một chiều qua stdout.
Cách này giúp module AI hoàn toàn tách biệt --- cập nhật mô hình không cần
recompile code robot.

\begin{figure}[htbp]
  \centering
  \begin{tikzpicture}[
    node distance=0.38cm and 0.5cm,
    comp/.style={draw, rounded corners=4pt, text width=3.1cm, align=center,
                 minimum height=0.82cm, font=\small},
    cpp/.style={comp, fill=cppblue!10, draw=cppblue!55},
    py/.style={comp,  fill=pygreen!10, draw=pygreen!60},
    tmp/.style={comp,  fill=orng!18, draw=orng!65, text width=2.5cm,
                minimum height=0.7cm, font=\scriptsize},
    arr/.style={->, >=Stealth, thick},
    ]

  %% ---- C++ side ----
  \node[cpp] (cap)   {Camera capture\\{\scriptsize vpImage{$<$}uchar{$>$}}};
  \node[cpp, below=of cap]  (save)  {Save frame\\{\scriptsize cv::imwrite()}};
  \node[cpp, below=of save] (pop)   {popen() subprocess\\{\scriptsize build cmd}};
  \node[cpp, below=of pop]  (sel)   {select() timeout 3s\\{\scriptsize POSIX wait}};
  \node[cpp, below=of sel]  (parse) {Parse stdout\\{\scriptsize sscanf SUCCESS}};
  \node[cpp, below=of parse](feat)  {vpFeatureBuilder\\{\scriptsize create feature}};

  \draw[arr](cap)--(save); \draw[arr](save)--(pop);
  \draw[arr](pop)--(sel);  \draw[arr](sel)--(parse); \draw[arr](parse)--(feat);

  \node[draw=cppblue!50, dashed, rounded corners=6pt,
        fit=(cap)(save)(pop)(sel)(parse)(feat),
        inner sep=7pt] (cbox) {};
  \node[above=2pt of cbox, font=\bfseries\small, text=cppblue]
       {C++ (ViSP main loop)};

  %% ---- Temp file ----
  \node[tmp, right=2.6cm of save] (jpg)
       {/tmp/visp\_ai\_frame.jpg};

  %% ---- Python side ----
  \node[py, right=2.6cm of pop]   (load) {Load model\\{\scriptsize fallback chain}};
  \node[py, below=of load]        (pre)  {Preprocess\\{\scriptsize resize 320$\times$320}};
  \node[py, below=of pre]         (inf)  {Inference\\{\scriptsize invoke()}};
  \node[py, below=of inf]         (post) {Post-process\\{\scriptsize dequant+filter}};
  \node[py, below=of post]        (out)  {Print stdout\\{\scriptsize SUCCESS u v conf}};

  \draw[arr](load)--(pre); \draw[arr](pre)--(inf);
  \draw[arr](inf)--(post); \draw[arr](post)--(out);

  \node[draw=pygreen!55, dashed, rounded corners=6pt,
        fit=(load)(pre)(inf)(post)(out),
        inner sep=7pt] (pbox) {};
  \node[above=2pt of pbox, font=\bfseries\small, text=pygreen!80!black]
       {Python subprocess};

  %% ---- Data flow arrows ----
  \draw[arr, orng!80]  (save.east) -- (jpg.west)  node[midway,above,font=\scriptsize]{write};
  \draw[arr, orng!80]  (jpg.east)  -- (load.west) node[midway,above,font=\scriptsize]{read};
  \draw[arr, cppblue!60](pop.east)  -- (load.west) node[midway,above,font=\scriptsize,yshift=2pt]{spawn};
  \draw[arr, pygreen!70, bend right=25](out.west)  to (parse.east)
        node[midway,below,font=\scriptsize]{stdout pipe};

  \end{tikzpicture}
  \caption{Kiến trúc tích hợp C++ ViSP $\leftrightarrow$ Python subprocess.
           C++ lưu frame hiện tại ra file tạm, spawn Python làm child process
           qua \texttt{popen()}, đọc kết quả từ stdout sau khi qua timeout guard.}
  \label{fig:ai3-arch}
\end{figure}

\subsubsection*{Protocol trao đổi dữ liệu}

Output của Python script là một dòng duy nhất trên stdout:

\begin{lstlisting}[language={}]
SUCCESS 272.86 190.74 0.8362 238.42 166.89 68.88 47.68
%        u      v      conf   x_min  y_min  width  height
\end{lstlisting}

hoặc \texttt{FAILURE~<lý~do>} nếu không phát hiện được. Tất cả debug log
đi ra \texttt{stderr} --- C++ chỉ đọc stdout --- tránh nhiễu kết quả parse.

Đáng chú ý: trong ViSP, \texttt{vpImagePoint} được định nghĩa là \texttt{(row,~col)}
tương ứng \texttt{(v,~u)}, không phải \texttt{(u,~v)} như thông thường:

\begin{lstlisting}[language=C++]
detected_center = vpImagePoint(v, u);  // row=v, col=u
\end{lstlisting}

Tham số \texttt{hint} cho phép Python chọn cylinder \textit{gần nhất với vị
trí đã biết} thay vì cylinder có confidence cao nhất --- cơ chế này tránh
hiện tượng servo ``nhảy'' sang cylinder khác khi scores thay đổi giữa các frame.


% ----------------------------------------------------------------
\subsection{State machine: PREINIT \texorpdfstring{$\to$}{→} INIT (AI) \texorpdfstring{$\to$}{→} JOINT \texorpdfstring{$\to$}{→} ...}
\label{subsec:ai3-sm}
% ----------------------------------------------------------------

Hệ thống vận hành theo vòng lặp chính liên tục với 6 trạng thái.
Mỗi iteration đọc một frame camera và thực thi logic trạng thái hiện tại.

\begin{figure}[htbp]
  \centering
  \begin{tikzpicture}[
    node distance=1.0cm and 2.0cm,
    st/.style={draw, rounded corners=4pt, fill=cppblue!10, draw=cppblue!55,
               text width=2.3cm, align=center, minimum height=1.0cm, font=\small},
    ai/.style={st, fill=pygreen!10, draw=pygreen!60},
    arr/.style={->, >=Stealth, thick},
    lbl/.style={font=\scriptsize, text=gray!70},
  ]

  %% Top row
  \node[st]  (pre)  {PREINIT\\{\scriptsize home pose}};
  \node[ai,  right=of pre]   (ini)  {INIT\\{\scriptsize AI detect}};
  \node[ai,  right=of ini]   (jnt)  {JOINT\\{\scriptsize AI tracking}};
  \node[st,  right=of jnt]   (app)  {APPROACH\\{\scriptsize chờ RC5}};

  %% Bottom row
  \node[st,  below=1.8cm of app]  (grp)  {GRIPPER\\{\scriptsize đóng gripper}};
  \node[st,  left=of grp]         (cls)  {CLASSIFIED\\{\scriptsize vận chuyển}};

  %% Arrows
  \draw[arr](pre)--(ini)  node[lbl,above,midway]{home\\reached};
  \draw[arr](ini)--(jnt)  node[lbl,above,midway]{AI init\\ok};
  \draw[arr](jnt)--(app)  node[lbl,above,midway]{$|e|<5{\times}10^{-3}$};
  \draw[arr](app)--(grp)  node[lbl,right,midway]{RC5\\signal};
  \draw[arr](grp)--(cls)  node[lbl,below,midway]{grasped};
  \draw[arr](cls.west) to[out=180,in=270,looseness=0.8] (pre.south)
        node[lbl,left,pos=0.5]{cycle\\restart};

  %% AI self-recovery
  \draw[arr, pygreen!70, dashed]
        (jnt.north) to[out=100,in=80,looseness=2.5] (jnt.north)
        node[above=0.7cm of jnt, font=\scriptsize, text=pygreen!80!black]
        {AI re-detect};

  %% Fallback from INIT
  \draw[arr, warnred!70, dashed]
        (ini.south) -- ++(0,-0.7)
        node[below, font=\scriptsize, text=warnred]
        {fallback: click thủ công};

  \end{tikzpicture}
  \caption{State machine của hệ thống pick \& place sau cải tiến.
           Trạng thái tô xanh lá là nơi AI được gọi.
           Vòng lặp tự phục hồi ở JOINT (mũi tên đứt nét) cho phép hệ thống
           tiếp tục mà không cần reset khi mất detection.}
  \label{fig:ai3-sm}
\end{figure}

\paragraph{PREINIT.}
Robot gửi lệnh joint position về home pose
$[0^\circ, 0^\circ, 90^\circ, 0^\circ, 90^\circ, 0^\circ]$.
Servo task được cấu hình với interaction matrix Desired, pseudo-inverse,
gain $\lambda = 0.4$.

\paragraph{INIT.}
Chờ robot đạt home pose (sai số từng joint $< 0.01$) và gripper mở xong.
Khi cả hai điều kiện thỏa, \texttt{detectCylinderWithAI()} được gọi
không có hint --- Python chọn cylinder có confidence cao nhất.
Kết quả trực tiếp tạo visual feature \texttt{vpFeaturePoint}
và add vào servo task. Nếu AI thất bại, \texttt{ai\_hint}
giữ giá trị mặc định (fallback thủ công).

\paragraph{JOINT.}
Khác với cách tiếp cận truyền thống dùng blob tracker mỗi frame,
hệ thống hiện tại \textbf{gọi AI trực tiếp ở mỗi iteration}
với \texttt{hint = \&ai\_hint}. Control law:
\begin{equation}
  \mathbf{v} = -\lambda\,\mathbf{L}_{s^*}^{+}\,\mathbf{e},
  \qquad \mathbf{q}_{k+1} = \mathbf{q}_k + \mathbf{v}\,\Delta t
  \label{eq:ibvs}
\end{equation}
với $\Delta t = 0.5\,$s. Điều kiện hội tụ: $|e_x| < 5{\times}10^{-3}$
và $|e_y| < 5{\times}10^{-3}$ (normalized image coordinates).
Khi đạt, gửi \texttt{"OKE\textbackslash r"} qua UART đến RC5 controller.


% ----------------------------------------------------------------
\subsection{Các vấn đề kỹ thuật và giải pháp}
\label{subsec:ai3-issues}
% ----------------------------------------------------------------

Một số vấn đề nảy sinh trong quá trình tích hợp mà khó thấy trước khi thực
sự chạy hệ thống.

\subsubsection*{P1: Xác định output tensor sai lệch}

\texttt{tflite\_model\_maker} 0.4.x đặt tên tensor theo thứ tự riêng của nó:

\begin{center}
\small
\begin{tabular}{ccc}
  \toprule
  Tensor & Shape & Nội dung \\
  \midrule
  \texttt{:0} & \texttt{[1]}     & num\_detections \\
  \texttt{:1} & \texttt{[1,\,N]} & scores \quad \textbf{← cần lấy cái này} \\
  \texttt{:2} & \texttt{[1,\,N]} & class IDs \quad \textbf{← không phải cái này} \\
  \texttt{:3} & \texttt{[1,\,N,\,4]} & boxes \\
  \bottomrule
\end{tabular}
\end{center}

Thứ tự này không khớp với API của pycoral, vốn giả định positional order
\texttt{boxes/classes/scores/count = 0/1/2/3}. Dùng API đó sẽ đọc
num\_detections (raw uint8 = 255) như scores và crash với \texttt{IndexError}.
Giải pháp: xác định tensor \textbf{theo shape động} --- boxes là
\texttt{[1,N,4]}, scores là tensor \texttt{[1,N]} có name suffix nhỏ nhất
(\texttt{:1} chứ không phải \texttt{:2}).

\subsubsection*{P2: Python 3.10 không có \texttt{tflite-runtime}}

Google không publish wheel \texttt{tflite-runtime} cho Python 3.10+.
Giải pháp là xây dựng \textbf{fallback chain 4 tầng}:

\begin{figure}[htbp]
  \centering
  \begin{tikzpicture}[
    node distance=0.4cm,
    fb/.style={draw, rounded corners=3pt, text width=7.5cm, align=left,
               minimum height=0.85cm, font=\small, inner sep=7pt},
    t1/.style={fb, fill=purple!12, draw=purple!55},
    t2/.style={fb, fill=teal!10,   draw=teal!50},
    t3/.style={fb, fill=cppblue!10, draw=cppblue!50},
    t4/.style={fb, fill=gray!12,   draw=gray!50},
    fa/.style={fb, fill=warnred!10, draw=warnred!50},
    arr/.style={->, >=Stealth, thick, dashed, gray!55},
    lbl/.style={font=\scriptsize, text=warnred!80},
  ]
  \node[t1] (a) {\textbf{1.}~\texttt{pycoral.utils.edgetpu}\quad{\scriptsize Coral EdgeTPU, Python 3.9}};
  \node[t2, below=of a] (b) {\textbf{2.}~\texttt{ai\_edge\_litert.interpreter}\quad{\scriptsize Google's new pkg, Python 3.8--3.12}};
  \node[t3, below=of b] (c) {\textbf{3.}~\texttt{tflite\_runtime.interpreter}\quad{\scriptsize Legacy, Python 3.9}};
  \node[t4, below=of c] (d) {\textbf{4.}~\texttt{tensorflow.lite.Interpreter}\quad{\scriptsize Heavy fallback}};
  \node[fa, below=of d] (e) {\textbf{FAILURE}:~\texttt{no\_tflite\_backend\_available}};

  \draw[arr](a)--(b) node[lbl,right,midway]{ImportError / Coral not connected};
  \draw[arr](b)--(c) node[lbl,right,midway]{ImportError};
  \draw[arr](c)--(d) node[lbl,right,midway]{ImportError};
  \draw[arr](d)--(e) node[lbl,right,midway]{ImportError};
  \end{tikzpicture}
  \caption{Python backend fallback chain trong \texttt{detect\_cylinder.py}.
           Script tự động thử từng backend theo thứ tự ưu tiên,
           tương thích với cả môi trường có Coral lẫn CPU-only.}
  \label{fig:ai3-fallback}
\end{figure}

\texttt{ai-edge-litert} là gói chính thức của Google thay thế \texttt{tflite-runtime},
hỗ trợ Python 3.8--3.12 với API giống hệt nhau --- là giải pháp cho vấn đề Python 3.10.

\subsubsection*{P3: JSON comment trong config}

\texttt{config.json} dùng comment \texttt{//} để ghi chú, nhưng cả Python
\texttt{json.load()} và \texttt{nlohmann::json} đều từ chối mặc định.
Phía C++: \texttt{nlohmann::json::parse(file, nullptr, true, true)} với
\texttt{ignore\_comments=true} (có từ nlohmann/json~3.9.0).
Phía Python: strip comment lines thủ công trước khi \texttt{json.loads()}.


% ----------------------------------------------------------------
\subsection{Thông số cấu hình visual servo}
\label{subsec:ai3-servo}
% ----------------------------------------------------------------

Hệ thống sử dụng \textbf{Image-Based Visual Servoing (IBVS)} với feature là
tọa độ 2D của cylinder trong không gian ảnh normalized $\mathbf{s}=(x,y)$,
desired position $\mathbf{s}^*=(0,0,1)$ --- cylinder phải nằm đúng tâm ảnh.

\begin{table}[htbp]
  \centering
  \caption{Thông số điều khiển visual servo}
  \label{tab:ai3-servo}
  \small
  \begin{tabular}{ll}
    \toprule
    \textbf{Tham số}            & \textbf{Giá trị} \\
    \midrule
    Loại điều khiển             & Eye-in-hand, \texttt{EYEINHAND\_L\_cVe\_eJe} \\
    Interaction matrix          & Desired ($\mathbf{L}_{s^*}$), pseudo-inverse \\
    Gain $\lambda$              & 0.4 \\
    Bước thời gian $\Delta t$   & 0.5\,s \\
    Ngưỡng hội tụ               & $|e_x| < 5{\times}10^{-3}$ và $|e_y| < 5{\times}10^{-3}$ \\
    Desired position            & $(x^*, y^*) = (0,\,0)$ --- tâm ảnh, $Z = 1$ \\
    \bottomrule
  \end{tabular}
\end{table}

Dùng $\mathbf{L}_{s^*}$ (tại vị trí mong muốn) thay vì $\mathbf{L}_{s_{cur}}$
giúp tránh singularity khi feature tiến gần desired point và giữ ổn định
theo ổn định hơn trong giai đoạn hội tụ. Velocity được saturate trước khi
gửi xuống robot để đảm bảo không vượt giới hạn vật lý từng khớp.
   %% mục 4.4
%%     ...
%%     \chapter{Kết quả và Đánh giá}
%%     ...
%%     %%=============================================================
%% chapter5_4.tex
%% Phần 5.4 — Kết quả đánh giá: AI #3 Phát hiện vật thể
%%
%% \input file — cần thêm vào preamble:
%%   \usepackage{tikz}
%%   \usetikzlibrary{arrows.meta,positioning,patterns}
%%   \usepackage{pgfplots}
%%   \pgfplotsset{compat=1.18}
%%   \usepackage{booktabs,multirow,subcaption,xcolor}
%%=============================================================

% ================================================================
\section{Kết quả và Đánh giá: AI \#3 --- Phát hiện vật thể}
\label{sec:ai3-results}
% ================================================================

% ----------------------------------------------------------------
\subsection{Kết quả detection model: COCO metrics}
\label{subsec:ai3-coco}
% ----------------------------------------------------------------

Mô hình được đánh giá trên cả validation set và test set độc lập
theo chuẩn \textbf{COCO evaluation protocol} --- tính trung bình mAP qua
các ngưỡng IoU từ 0.50 đến 0.95 với bước 0.05.

\begin{figure}[htbp]
  \centering
  \begin{tikzpicture}
    \begin{axis}[
      ybar, bar width=0.45cm,
      width=10cm, height=7cm,
      symbolic x coords={mAP (0.50:0.95), AP50, AP75, AR},
      xtick=data,
      x tick label style={font=\small, align=center},
      ymin=0, ymax=1.15,
      ylabel={Score},
      ylabel style={font=\small},
      legend style={at={(0.5,1.03)}, anchor=south, legend columns=2, font=\small},
      legend cell align=left,
      nodes near coords,
      every node near coord/.style={font=\scriptsize, /pgf/number format/fixed,
                                    /pgf/number format/precision=3},
      enlarge x limits=0.25,
      grid=major, grid style={dashed, gray!30},
      axis line style={gray!60},
      tick style={gray!60},
    ]
    \addplot[fill=cppblue!70, draw=cppblue!90] coordinates {
      (mAP (0.50:0.95), 0.791)
      (AP50, 0.990)
      (AP75, 0.990)
      (AR,   0.869)
    };
    \addplot[fill=orng!80, draw=orng!90] coordinates {
      (mAP (0.50:0.95), 0.799)
      (AP50, 1.000)
      (AP75, 0.969)
      (AR,   0.856)
    };
    \legend{Validation, Test}
    \end{axis}
  \end{tikzpicture}
  \caption{COCO metrics trên validation set và test set. AP50 trên test set
           đạt 1.000 cho thấy 100\% cylinder được phát hiện đúng ở ngưỡng IoU
           lỏng lẻ (50\%). mAP$\approx$0.80 phản ánh vẫn còn đây đó những trường hợp
           định vị bounding box chưa được chính xác ở ngưỡng IoU cao.}
  \label{fig:ai3-coco}
\end{figure}

Bảng~\ref{tab:ai3-metrics} tổng hợp số liệu chi tiết:

\begin{table}[htbp]
  \centering
  \caption{COCO metrics trên validation và test set}
  \label{tab:ai3-metrics}
  \small
  \begin{tabular}{lcc}
    \toprule
    \textbf{Metric}                    & \textbf{Validation} & \textbf{Test} \\
    \midrule
    mAP (IoU 0.50:0.95)               & 0.791               & \textbf{0.799} \\
    AP50 (IoU $\geq$ 0.50)            & 0.990               & \textbf{1.000} \\
    AP75 (IoU $\geq$ 0.75)            & 0.990               & 0.969 \\
    AR (maxDets=100)                   & 0.869               & 0.856 \\
    \midrule
    Precision (IoU=0.50, conf=0.50)    & ---                 & \textbf{1.000} \\
    Recall    (IoU=0.50, conf=0.50)    & ---                 & 0.990 \\
    F1-Score                           & ---                 & \textbf{0.995} \\
    TP / FP / FN                       & ---                 & 95 / 0 / 1 \\
    \bottomrule
  \end{tabular}
\end{table}

\begin{figure}[htbp]
  \centering
  \begin{tikzpicture}
    \begin{axis}[
      ybar, bar width=0.55cm,
      width=8.5cm, height=6cm,
      symbolic x coords={Precision, Recall, F1-Score},
      xtick=data,
      x tick label style={font=\small},
      ymin=0.96, ymax=1.04,
      ylabel={Score},
      ylabel style={font=\small},
      nodes near coords,
      every node near coord/.style={font=\small, /pgf/number format/fixed,
                                    /pgf/number format/precision=3},
      enlarge x limits=0.3,
      grid=major, grid style={dashed, gray!30},
      axis line style={gray!60},
    ]
    \addplot[fill=pygreen!65, draw=pygreen!80] coordinates {
      (Precision, 1.000)
      (Recall,    0.990)
      (F1-Score,  0.995)
    };
    \end{axis}
  \end{tikzpicture}
  \caption{Precision, Recall và F1-Score trên test set (PASCAL VOC protocol,
           IoU$\geq$0.50, confidence$\geq$0.50, chỉ lớp cylinder).
           Precision = 1.000 và FP = 0 chứng tỏ không có false positive
           nào ở ngưỡng 0.50. FN = 1 (một cylinder bị bỏ sót) giải thích
           Recall = 0.990.}
  \label{fig:ai3-prf}
\end{figure}

Giá trị AP50 = 1.000 trên test set à ngưỡng IoU lỏng lẻ (50\%) phản ánh
mô hình phát hiện được toàn bộ cylinder. mAP = 0.799 (trung bình
qua IoU 0.50--0.95) thấp hơn một chút do mô hình EfficientDet-Lite0
đơn giản vẫn còn giới hạn trong việc định vị bounding box chính xác
ở các ngưỡng IoU cao ($\geq$0.75) --- điều phổ biến với các kiến trúc
nhe hơn.


% ----------------------------------------------------------------
\subsection{Phân tích confidence score: true positive vs false positive}
\label{subsec:ai3-conf}
% ----------------------------------------------------------------

Một trong những điều quan trọng khi triển khai thực tế là khoảng cách
phân tách giữa true positive và false positive --- điều này quyết định
ngưỡng confidence nào thực tế dùng được mà không sai kết quả.

Kiểm tra trên 218 ảnh cylinder (true positives) và 10 ảnh cube
(false positives):

\begin{figure}[htbp]
  \centering
  \begin{tikzpicture}
    \begin{axis}[
      width=11cm, height=5.5cm,
      xlabel={Confidence score},
      ylabel={},
      xmin=0, xmax=1.05,
      ymin=0, ymax=3,
      xtick={0,0.1,...,1.0},
      ytick={0.7, 2.0},
      yticklabels={\small False Positive\\(cube, $n$=10), \small True Positive\\(cylinder, $n$=218)},
      y tick label style={align=right, font=\small},
      x tick label style={font=\small},
      xlabel style={font=\small},
      grid=both, grid style={dashed, gray!25},
      axis line style={gray!60},
    ]

    %% --- True Positive band (y=2) ---
    %% Most scores: 0.85-0.98 (majority)
    \addplot[only marks, mark=|, mark size=5pt, mark options={draw=pygreen!70, line width=1.2pt},
             opacity=0.5]
    coordinates {
      (0.29,2)(0.61,2)(0.74,2)
      (0.82,2)(0.83,2)(0.84,2)(0.85,2)(0.86,2)(0.87,2)(0.88,2)
      (0.89,2)(0.90,2)(0.91,2)(0.92,2)(0.93,2)(0.94,2)(0.94,2)
      (0.95,2)(0.95,2)(0.96,2)(0.96,2)(0.97,2)(0.97,2)(0.98,2)
    };

    %% --- False Positive band (y=0.7) ---
    \addplot[only marks, mark=|, mark size=5pt, mark options={draw=warnred!70, line width=1.2pt},
             opacity=0.6]
    coordinates {
      (0.08,0.7)(0.13,0.7)(0.17,0.7)(0.21,0.7)
      (0.25,0.7)(0.28,0.7)(0.30,0.7)(0.33,0.7)(0.35,0.7)(0.36,0.7)
    };

    %% --- Threshold line at 0.50 ---
    \addplot[thick, dashed, draw=gray!80, domain=0:3]
    coordinates {(0.50, 0) (0.50, 3)};
    \node[font=\scriptsize, text=gray!80, rotate=90]
         at (axis cs: 0.52, 1.5) {threshold = 0.50};

    %% --- Gap annotation ---
    \draw[->, >=Stealth, thick, gray!70]
         (axis cs:0.38, 2.5) -- (axis cs:0.48, 2.5);
    \draw[->, >=Stealth, thick, gray!70]
         (axis cs:0.34, 2.5) -- (axis cs:0.27, 2.5);
    \node[font=\scriptsize, text=gray!80] at (axis cs:0.37, 2.7) {gap = 0.34};

    %% --- Median annotation for TP ---
    \addplot[only marks, mark=diamond*, mark size=4pt,
             mark options={draw=pygreen!80!black, fill=pygreen!60}]
    coordinates {(0.94, 2)};
    \node[font=\scriptsize, text=pygreen!80!black, above]
         at (axis cs:0.94, 2) {median 0.94};

    %% --- Max FP annotation ---
    \addplot[only marks, mark=diamond*, mark size=4pt,
             mark options={draw=warnred!80, fill=warnred!50}]
    coordinates {(0.36, 0.7)};
    \node[font=\scriptsize, text=warnred, below]
         at (axis cs:0.36, 0.7) {max 0.36};

    \end{axis}
  \end{tikzpicture}
  \caption{Phân tách confidence score giữa true positive (cylinder, xanh lá)
           và false positive (cube, đỏ). Mỗi vạch đứng đại diện một detection.
           Ngưỡng 0.50 nằm chính giữa khoảng gap 0.34,
           loại bỏ toàn bộ false positive mà vẫn chấp nhận 99.5\% true positive.}
  \label{fig:ai3-conf}
\end{figure}

Khoảng cách 0.34 giữa điểm false positive cao nhất (0.36) và 99.1\%
true positive ($\geq$0.70) xác nhận mô hình đã học được đặc trưng
phân biệt mạnh giữa cylinder và cube. Ngưỡng hiện tại 0.50 nằm an toàn
giữa hai vùng phân bố này.

Ba ảnh outlier có score thấp (0.29, 0.61, 0.74) tương ứng với cylinder
bị nhìn từ góc bất thường hoặc bị che khuất một phần. Ỡ runtime,
các trường hợp này fallback về click thủ công --- hành vi an toàn.

\begin{table}[htbp]
  \centering
  \caption{Phân tích confidence score trên 218 ảnh cylinder + 10 ảnh cube}
  \label{tab:ai3-conf}
  \small
  \begin{tabular}{lcc}
    \toprule
    & \textbf{True Positive (cylinder)} & \textbf{False Positive (cube)} \\
    \midrule
    Score range      & 0.29 -- 0.98          & 0.08 -- 0.36 \\
    Median           & \textbf{0.94}         & --- \\
    $\geq$ 0.50      & 99.5\% (217/218)      & \textbf{0/10 (0\%)} \\
    $\geq$ 0.70      & 99.1\% (216/218)      & 0/10 (0\%) \\
    $\geq$ 0.80      & 95.9\% (209/218)      & 0/10 (0\%) \\
    \midrule
    Gap ph\^an t\'ach & \multicolumn{2}{c}{$0.70 - 0.36 = 0.34$} \\
    \bottomrule
  \end{tabular}
\end{table}


% ----------------------------------------------------------------
\subsection{Kết quả tích hợp hệ thống}
\label{subsec:ai3-integration}
% ----------------------------------------------------------------

\begin{table}[htbp]
  \centering
  \caption{So sánh khả năng trước và sau khi tích hợp AI}
  \label{tab:ai3-comparison}
  \small
  \begin{tabular}{lll}
    \toprule
    \textbf{Tiêu chí}                    & \textbf{Hệ thống gốc} & \textbf{Sau cải tiến} \\
    \midrule
    Khởi tạo mỗi chu kỳ               & Click thủ công (bắt buộc) & Tự động (AI) \\
    Tỷ lệ init thành công             & 100\% (do manual)  & $\sim$99.5\% (AI) \\
    Fallback khi AI thất bại           & Không có              & Click thủ công \\
    Phục hồi khi mất tracking          & Reset toàn bộ        & AI re-detect tự động \\
    Chạy không giám sát                & Không thể              & Có (trường hợp bình thường) \\
    \bottomrule
  \end{tabular}
\end{table}

Thiết kế hiện tại gọi AI trực tiếp \textbf{mỗi iteration} của servo
loop thay vì dùng blob tracker. Điều này có một hàm ý đáng chú ý:
\textit{không có tracking failure} theo nghĩa truyền thống --- mỗi frame
là một lần detection độc lập. Tính ổn định của servo phụ thuộc vào
sự nhất quán của AI detection qua các frame liên tiếp. Cơ chế \texttt{--hint}
(truyền vị trí frame trước) giúp được phần nào khi có nhiều cylinder
trong cảnh.


% ----------------------------------------------------------------
\subsection{Phân tích latency: subprocess CPU vs EdgeTPU}
\label{subsec:ai3-latency}
% ----------------------------------------------------------------

Mỗi lần gọi \texttt{detectCylinderWithAI()} bao gồm nhiều thành phần
latency khác nhau. Điều ít người để ý là: phần overhead lớn nhất
không phải là inference mà là \textbf{Python startup time} --- chi phí
mở một process mới mỗi lần gọi \texttt{popen()}.

\begin{figure}[htbp]
  \centering
  \begin{tikzpicture}
    \begin{axis}[
      xbar stacked,
      bar width=0.55cm,
      width=11cm, height=6.5cm,
      symbolic y coords={
        {EdgeTPU\\(persistent, tương lai)},
        {CPU\\(persistent, tương lai)},
        {EdgeTPU\\(subprocess hiện tại)},
        {CPU\\(subprocess hiện tại)}
      },
      ytick=data,
      y tick label style={align=right, font=\small},
      xlabel={Thời gian (ms)},
      xlabel style={font=\small},
      xmin=0, xmax=780,
      legend style={at={(0.5,-0.18)}, anchor=north, legend columns=4,
                    font=\scriptsize},
      legend cell align=left,
      grid=major, grid style={dashed, gray!25},
      axis line style={gray!60},
      enlarge y limits=0.18,
      nodes near coords, nodes near coords align={horizontal},
      point meta=explicit symbolic,
      every node near coord/.style={font=\scriptsize},
    ]
    %%  stacks: I/O | Inference | Model load | Python startup
    %% CPU subprocess: 10 + 150 + 100 + 350 = 610
    \addplot[fill=gray!40, draw=gray!60] coordinates {
      (10, {CPU\\(subprocess hiện tại)})
      (10, {EdgeTPU\\(subprocess hiện tại)})
      (10, {CPU\\(persistent, tương lai)})
      (5,  {EdgeTPU\\(persistent, tương lai)})
    };  %% I/O
    \addplot[fill=warnred!55, draw=warnred!70] coordinates {
      (150, {CPU\\(subprocess hiện tại)})
      (10,  {EdgeTPU\\(subprocess hiện tại)})
      (150, {CPU\\(persistent, tương lai)})
      (10,  {EdgeTPU\\(persistent, tương lai)})
    };  %% Inference
    \addplot[fill=orng!65, draw=orng!80] coordinates {
      (100, {CPU\\(subprocess hiện tại)})
      (50,  {EdgeTPU\\(subprocess hiện tại)})
      (0,   {CPU\\(persistent, tương lai)})
      (0,   {EdgeTPU\\(persistent, tương lai)})
    };  %% Model load
    \addplot[fill=cppblue!55, draw=cppblue!70] coordinates {
      (350, {CPU\\(subprocess hiện tại)})
      (350, {EdgeTPU\\(subprocess hiện tại)})
      (0,   {CPU\\(persistent, tương lai)})
      (0,   {EdgeTPU\\(persistent, tương lai)})
    };  %% Python startup
    \legend{I/O (JPEG), Inference, Model load, Python startup}
    \end{axis}
  \end{tikzpicture}
  \caption{Phân tích latency mỗi lần gọi AI detection.
           Python startup time ($\sim$350~ms) chiếm phần lẫn tổng overhead
           trong kiến trúc subprocess hiện tại --- lớn hơn nhiều so với
           thời gian inference thực tế. Kiến trúc persistent process (tương lai)
           loại bỏ hoàn toàn chi phí này.}
  \label{fig:ai3-latency}
\end{figure}

\begin{table}[htbp]
  \centering
  \caption{So sánh tổng latency giữa các cấu hình triển khai}
  \label{tab:ai3-latency}
  \small
  \begin{tabular}{lcc}
    \toprule
    \textbf{Cấu hình}                    & \textbf{Inference} & \textbf{Tổng / lần gọi} \\
    \midrule
    CPU subprocess (hiện tại)            & $\sim$100--200~ms  & $\sim$600--700~ms \\
    EdgeTPU subprocess (hiện tại)       & $\sim$4--15~ms     & $\sim$420--430~ms \\
    CPU persistent (tương lai)           & $\sim$100--200~ms  & $\sim$160--210~ms \\
    EdgeTPU persistent (tương lai)      & $\sim$4--15~ms     & $\sim$15--20~ms  \\
    \bottomrule
  \end{tabular}
\end{table}

Kết quả trên cho thấy rõ điểm nghẻn hiện tại:
EdgeTPU inference chỉ mất 4--15~ms nhưng toàn bộ lá gọi
vẫn tốn 420~ms do phải khởi động Python mỗi lần. Với
kiến trúc \textbf{persistent Python process} --- một process
duy nhất khởi động một lần, load mô hình vào bộ nhớ, sau
đó tiếp nhận request từ C++ qua stdin/stdout pipe --- latency
mỗi lần gọi giảm xuống còn $\sim$15--20~ms với EdgeTPU,
mở ra khả năng chạy AI detection \textit{mỗi frame} camera
thay vì chỉ khi cần.

% Chú thích về nguồn số liệu
\smallskip
\noindent\textit{\small
Lưu ý: COCO metrics trong Bảng~\ref{tab:ai3-metrics} và Hình~\ref{fig:ai3-coco}
lấy từ kết quả đánh giá hai lớp (cylinder + cube) trên repository huấn luyện.
Phân tích confidence score trong Bảng~\ref{tab:ai3-conf} lấy từ kiểm tra
độc lập trên 218 ảnh cylinder và 10 ảnh cube (CPU inference via ai-edge-litert).
}
   %% mục 5.4
%%     ...
%%   \end{document}
%%
%% Compile: pdflatex -> pdflatex (chạy 2 lần để có label/ref đúng)
%%=============================================================

%%   \begin{document}
%%     ...
%%     \chapter{Thiết kế và Hiện thực}
%%     ...
%%     %%=============================================================
%% chapter4_4.tex
%% Phần 4.4 — AI #3: Phát hiện vật thể cho pick & place
%%
%% Đây là \input file — thêm vào preamble của main.tex:
%%   \usepackage{tikz}
%%   \usetikzlibrary{shapes.geometric,arrows.meta,positioning,
%%     fit,backgrounds,calc,decorations.pathreplacing}
%%   \usepackage{pgfplots}
%%   \pgfplotsset{compat=1.18}
%%   \usepackage{listings,xcolor,booktabs,multirow,subcaption}
%%=============================================================

\definecolor{cppblue}{RGB}{30,100,200}
\definecolor{pygreen}{RGB}{40,150,80}
\definecolor{warnred}{RGB}{200,60,60}
\definecolor{orng}{RGB}{220,130,30}

\lstset{
  basicstyle=\ttfamily\small,
  keywordstyle=\color{cppblue},
  commentstyle=\color{gray!70},
  stringstyle=\color{orng},
  breaklines=true,
  frame=single,
  rulecolor=\color{gray!40},
  framexleftmargin=4pt,
  xleftmargin=8pt,
  numbers=none,
}

% ================================================================
\section{AI \#3 --- Tự động hoá khâu khởi tạo trong hệ thống pick \& place}
\label{sec:ai3}
% ================================================================

% ----------------------------------------------------------------
\subsection{Bối cảnh: vấn đề của hệ thống visual servoing ban đầu}
\label{subsec:ai3-context}
% ----------------------------------------------------------------

Hệ thống pick \& place chạy trên cánh tay robot Denso VS-6577E 6 bậc tự do,
sử dụng thư viện ViSP (Visual Servoing Platform) để điều khiển bằng thị giác.
Camera USB được gắn trực tiếp lên đầu công tác theo cấu hình \textit{eye-in-hand} ---
camera di chuyển cùng robot và nhìn thẳng về phía vật thể cần gắp trong suốt
quá trình tiếp cận.

Vấn đề cốt lõi của hệ thống ban đầu nằm ở bước khởi tạo: sau khi robot về
home pose cố định, \textbf{người vận hành phải click chuột vào vị trí cylinder}
trên cửa sổ hiển thị camera thì blob tracker (\texttt{vpDot2}) mới bắt đầu
hoạt động được. Visual servo chỉ chạy sau thao tác thủ công đó. Ngoài ra,
nếu tracker mất target trong lúc servo --- điều hoàn toàn có thể xảy ra khi
cylinder bị che khuất hoặc di chuyển ra ngoài vùng sáng --- hệ thống reset
hoàn toàn và lại đòi hỏi click thủ công ở lần tiếp theo. Mỗi chu kỳ
pick \& place đều cần sự can thiệp của người vận hành, loại bỏ khả năng
vận hành tự động không giám sát.

\begin{figure}[htbp]
  \centering
  \begin{tikzpicture}[
    node distance=0.8cm and 0.6cm,
    bx/.style={draw, rounded corners=4pt, text width=3.0cm, align=center,
               minimum height=0.9cm, font=\small},
    robot/.style={bx, draw=cppblue!70, fill=cppblue!8},
    manual/.style={bx, draw=warnred!70, fill=warnred!8},
    aibox/.style={bx, draw=pygreen!70, fill=pygreen!8},
    arr/.style={->, >=Stealth, thick},
    rarr/.style={arr, draw=warnred!70},
    garr/.style={arr, draw=pygreen!70},
  ]

  %% --- CỘT TRÁI: hệ thống gốc ---
  \node[robot]                    (L1) {PREINIT\\{\scriptsize robot $\to$ home pose}};
  \node[manual, below=of L1]      (L2) {Click thủ công\\{\scriptsize chọn cylinder}};
  \node[robot,  below=of L2]      (L3) {JOINT\\{\scriptsize vpDot2 tracking}};
  \node[manual, below=of L3]      (L4) {Mất tracking\\{\scriptsize $\to$ reset toàn bộ}};

  \draw[arr]  (L1) -- (L2);
  \draw[arr]  (L2) -- (L3);
  \draw[rarr] (L3) -- (L4);
  \draw[rarr, dashed] (L4.west)
        to[out=190,in=170,looseness=1.3] (L1.west)
        node[midway, left, font=\scriptsize, text=warnred!70] {reset};

  \node[font=\bfseries\small, text=warnred, above=0.5cm of L1]
       {\textit{Trước cải tiến}};

  %% --- CỘT PHẢI: hệ thống mới ---
  \node[robot,  right=3.8cm of L1] (R1) {PREINIT\\{\scriptsize robot $\to$ home pose}};
  \node[aibox, below=of R1]        (R2) {AI detect\\{\scriptsize EfficientDet-Lite0}};
  \node[robot,  below=of R2]       (R3) {JOINT\\{\scriptsize AI tracking mỗi frame}};
  \node[aibox, below=of R3]        (R4) {AI re-detect\\{\scriptsize tự phục hồi}};

  \draw[arr]  (R1) -- (R2);
  \draw[garr] (R2) -- (R3);
  \draw[garr] (R3) -- (R4);
  \draw[garr, dashed] (R4.east)
        to[out=0,in=0,looseness=1.5] (R3.east)
        node[midway, right, font=\scriptsize, text=pygreen!70] {tiếp tục};

  %% fallback arrow từ AI detect
  \node[manual, right=1.2cm of R3] (FB)
       {{\scriptsize fallback:}\\click thủ công};
  \draw[rarr, dashed] (R2.east) -| (FB.north)
        node[near start, above, font=\scriptsize, text=warnred] {AI fail};

  \node[font=\bfseries\small, text=pygreen!80!black, above=0.5cm of R1]
       {\textit{Sau cải tiến}};

  %% Đường phân cách
  \draw[thick, dashed, gray!50]
      ($(L1.east)!0.5!(R1.west)+(0,0.8)$) -- ++(0,-6.5);

  \end{tikzpicture}
  \caption{So sánh luồng vận hành trước và sau khi tích hợp module AI.
           Bước yêu cầu thao tác thủ công (đỏ) được thay thế bằng AI tự động (xanh lá).
           Fallback về click thủ công vẫn được giữ lại cho trường hợp AI thất bại.}
  \label{fig:ai3-before-after}
\end{figure}

Hình~\ref{fig:ai3-before-after} minh hoạ sự thay đổi này. Module AI đảm nhiệm
toàn bộ việc phát hiện cylinder ở bước khởi tạo và tự phục hồi khi mất target
trong quá trình servo. Hệ thống vẫn giữ nguyên fallback về click thủ công để
đảm bảo tương thích ngược.


% ----------------------------------------------------------------
\subsection{Xây dựng và huấn luyện mô hình EfficientDet-Lite0}
\label{subsec:ai3-model}
% ----------------------------------------------------------------

\subsubsection*{Lựa chọn kiến trúc}

Yêu cầu đặt ra khá cụ thể: mô hình phải đủ nhỏ để triển khai được trên
Google Coral EdgeTPU, phải hỗ trợ INT8 quantization, và tương thích với
\texttt{tflite\_model\_maker} --- framework đã có sẵn trong môi trường phát triển.
\textbf{EfficientDet-Lite0} đáp ứng tất cả: đây là biến thể nhỏ nhất trong
họ EfficientDet được Google tối ưu cho thiết bị edge, input cố định
$320 \times 320$ pixel, backbone là EfficientNet-B0 đã được simplify.

\begin{table}[htbp]
  \centering
  \caption{Thông số mô hình EfficientDet-Lite0}
  \label{tab:ai3-model}
  \small
  \begin{tabular}{ll}
    \toprule
    \textbf{Thuộc tính}       & \textbf{Giá trị} \\
    \midrule
    Kiến trúc                 & EfficientDet-Lite0 \\
    Framework training        & TensorFlow~2.8.4 / \texttt{tflite\_model\_maker}~0.4.3 \\
    Input size                & $320 \times 320$~px, uint8 \\
    Số lớp đối tượng          & 2 (cylinder~--~lớp~0, cube~--~lớp~1) \\
    Số epochs / batch size    & 50 / 8 \\
    Transfer learning         & Frozen backbone, fine-tune detection head \\
    Target hardware           & Google Coral USB Accelerator (EdgeTPU) \\
    \bottomrule
  \end{tabular}
\end{table}

Mô hình được huấn luyện để phát hiện \textbf{hai lớp}: cylinder và cube nhựa PLA
in 3D. Việc đưa thêm cube vào tập huấn luyện là có chủ ý --- giúp mô hình học
cách phân biệt hai loại vật thể, từ đó giảm false positive khi cả hai cùng
xuất hiện trong cảnh. Ở phía tích hợp với ViSP, chỉ kết quả cylinder
(\texttt{class\_id~==~0}) được dùng; dự đoán cube bị lọc bỏ.

\subsubsection*{Tập dữ liệu}

Toàn bộ ảnh được thu trực tiếp từ camera eye-in-hand gắn trên robot Denso,
đảm bảo điều kiện thực tế sát với lúc triển khai. Ảnh bao gồm nhiều cấu hình:
góc chụp khác nhau, khoảng cách thay đổi, và số lượng vật thể từ 1 đến 5 trong
một cảnh.

\begin{figure}[htbp]
  \centering
  \begin{tikzpicture}
    \begin{axis}[
      xbar, bar width=0.42cm,
      xlabel={Số lượng ảnh (xấp xỉ)},
      ylabel={Danh mục},
      symbolic y coords={
        Background,
        2 cube + 3 cyl,
        2 cylinders,
        3 cylinders,
        1 cyl + 1 cube,
        1 cube,
        1 cylinder,
        Cylinder (gốc)
      },
      ytick=data,
      xmin=0, xmax=175,
      width=9.5cm, height=6.8cm,
      nodes near coords,
      nodes near coords align={horizontal},
      every node near coord/.style={font=\scriptsize},
      tick label style={font=\small},
      xlabel style={font=\small},
      ylabel style={font=\small},
      enlarge y limits=0.08,
    ]
    \addplot[fill=cppblue!55, draw=cppblue!80] coordinates {
      (38,  Background)
      (28,  2 cube + 3 cyl)
      (52,  2 cylinders)
      (43,  3 cylinders)
      (68,  1 cyl + 1 cube)
      (38,  1 cube)
      (108, 1 cylinder)
      (144, Cylinder (gốc))
    };
    \end{axis}
  \end{tikzpicture}
  \caption{Phân bố tập dữ liệu theo danh mục ảnh (~534 ảnh tổng,
           633 annotation cylinder, 170 annotation cube).
           Các cảnh đa vật thể giúp mô hình học phát hiện từng đối tượng
           riêng lẻ trong môi trường phức tạp.}
  \label{fig:ai3-dataset}
\end{figure}

Annotation được thực hiện bằng \texttt{labelme} theo format bounding box hình chữ
nhật. Workflow:
\begin{enumerate}
  \item Vẽ bounding box cho từng cylinder và cube bằng \texttt{labelme}
  \item Chuyển đổi labelme JSON $\to$ COCO JSON thống nhất
        (script \texttt{labelme\_to\_coco.py})
  \item Tập dữ liệu được split theo tỉ lệ $70/15/15$ (train/val/test),
        seed cố định để đảm bảo tính tái lập
\end{enumerate}

Phân chia ba chiều thay vì hai chiều đảm bảo tập test hoàn toàn độc lập ---
không được nhìn thấy trong suốt quá trình training hoặc điều chỉnh siêu tham số.

\subsubsection*{Pipeline huấn luyện}

\begin{figure}[htbp]
  \centering
  \begin{tikzpicture}[
    node distance=0.45cm and 0.15cm,
    st/.style={draw, rounded corners=3pt, fill=gray!10, draw=gray!55,
               text width=2.0cm, align=center, minimum height=0.85cm, font=\small},
    blue_st/.style={st, fill=cppblue!10, draw=cppblue!55},
    out_st/.style={st, fill=orng!20, draw=orng!70},
    arr/.style={->, >=Stealth, thick, gray!60},
    lbl/.style={font=\scriptsize, text=gray!80},
  ]

  \node[blue_st] (s1) {Annotation\\{\scriptsize labelme}};
  \node[blue_st, right=of s1] (s2) {COCO\\JSON};
  \node[blue_st, right=of s2] (s3) {Pascal\\VOC XML};
  \node[blue_st, right=of s3] (s4) {Training\\{\scriptsize 50 epochs}};
  \node[out_st, right=of s4] (s5) {SavedModel\\{\scriptsize Float32}};
  \node[out_st, below=0.55cm of s5] (s6) {INT8\\TFLite};
  \node[out_st, below=0.55cm of s6] (s7) {EdgeTPU\\TFLite};

  \draw[arr] (s1) -- (s2) node[lbl,above,midway] {to\_coco.py};
  \draw[arr] (s2) -- (s3) node[lbl,above,midway] {tflite\_mm};
  \draw[arr] (s3) -- (s4);
  \draw[arr] (s4) -- (s5) node[lbl,above,midway] {export};
  \draw[arr] (s5) -- (s6) node[lbl,right,midway] {PTQ};
  \draw[arr] (s6) -- (s7) node[lbl,right,midway] {edgetpu\_compiler};

  \end{tikzpicture}
  \caption{Pipeline huấn luyện và xuất mô hình. Từ ảnh thô đến ba file
           mô hình phục vụ các môi trường triển khai khác nhau.
           Bước chuyển đổi sang Pascal VOC XML là yêu cầu bắt buộc của
           \texttt{tflite\_model\_maker}.}
  \label{fig:ai3-pipeline}
\end{figure}

Transfer learning với \texttt{train\_whole\_model=False} --- chỉ fine-tune detection
head, giữ nguyên backbone đã pretrain --- phù hợp với tập dữ liệu tương đối nhỏ
($\sim$374 ảnh train) và giảm nguy cơ overfitting đáng kể so với
huấn luyện từ đầu.


% ----------------------------------------------------------------
\subsection{Quantization INT8 và biên dịch cho EdgeTPU}
\label{subsec:ai3-quant}
% ----------------------------------------------------------------

Mô hình sau training là Float32 --- không chạy được trực tiếp trên Coral EdgeTPU.
EdgeTPU chỉ thực thi mô hình INT8 quantized, trong đó trọng số và activation
được biểu diễn bằng số nguyên 8-bit. Lợi ích: kích thước giảm $\sim4\times$,
và phần cứng EdgeTPU đạt tốc độ inference cực nhanh ($\sim$4--15~ms/frame)
nhờ tận dụng phép tính số nguyên.

Post-Training Quantization (PTQ) sử dụng tập train làm \textit{representative
dataset} để calibrate scale và zero\_point cho từng layer --- đây là bước bắt
buộc để duy trì độ chính xác của INT8 model gần với Float32 baseline.

File \texttt{\_edgetpu.tflite} sau khi biên dịch chứa custom op
\texttt{edgetpu-custom-op} mà TFLite runtime thông thường
\textit{không thể thực thi trên CPU}, gây lỗi:

\begin{lstlisting}[language={}]
RuntimeError: Encountered unresolved custom op: edgetpu-custom-op
\end{lstlisting}

Do đó hệ thống \textbf{duy trì hai file riêng biệt} và lựa chọn tự động:

\begin{center}
\small
\begin{tabular}{lll}
  \toprule
  File & Runtime & Deployment \\
  \midrule
  \texttt{cylinder\_detector\_int8\_edgetpu.tflite} & pycoral   & Coral EdgeTPU \\
  \texttt{cylinder\_detector\_int8.tflite}          & TFLite/ai-edge-litert & CPU \\
  \bottomrule
\end{tabular}
\end{center}


% ----------------------------------------------------------------
\subsection{Kiến trúc tích hợp: C++ ViSP ↔ Python subprocess}
\label{subsec:ai3-arch}
% ----------------------------------------------------------------

\subsubsection*{Quyết định thiết kế}

Thay vì link trực tiếp thư viện TFLite C++ vào build system của ViSP
(đòi hỏi thay đổi lớn vào CMakeLists và phụ thuộc version build cụ thể),
hệ thống áp dụng kiến trúc \textbf{subprocess model}: C++ gọi Python script
như một tiến trình con qua \texttt{popen()}, giao tiếp một chiều qua stdout.
Cách này giúp module AI hoàn toàn tách biệt --- cập nhật mô hình không cần
recompile code robot.

\begin{figure}[htbp]
  \centering
  \begin{tikzpicture}[
    node distance=0.38cm and 0.5cm,
    comp/.style={draw, rounded corners=4pt, text width=3.1cm, align=center,
                 minimum height=0.82cm, font=\small},
    cpp/.style={comp, fill=cppblue!10, draw=cppblue!55},
    py/.style={comp,  fill=pygreen!10, draw=pygreen!60},
    tmp/.style={comp,  fill=orng!18, draw=orng!65, text width=2.5cm,
                minimum height=0.7cm, font=\scriptsize},
    arr/.style={->, >=Stealth, thick},
    ]

  %% ---- C++ side ----
  \node[cpp] (cap)   {Camera capture\\{\scriptsize vpImage{$<$}uchar{$>$}}};
  \node[cpp, below=of cap]  (save)  {Save frame\\{\scriptsize cv::imwrite()}};
  \node[cpp, below=of save] (pop)   {popen() subprocess\\{\scriptsize build cmd}};
  \node[cpp, below=of pop]  (sel)   {select() timeout 3s\\{\scriptsize POSIX wait}};
  \node[cpp, below=of sel]  (parse) {Parse stdout\\{\scriptsize sscanf SUCCESS}};
  \node[cpp, below=of parse](feat)  {vpFeatureBuilder\\{\scriptsize create feature}};

  \draw[arr](cap)--(save); \draw[arr](save)--(pop);
  \draw[arr](pop)--(sel);  \draw[arr](sel)--(parse); \draw[arr](parse)--(feat);

  \node[draw=cppblue!50, dashed, rounded corners=6pt,
        fit=(cap)(save)(pop)(sel)(parse)(feat),
        inner sep=7pt] (cbox) {};
  \node[above=2pt of cbox, font=\bfseries\small, text=cppblue]
       {C++ (ViSP main loop)};

  %% ---- Temp file ----
  \node[tmp, right=2.6cm of save] (jpg)
       {/tmp/visp\_ai\_frame.jpg};

  %% ---- Python side ----
  \node[py, right=2.6cm of pop]   (load) {Load model\\{\scriptsize fallback chain}};
  \node[py, below=of load]        (pre)  {Preprocess\\{\scriptsize resize 320$\times$320}};
  \node[py, below=of pre]         (inf)  {Inference\\{\scriptsize invoke()}};
  \node[py, below=of inf]         (post) {Post-process\\{\scriptsize dequant+filter}};
  \node[py, below=of post]        (out)  {Print stdout\\{\scriptsize SUCCESS u v conf}};

  \draw[arr](load)--(pre); \draw[arr](pre)--(inf);
  \draw[arr](inf)--(post); \draw[arr](post)--(out);

  \node[draw=pygreen!55, dashed, rounded corners=6pt,
        fit=(load)(pre)(inf)(post)(out),
        inner sep=7pt] (pbox) {};
  \node[above=2pt of pbox, font=\bfseries\small, text=pygreen!80!black]
       {Python subprocess};

  %% ---- Data flow arrows ----
  \draw[arr, orng!80]  (save.east) -- (jpg.west)  node[midway,above,font=\scriptsize]{write};
  \draw[arr, orng!80]  (jpg.east)  -- (load.west) node[midway,above,font=\scriptsize]{read};
  \draw[arr, cppblue!60](pop.east)  -- (load.west) node[midway,above,font=\scriptsize,yshift=2pt]{spawn};
  \draw[arr, pygreen!70, bend right=25](out.west)  to (parse.east)
        node[midway,below,font=\scriptsize]{stdout pipe};

  \end{tikzpicture}
  \caption{Kiến trúc tích hợp C++ ViSP $\leftrightarrow$ Python subprocess.
           C++ lưu frame hiện tại ra file tạm, spawn Python làm child process
           qua \texttt{popen()}, đọc kết quả từ stdout sau khi qua timeout guard.}
  \label{fig:ai3-arch}
\end{figure}

\subsubsection*{Protocol trao đổi dữ liệu}

Output của Python script là một dòng duy nhất trên stdout:

\begin{lstlisting}[language={}]
SUCCESS 272.86 190.74 0.8362 238.42 166.89 68.88 47.68
%        u      v      conf   x_min  y_min  width  height
\end{lstlisting}

hoặc \texttt{FAILURE~<lý~do>} nếu không phát hiện được. Tất cả debug log
đi ra \texttt{stderr} --- C++ chỉ đọc stdout --- tránh nhiễu kết quả parse.

Đáng chú ý: trong ViSP, \texttt{vpImagePoint} được định nghĩa là \texttt{(row,~col)}
tương ứng \texttt{(v,~u)}, không phải \texttt{(u,~v)} như thông thường:

\begin{lstlisting}[language=C++]
detected_center = vpImagePoint(v, u);  // row=v, col=u
\end{lstlisting}

Tham số \texttt{hint} cho phép Python chọn cylinder \textit{gần nhất với vị
trí đã biết} thay vì cylinder có confidence cao nhất --- cơ chế này tránh
hiện tượng servo ``nhảy'' sang cylinder khác khi scores thay đổi giữa các frame.


% ----------------------------------------------------------------
\subsection{State machine: PREINIT \texorpdfstring{$\to$}{→} INIT (AI) \texorpdfstring{$\to$}{→} JOINT \texorpdfstring{$\to$}{→} ...}
\label{subsec:ai3-sm}
% ----------------------------------------------------------------

Hệ thống vận hành theo vòng lặp chính liên tục với 6 trạng thái.
Mỗi iteration đọc một frame camera và thực thi logic trạng thái hiện tại.

\begin{figure}[htbp]
  \centering
  \begin{tikzpicture}[
    node distance=1.0cm and 2.0cm,
    st/.style={draw, rounded corners=4pt, fill=cppblue!10, draw=cppblue!55,
               text width=2.3cm, align=center, minimum height=1.0cm, font=\small},
    ai/.style={st, fill=pygreen!10, draw=pygreen!60},
    arr/.style={->, >=Stealth, thick},
    lbl/.style={font=\scriptsize, text=gray!70},
  ]

  %% Top row
  \node[st]  (pre)  {PREINIT\\{\scriptsize home pose}};
  \node[ai,  right=of pre]   (ini)  {INIT\\{\scriptsize AI detect}};
  \node[ai,  right=of ini]   (jnt)  {JOINT\\{\scriptsize AI tracking}};
  \node[st,  right=of jnt]   (app)  {APPROACH\\{\scriptsize chờ RC5}};

  %% Bottom row
  \node[st,  below=1.8cm of app]  (grp)  {GRIPPER\\{\scriptsize đóng gripper}};
  \node[st,  left=of grp]         (cls)  {CLASSIFIED\\{\scriptsize vận chuyển}};

  %% Arrows
  \draw[arr](pre)--(ini)  node[lbl,above,midway]{home\\reached};
  \draw[arr](ini)--(jnt)  node[lbl,above,midway]{AI init\\ok};
  \draw[arr](jnt)--(app)  node[lbl,above,midway]{$|e|<5{\times}10^{-3}$};
  \draw[arr](app)--(grp)  node[lbl,right,midway]{RC5\\signal};
  \draw[arr](grp)--(cls)  node[lbl,below,midway]{grasped};
  \draw[arr](cls.west) to[out=180,in=270,looseness=0.8] (pre.south)
        node[lbl,left,pos=0.5]{cycle\\restart};

  %% AI self-recovery
  \draw[arr, pygreen!70, dashed]
        (jnt.north) to[out=100,in=80,looseness=2.5] (jnt.north)
        node[above=0.7cm of jnt, font=\scriptsize, text=pygreen!80!black]
        {AI re-detect};

  %% Fallback from INIT
  \draw[arr, warnred!70, dashed]
        (ini.south) -- ++(0,-0.7)
        node[below, font=\scriptsize, text=warnred]
        {fallback: click thủ công};

  \end{tikzpicture}
  \caption{State machine của hệ thống pick \& place sau cải tiến.
           Trạng thái tô xanh lá là nơi AI được gọi.
           Vòng lặp tự phục hồi ở JOINT (mũi tên đứt nét) cho phép hệ thống
           tiếp tục mà không cần reset khi mất detection.}
  \label{fig:ai3-sm}
\end{figure}

\paragraph{PREINIT.}
Robot gửi lệnh joint position về home pose
$[0^\circ, 0^\circ, 90^\circ, 0^\circ, 90^\circ, 0^\circ]$.
Servo task được cấu hình với interaction matrix Desired, pseudo-inverse,
gain $\lambda = 0.4$.

\paragraph{INIT.}
Chờ robot đạt home pose (sai số từng joint $< 0.01$) và gripper mở xong.
Khi cả hai điều kiện thỏa, \texttt{detectCylinderWithAI()} được gọi
không có hint --- Python chọn cylinder có confidence cao nhất.
Kết quả trực tiếp tạo visual feature \texttt{vpFeaturePoint}
và add vào servo task. Nếu AI thất bại, \texttt{ai\_hint}
giữ giá trị mặc định (fallback thủ công).

\paragraph{JOINT.}
Khác với cách tiếp cận truyền thống dùng blob tracker mỗi frame,
hệ thống hiện tại \textbf{gọi AI trực tiếp ở mỗi iteration}
với \texttt{hint = \&ai\_hint}. Control law:
\begin{equation}
  \mathbf{v} = -\lambda\,\mathbf{L}_{s^*}^{+}\,\mathbf{e},
  \qquad \mathbf{q}_{k+1} = \mathbf{q}_k + \mathbf{v}\,\Delta t
  \label{eq:ibvs}
\end{equation}
với $\Delta t = 0.5\,$s. Điều kiện hội tụ: $|e_x| < 5{\times}10^{-3}$
và $|e_y| < 5{\times}10^{-3}$ (normalized image coordinates).
Khi đạt, gửi \texttt{"OKE\textbackslash r"} qua UART đến RC5 controller.


% ----------------------------------------------------------------
\subsection{Các vấn đề kỹ thuật và giải pháp}
\label{subsec:ai3-issues}
% ----------------------------------------------------------------

Một số vấn đề nảy sinh trong quá trình tích hợp mà khó thấy trước khi thực
sự chạy hệ thống.

\subsubsection*{P1: Xác định output tensor sai lệch}

\texttt{tflite\_model\_maker} 0.4.x đặt tên tensor theo thứ tự riêng của nó:

\begin{center}
\small
\begin{tabular}{ccc}
  \toprule
  Tensor & Shape & Nội dung \\
  \midrule
  \texttt{:0} & \texttt{[1]}     & num\_detections \\
  \texttt{:1} & \texttt{[1,\,N]} & scores \quad \textbf{← cần lấy cái này} \\
  \texttt{:2} & \texttt{[1,\,N]} & class IDs \quad \textbf{← không phải cái này} \\
  \texttt{:3} & \texttt{[1,\,N,\,4]} & boxes \\
  \bottomrule
\end{tabular}
\end{center}

Thứ tự này không khớp với API của pycoral, vốn giả định positional order
\texttt{boxes/classes/scores/count = 0/1/2/3}. Dùng API đó sẽ đọc
num\_detections (raw uint8 = 255) như scores và crash với \texttt{IndexError}.
Giải pháp: xác định tensor \textbf{theo shape động} --- boxes là
\texttt{[1,N,4]}, scores là tensor \texttt{[1,N]} có name suffix nhỏ nhất
(\texttt{:1} chứ không phải \texttt{:2}).

\subsubsection*{P2: Python 3.10 không có \texttt{tflite-runtime}}

Google không publish wheel \texttt{tflite-runtime} cho Python 3.10+.
Giải pháp là xây dựng \textbf{fallback chain 4 tầng}:

\begin{figure}[htbp]
  \centering
  \begin{tikzpicture}[
    node distance=0.4cm,
    fb/.style={draw, rounded corners=3pt, text width=7.5cm, align=left,
               minimum height=0.85cm, font=\small, inner sep=7pt},
    t1/.style={fb, fill=purple!12, draw=purple!55},
    t2/.style={fb, fill=teal!10,   draw=teal!50},
    t3/.style={fb, fill=cppblue!10, draw=cppblue!50},
    t4/.style={fb, fill=gray!12,   draw=gray!50},
    fa/.style={fb, fill=warnred!10, draw=warnred!50},
    arr/.style={->, >=Stealth, thick, dashed, gray!55},
    lbl/.style={font=\scriptsize, text=warnred!80},
  ]
  \node[t1] (a) {\textbf{1.}~\texttt{pycoral.utils.edgetpu}\quad{\scriptsize Coral EdgeTPU, Python 3.9}};
  \node[t2, below=of a] (b) {\textbf{2.}~\texttt{ai\_edge\_litert.interpreter}\quad{\scriptsize Google's new pkg, Python 3.8--3.12}};
  \node[t3, below=of b] (c) {\textbf{3.}~\texttt{tflite\_runtime.interpreter}\quad{\scriptsize Legacy, Python 3.9}};
  \node[t4, below=of c] (d) {\textbf{4.}~\texttt{tensorflow.lite.Interpreter}\quad{\scriptsize Heavy fallback}};
  \node[fa, below=of d] (e) {\textbf{FAILURE}:~\texttt{no\_tflite\_backend\_available}};

  \draw[arr](a)--(b) node[lbl,right,midway]{ImportError / Coral not connected};
  \draw[arr](b)--(c) node[lbl,right,midway]{ImportError};
  \draw[arr](c)--(d) node[lbl,right,midway]{ImportError};
  \draw[arr](d)--(e) node[lbl,right,midway]{ImportError};
  \end{tikzpicture}
  \caption{Python backend fallback chain trong \texttt{detect\_cylinder.py}.
           Script tự động thử từng backend theo thứ tự ưu tiên,
           tương thích với cả môi trường có Coral lẫn CPU-only.}
  \label{fig:ai3-fallback}
\end{figure}

\texttt{ai-edge-litert} là gói chính thức của Google thay thế \texttt{tflite-runtime},
hỗ trợ Python 3.8--3.12 với API giống hệt nhau --- là giải pháp cho vấn đề Python 3.10.

\subsubsection*{P3: JSON comment trong config}

\texttt{config.json} dùng comment \texttt{//} để ghi chú, nhưng cả Python
\texttt{json.load()} và \texttt{nlohmann::json} đều từ chối mặc định.
Phía C++: \texttt{nlohmann::json::parse(file, nullptr, true, true)} với
\texttt{ignore\_comments=true} (có từ nlohmann/json~3.9.0).
Phía Python: strip comment lines thủ công trước khi \texttt{json.loads()}.


% ----------------------------------------------------------------
\subsection{Thông số cấu hình visual servo}
\label{subsec:ai3-servo}
% ----------------------------------------------------------------

Hệ thống sử dụng \textbf{Image-Based Visual Servoing (IBVS)} với feature là
tọa độ 2D của cylinder trong không gian ảnh normalized $\mathbf{s}=(x,y)$,
desired position $\mathbf{s}^*=(0,0,1)$ --- cylinder phải nằm đúng tâm ảnh.

\begin{table}[htbp]
  \centering
  \caption{Thông số điều khiển visual servo}
  \label{tab:ai3-servo}
  \small
  \begin{tabular}{ll}
    \toprule
    \textbf{Tham số}            & \textbf{Giá trị} \\
    \midrule
    Loại điều khiển             & Eye-in-hand, \texttt{EYEINHAND\_L\_cVe\_eJe} \\
    Interaction matrix          & Desired ($\mathbf{L}_{s^*}$), pseudo-inverse \\
    Gain $\lambda$              & 0.4 \\
    Bước thời gian $\Delta t$   & 0.5\,s \\
    Ngưỡng hội tụ               & $|e_x| < 5{\times}10^{-3}$ và $|e_y| < 5{\times}10^{-3}$ \\
    Desired position            & $(x^*, y^*) = (0,\,0)$ --- tâm ảnh, $Z = 1$ \\
    \bottomrule
  \end{tabular}
\end{table}

Dùng $\mathbf{L}_{s^*}$ (tại vị trí mong muốn) thay vì $\mathbf{L}_{s_{cur}}$
giúp tránh singularity khi feature tiến gần desired point và giữ ổn định
theo ổn định hơn trong giai đoạn hội tụ. Velocity được saturate trước khi
gửi xuống robot để đảm bảo không vượt giới hạn vật lý từng khớp.
   %% mục 4.4
%%     ...
%%     \chapter{Kết quả và Đánh giá}
%%     ...
%%     %%=============================================================
%% chapter5_4.tex
%% Phần 5.4 — Kết quả đánh giá: AI #3 Phát hiện vật thể
%%
%% \input file — cần thêm vào preamble:
%%   \usepackage{tikz}
%%   \usetikzlibrary{arrows.meta,positioning,patterns}
%%   \usepackage{pgfplots}
%%   \pgfplotsset{compat=1.18}
%%   \usepackage{booktabs,multirow,subcaption,xcolor}
%%=============================================================

% ================================================================
\section{Kết quả và Đánh giá: AI \#3 --- Phát hiện vật thể}
\label{sec:ai3-results}
% ================================================================

% ----------------------------------------------------------------
\subsection{Kết quả detection model: COCO metrics}
\label{subsec:ai3-coco}
% ----------------------------------------------------------------

Mô hình được đánh giá trên cả validation set và test set độc lập
theo chuẩn \textbf{COCO evaluation protocol} --- tính trung bình mAP qua
các ngưỡng IoU từ 0.50 đến 0.95 với bước 0.05.

\begin{figure}[htbp]
  \centering
  \begin{tikzpicture}
    \begin{axis}[
      ybar, bar width=0.45cm,
      width=10cm, height=7cm,
      symbolic x coords={mAP (0.50:0.95), AP50, AP75, AR},
      xtick=data,
      x tick label style={font=\small, align=center},
      ymin=0, ymax=1.15,
      ylabel={Score},
      ylabel style={font=\small},
      legend style={at={(0.5,1.03)}, anchor=south, legend columns=2, font=\small},
      legend cell align=left,
      nodes near coords,
      every node near coord/.style={font=\scriptsize, /pgf/number format/fixed,
                                    /pgf/number format/precision=3},
      enlarge x limits=0.25,
      grid=major, grid style={dashed, gray!30},
      axis line style={gray!60},
      tick style={gray!60},
    ]
    \addplot[fill=cppblue!70, draw=cppblue!90] coordinates {
      (mAP (0.50:0.95), 0.791)
      (AP50, 0.990)
      (AP75, 0.990)
      (AR,   0.869)
    };
    \addplot[fill=orng!80, draw=orng!90] coordinates {
      (mAP (0.50:0.95), 0.799)
      (AP50, 1.000)
      (AP75, 0.969)
      (AR,   0.856)
    };
    \legend{Validation, Test}
    \end{axis}
  \end{tikzpicture}
  \caption{COCO metrics trên validation set và test set. AP50 trên test set
           đạt 1.000 cho thấy 100\% cylinder được phát hiện đúng ở ngưỡng IoU
           lỏng lẻ (50\%). mAP$\approx$0.80 phản ánh vẫn còn đây đó những trường hợp
           định vị bounding box chưa được chính xác ở ngưỡng IoU cao.}
  \label{fig:ai3-coco}
\end{figure}

Bảng~\ref{tab:ai3-metrics} tổng hợp số liệu chi tiết:

\begin{table}[htbp]
  \centering
  \caption{COCO metrics trên validation và test set}
  \label{tab:ai3-metrics}
  \small
  \begin{tabular}{lcc}
    \toprule
    \textbf{Metric}                    & \textbf{Validation} & \textbf{Test} \\
    \midrule
    mAP (IoU 0.50:0.95)               & 0.791               & \textbf{0.799} \\
    AP50 (IoU $\geq$ 0.50)            & 0.990               & \textbf{1.000} \\
    AP75 (IoU $\geq$ 0.75)            & 0.990               & 0.969 \\
    AR (maxDets=100)                   & 0.869               & 0.856 \\
    \midrule
    Precision (IoU=0.50, conf=0.50)    & ---                 & \textbf{1.000} \\
    Recall    (IoU=0.50, conf=0.50)    & ---                 & 0.990 \\
    F1-Score                           & ---                 & \textbf{0.995} \\
    TP / FP / FN                       & ---                 & 95 / 0 / 1 \\
    \bottomrule
  \end{tabular}
\end{table}

\begin{figure}[htbp]
  \centering
  \begin{tikzpicture}
    \begin{axis}[
      ybar, bar width=0.55cm,
      width=8.5cm, height=6cm,
      symbolic x coords={Precision, Recall, F1-Score},
      xtick=data,
      x tick label style={font=\small},
      ymin=0.96, ymax=1.04,
      ylabel={Score},
      ylabel style={font=\small},
      nodes near coords,
      every node near coord/.style={font=\small, /pgf/number format/fixed,
                                    /pgf/number format/precision=3},
      enlarge x limits=0.3,
      grid=major, grid style={dashed, gray!30},
      axis line style={gray!60},
    ]
    \addplot[fill=pygreen!65, draw=pygreen!80] coordinates {
      (Precision, 1.000)
      (Recall,    0.990)
      (F1-Score,  0.995)
    };
    \end{axis}
  \end{tikzpicture}
  \caption{Precision, Recall và F1-Score trên test set (PASCAL VOC protocol,
           IoU$\geq$0.50, confidence$\geq$0.50, chỉ lớp cylinder).
           Precision = 1.000 và FP = 0 chứng tỏ không có false positive
           nào ở ngưỡng 0.50. FN = 1 (một cylinder bị bỏ sót) giải thích
           Recall = 0.990.}
  \label{fig:ai3-prf}
\end{figure}

Giá trị AP50 = 1.000 trên test set à ngưỡng IoU lỏng lẻ (50\%) phản ánh
mô hình phát hiện được toàn bộ cylinder. mAP = 0.799 (trung bình
qua IoU 0.50--0.95) thấp hơn một chút do mô hình EfficientDet-Lite0
đơn giản vẫn còn giới hạn trong việc định vị bounding box chính xác
ở các ngưỡng IoU cao ($\geq$0.75) --- điều phổ biến với các kiến trúc
nhe hơn.


% ----------------------------------------------------------------
\subsection{Phân tích confidence score: true positive vs false positive}
\label{subsec:ai3-conf}
% ----------------------------------------------------------------

Một trong những điều quan trọng khi triển khai thực tế là khoảng cách
phân tách giữa true positive và false positive --- điều này quyết định
ngưỡng confidence nào thực tế dùng được mà không sai kết quả.

Kiểm tra trên 218 ảnh cylinder (true positives) và 10 ảnh cube
(false positives):

\begin{figure}[htbp]
  \centering
  \begin{tikzpicture}
    \begin{axis}[
      width=11cm, height=5.5cm,
      xlabel={Confidence score},
      ylabel={},
      xmin=0, xmax=1.05,
      ymin=0, ymax=3,
      xtick={0,0.1,...,1.0},
      ytick={0.7, 2.0},
      yticklabels={\small False Positive\\(cube, $n$=10), \small True Positive\\(cylinder, $n$=218)},
      y tick label style={align=right, font=\small},
      x tick label style={font=\small},
      xlabel style={font=\small},
      grid=both, grid style={dashed, gray!25},
      axis line style={gray!60},
    ]

    %% --- True Positive band (y=2) ---
    %% Most scores: 0.85-0.98 (majority)
    \addplot[only marks, mark=|, mark size=5pt, mark options={draw=pygreen!70, line width=1.2pt},
             opacity=0.5]
    coordinates {
      (0.29,2)(0.61,2)(0.74,2)
      (0.82,2)(0.83,2)(0.84,2)(0.85,2)(0.86,2)(0.87,2)(0.88,2)
      (0.89,2)(0.90,2)(0.91,2)(0.92,2)(0.93,2)(0.94,2)(0.94,2)
      (0.95,2)(0.95,2)(0.96,2)(0.96,2)(0.97,2)(0.97,2)(0.98,2)
    };

    %% --- False Positive band (y=0.7) ---
    \addplot[only marks, mark=|, mark size=5pt, mark options={draw=warnred!70, line width=1.2pt},
             opacity=0.6]
    coordinates {
      (0.08,0.7)(0.13,0.7)(0.17,0.7)(0.21,0.7)
      (0.25,0.7)(0.28,0.7)(0.30,0.7)(0.33,0.7)(0.35,0.7)(0.36,0.7)
    };

    %% --- Threshold line at 0.50 ---
    \addplot[thick, dashed, draw=gray!80, domain=0:3]
    coordinates {(0.50, 0) (0.50, 3)};
    \node[font=\scriptsize, text=gray!80, rotate=90]
         at (axis cs: 0.52, 1.5) {threshold = 0.50};

    %% --- Gap annotation ---
    \draw[->, >=Stealth, thick, gray!70]
         (axis cs:0.38, 2.5) -- (axis cs:0.48, 2.5);
    \draw[->, >=Stealth, thick, gray!70]
         (axis cs:0.34, 2.5) -- (axis cs:0.27, 2.5);
    \node[font=\scriptsize, text=gray!80] at (axis cs:0.37, 2.7) {gap = 0.34};

    %% --- Median annotation for TP ---
    \addplot[only marks, mark=diamond*, mark size=4pt,
             mark options={draw=pygreen!80!black, fill=pygreen!60}]
    coordinates {(0.94, 2)};
    \node[font=\scriptsize, text=pygreen!80!black, above]
         at (axis cs:0.94, 2) {median 0.94};

    %% --- Max FP annotation ---
    \addplot[only marks, mark=diamond*, mark size=4pt,
             mark options={draw=warnred!80, fill=warnred!50}]
    coordinates {(0.36, 0.7)};
    \node[font=\scriptsize, text=warnred, below]
         at (axis cs:0.36, 0.7) {max 0.36};

    \end{axis}
  \end{tikzpicture}
  \caption{Phân tách confidence score giữa true positive (cylinder, xanh lá)
           và false positive (cube, đỏ). Mỗi vạch đứng đại diện một detection.
           Ngưỡng 0.50 nằm chính giữa khoảng gap 0.34,
           loại bỏ toàn bộ false positive mà vẫn chấp nhận 99.5\% true positive.}
  \label{fig:ai3-conf}
\end{figure}

Khoảng cách 0.34 giữa điểm false positive cao nhất (0.36) và 99.1\%
true positive ($\geq$0.70) xác nhận mô hình đã học được đặc trưng
phân biệt mạnh giữa cylinder và cube. Ngưỡng hiện tại 0.50 nằm an toàn
giữa hai vùng phân bố này.

Ba ảnh outlier có score thấp (0.29, 0.61, 0.74) tương ứng với cylinder
bị nhìn từ góc bất thường hoặc bị che khuất một phần. Ỡ runtime,
các trường hợp này fallback về click thủ công --- hành vi an toàn.

\begin{table}[htbp]
  \centering
  \caption{Phân tích confidence score trên 218 ảnh cylinder + 10 ảnh cube}
  \label{tab:ai3-conf}
  \small
  \begin{tabular}{lcc}
    \toprule
    & \textbf{True Positive (cylinder)} & \textbf{False Positive (cube)} \\
    \midrule
    Score range      & 0.29 -- 0.98          & 0.08 -- 0.36 \\
    Median           & \textbf{0.94}         & --- \\
    $\geq$ 0.50      & 99.5\% (217/218)      & \textbf{0/10 (0\%)} \\
    $\geq$ 0.70      & 99.1\% (216/218)      & 0/10 (0\%) \\
    $\geq$ 0.80      & 95.9\% (209/218)      & 0/10 (0\%) \\
    \midrule
    Gap ph\^an t\'ach & \multicolumn{2}{c}{$0.70 - 0.36 = 0.34$} \\
    \bottomrule
  \end{tabular}
\end{table}


% ----------------------------------------------------------------
\subsection{Kết quả tích hợp hệ thống}
\label{subsec:ai3-integration}
% ----------------------------------------------------------------

\begin{table}[htbp]
  \centering
  \caption{So sánh khả năng trước và sau khi tích hợp AI}
  \label{tab:ai3-comparison}
  \small
  \begin{tabular}{lll}
    \toprule
    \textbf{Tiêu chí}                    & \textbf{Hệ thống gốc} & \textbf{Sau cải tiến} \\
    \midrule
    Khởi tạo mỗi chu kỳ               & Click thủ công (bắt buộc) & Tự động (AI) \\
    Tỷ lệ init thành công             & 100\% (do manual)  & $\sim$99.5\% (AI) \\
    Fallback khi AI thất bại           & Không có              & Click thủ công \\
    Phục hồi khi mất tracking          & Reset toàn bộ        & AI re-detect tự động \\
    Chạy không giám sát                & Không thể              & Có (trường hợp bình thường) \\
    \bottomrule
  \end{tabular}
\end{table}

Thiết kế hiện tại gọi AI trực tiếp \textbf{mỗi iteration} của servo
loop thay vì dùng blob tracker. Điều này có một hàm ý đáng chú ý:
\textit{không có tracking failure} theo nghĩa truyền thống --- mỗi frame
là một lần detection độc lập. Tính ổn định của servo phụ thuộc vào
sự nhất quán của AI detection qua các frame liên tiếp. Cơ chế \texttt{--hint}
(truyền vị trí frame trước) giúp được phần nào khi có nhiều cylinder
trong cảnh.


% ----------------------------------------------------------------
\subsection{Phân tích latency: subprocess CPU vs EdgeTPU}
\label{subsec:ai3-latency}
% ----------------------------------------------------------------

Mỗi lần gọi \texttt{detectCylinderWithAI()} bao gồm nhiều thành phần
latency khác nhau. Điều ít người để ý là: phần overhead lớn nhất
không phải là inference mà là \textbf{Python startup time} --- chi phí
mở một process mới mỗi lần gọi \texttt{popen()}.

\begin{figure}[htbp]
  \centering
  \begin{tikzpicture}
    \begin{axis}[
      xbar stacked,
      bar width=0.55cm,
      width=11cm, height=6.5cm,
      symbolic y coords={
        {EdgeTPU\\(persistent, tương lai)},
        {CPU\\(persistent, tương lai)},
        {EdgeTPU\\(subprocess hiện tại)},
        {CPU\\(subprocess hiện tại)}
      },
      ytick=data,
      y tick label style={align=right, font=\small},
      xlabel={Thời gian (ms)},
      xlabel style={font=\small},
      xmin=0, xmax=780,
      legend style={at={(0.5,-0.18)}, anchor=north, legend columns=4,
                    font=\scriptsize},
      legend cell align=left,
      grid=major, grid style={dashed, gray!25},
      axis line style={gray!60},
      enlarge y limits=0.18,
      nodes near coords, nodes near coords align={horizontal},
      point meta=explicit symbolic,
      every node near coord/.style={font=\scriptsize},
    ]
    %%  stacks: I/O | Inference | Model load | Python startup
    %% CPU subprocess: 10 + 150 + 100 + 350 = 610
    \addplot[fill=gray!40, draw=gray!60] coordinates {
      (10, {CPU\\(subprocess hiện tại)})
      (10, {EdgeTPU\\(subprocess hiện tại)})
      (10, {CPU\\(persistent, tương lai)})
      (5,  {EdgeTPU\\(persistent, tương lai)})
    };  %% I/O
    \addplot[fill=warnred!55, draw=warnred!70] coordinates {
      (150, {CPU\\(subprocess hiện tại)})
      (10,  {EdgeTPU\\(subprocess hiện tại)})
      (150, {CPU\\(persistent, tương lai)})
      (10,  {EdgeTPU\\(persistent, tương lai)})
    };  %% Inference
    \addplot[fill=orng!65, draw=orng!80] coordinates {
      (100, {CPU\\(subprocess hiện tại)})
      (50,  {EdgeTPU\\(subprocess hiện tại)})
      (0,   {CPU\\(persistent, tương lai)})
      (0,   {EdgeTPU\\(persistent, tương lai)})
    };  %% Model load
    \addplot[fill=cppblue!55, draw=cppblue!70] coordinates {
      (350, {CPU\\(subprocess hiện tại)})
      (350, {EdgeTPU\\(subprocess hiện tại)})
      (0,   {CPU\\(persistent, tương lai)})
      (0,   {EdgeTPU\\(persistent, tương lai)})
    };  %% Python startup
    \legend{I/O (JPEG), Inference, Model load, Python startup}
    \end{axis}
  \end{tikzpicture}
  \caption{Phân tích latency mỗi lần gọi AI detection.
           Python startup time ($\sim$350~ms) chiếm phần lẫn tổng overhead
           trong kiến trúc subprocess hiện tại --- lớn hơn nhiều so với
           thời gian inference thực tế. Kiến trúc persistent process (tương lai)
           loại bỏ hoàn toàn chi phí này.}
  \label{fig:ai3-latency}
\end{figure}

\begin{table}[htbp]
  \centering
  \caption{So sánh tổng latency giữa các cấu hình triển khai}
  \label{tab:ai3-latency}
  \small
  \begin{tabular}{lcc}
    \toprule
    \textbf{Cấu hình}                    & \textbf{Inference} & \textbf{Tổng / lần gọi} \\
    \midrule
    CPU subprocess (hiện tại)            & $\sim$100--200~ms  & $\sim$600--700~ms \\
    EdgeTPU subprocess (hiện tại)       & $\sim$4--15~ms     & $\sim$420--430~ms \\
    CPU persistent (tương lai)           & $\sim$100--200~ms  & $\sim$160--210~ms \\
    EdgeTPU persistent (tương lai)      & $\sim$4--15~ms     & $\sim$15--20~ms  \\
    \bottomrule
  \end{tabular}
\end{table}

Kết quả trên cho thấy rõ điểm nghẻn hiện tại:
EdgeTPU inference chỉ mất 4--15~ms nhưng toàn bộ lá gọi
vẫn tốn 420~ms do phải khởi động Python mỗi lần. Với
kiến trúc \textbf{persistent Python process} --- một process
duy nhất khởi động một lần, load mô hình vào bộ nhớ, sau
đó tiếp nhận request từ C++ qua stdin/stdout pipe --- latency
mỗi lần gọi giảm xuống còn $\sim$15--20~ms với EdgeTPU,
mở ra khả năng chạy AI detection \textit{mỗi frame} camera
thay vì chỉ khi cần.

% Chú thích về nguồn số liệu
\smallskip
\noindent\textit{\small
Lưu ý: COCO metrics trong Bảng~\ref{tab:ai3-metrics} và Hình~\ref{fig:ai3-coco}
lấy từ kết quả đánh giá hai lớp (cylinder + cube) trên repository huấn luyện.
Phân tích confidence score trong Bảng~\ref{tab:ai3-conf} lấy từ kiểm tra
độc lập trên 218 ảnh cylinder và 10 ảnh cube (CPU inference via ai-edge-litert).
}
   %% mục 5.4
%%     ...
%%   \end{document}
%%
%% Compile: pdflatex -> pdflatex (chạy 2 lần để có label/ref đúng)
%%=============================================================
